% !TeX program = lualatex
\documentclass[11pt,a4paper]{ltjsarticle}
\usepackage[margin=25mm,headheight=17pt]{geometry}
\usepackage[haranoaji]{luatexja-preset}
\usepackage{amsmath,amssymb,amsthm,mathtools,mathrsfs,booktabs,array,enumitem,longtable}
\usepackage[unicode,hidelinks]{hyperref}
\usepackage{fancyhdr}
\pagestyle{fancy}\fancyhf{}\fancyfoot[C]{\thepage}
\fancyhead[L]{\small\nouppercase{\leftmark}}
\renewcommand{\sectionmark}[1]{\markboth{#1}{}}
\setlength{\emergencystretch}{3em}\allowdisplaybreaks[2]
\setlength{\parskip}{.15em}
\newtheoremstyle{jp}{.6em}{.6em}{\normalfont}{}{\bfseries}{.}{.5em}{}
\theoremstyle{jp}
\newtheorem{theorem}{定理}[section]
\newtheorem{proposition}[theorem]{命題}
\newtheorem{lemma}[theorem]{補題}
\newtheorem{corollary}[theorem]{系}
\newtheorem{definition}[theorem]{定義}
\newtheorem{remark}[theorem]{注意}
\newtheorem{example}[theorem]{例}
\newtheorem{question}[theorem]{問題}
\renewcommand{\proofname}{証明}
\newcommand{\kk}{\mathbb C}\newcommand{\PP}{\mathbb P}
\newcommand{\OO}{\mathcal O}\newcommand{\A}{\mathcal A}
\newcommand{\B}{B_{\mathrm{VK}}}\newcommand{\Blo}{B_{p,q}}
\newcommand{\HH}{\operatorname{HH}}\newcommand{\Hom}{\operatorname{Hom}}
\newcommand{\End}{\operatorname{End}}\newcommand{\Sym}{\operatorname{Sym}}
\newcommand{\Perf}{\operatorname{Perf}}\newcommand{\Der}{\operatorname{Der}}
\newcommand{\Inn}{\operatorname{Inn}}\newcommand{\gr}{\operatorname{gr}}
\newcommand{\rk}{\operatorname{rank}}\newcommand{\Spec}{\operatorname{Spec}}
\newcommand{\id}{\operatorname{id}}\newcommand{\Ext}{\operatorname{Ext}}
\newcommand{\Tor}{\operatorname{Tor}}\newcommand{\Ann}{\operatorname{Ann}}
\newcommand{\coker}{\operatorname{coker}}\newcommand{\Supp}{\operatorname{Supp}}
\newcommand{\pd}{\operatorname{pd}}\newcommand{\depth}{\operatorname{depth}}
\newcommand{\Frac}{\operatorname{Frac}}\newcommand{\Aut}{\operatorname{Aut}}
\newcommand{\Out}{\operatorname{Out}}\newcommand{\GL}{\operatorname{GL}}
\newcommand{\PGL}{\operatorname{PGL}}\newcommand{\Hilb}{\operatorname{Hilb}}
\newcommand{\dashsim}{\mathrel{\overset{\mathrm{bir}}{\sim}}}
\setcounter{tocdepth}{1}

\hypersetup{pdftitle={重み付きPoisson変形の非被約成分とCalabi--Yau量子化}}
\title{重み付きPoisson変形の非被約成分と\\Calabi--Yau量子化\\
\large 共鳴系列の一般点長さとDu Val中心}
\author{}\date{2026年10月29日\quad 日本語公表用原稿IX・v34}
\begin{document}\maketitle
\begin{abstract}
重み付き射影スタック $\mathbb P(1,1,1,m)$、$m\ge3$ のPoissonスキームを
非被約構造を保持して調べる。
接線共鳴から、同じ重みで一般点の長さが $1,\ldots,\lceil m/2\rceil$ となる
相異なる成分の系列を構成し、全形式局所環を $(\kappa u,\gamma^{j+1}u)$ と決定する。
対応する非線形アフィンLie作用のCGG量子化について、Nakayama自己同型を全パラメータで求め、
通常の4-Calabi--Yau代表とその多項式中心を分類する。
奇数重みの端では中心は $D_{m+4}$ 型のDu Val曲面となる。
引戻し成分には別の長さ $m+3$ のArtin環が現れ、
これらの厚みは標準例外列代数のminiversal変形基底にも残る。
量子化は標準例外列を持ち、次数零の関数斜体は共通のWeyl斜体である。
重み三では少なくとも十成分を認定し、第二の十八次元成分の長さ二を
独立な有理係数消去でも確認する。完全な成分分類は主張しない。
\end{abstract}
\tableofcontents\clearpage
\section{序論と主結果}

本稿は、重み付き射影三次元スタック
\[
 X_m=\mathbb P_{\mathrm{stack}}(1,1,1,m),\qquad m\ge3
\]
のPoisson構造のパラメータ空間と、その明示的な非可換量子化を調べる。
大域的二重ベクトル場のアフィン空間にJacobi方程式を課したスキームを
$\mathscr P_m$ と書く。スカラー倍と重み付き座標変換は商せず、
方程式のイデアルを根基化しない。
この規約の下では、被約な族が一つの座標変換軌道で稠密に覆われる場合にも、
パラメータ空間の一般点に非自明な冪零方向が残る。

主結果は、接する直線と二次曲線から生じる共鳴系列の全形式近傍と、
その一般点を量子化する通常の4-Calabi--Yau代数の中心を決定することである。
幾何側には任意に大きい一般点の長さが現れ、
量子代数側には奇数重みの端でDu Val曲面が現れる。
これらを、別の引戻し成分の厚み、および重み三の成分・境界の計算と合わせて記述する。

\subsection{接線共鳴の成分系列と全形式近傍}

$S_d=\mathbb C[x,y,z]_d$、$N_m=\binom{m+2}{2}$ と置く。
Poisson構造の標準形には
$c\in S_3,D\in\mathfrak{sl}_3,[F]\in S_{m+1}^{\oplus3}/\mathbf xS_m$
を用いる。その全Jacobi式は
\[
 D(c)=0,\qquad
 \frac{m+3}{3}F(c)-\frac m3(\mathbf x\times D)\cdot F-c\operatorname{div}F=0
\]
である。

\begin{theorem}[接線共鳴系列]
\label{thm:main-tangent-v34}
$0\le j\le\lfloor(m-1)/2\rfloor$、$r=m-2j$ とし、
\[
 q=xz-y^2,\quad c=xq,\quad
 D=-\frac{4m+3-6j}{2m}(x\partial_y+2y\partial_z),\quad
 F_j=q^j(0,yz^r/2,z^{r+1})
\]
と置く。組 $(c,D)$ の $\operatorname{SL}_3$ 軌道上で全混合核を動かした族の閉包
$\mathscr E_{m,j}$ は、次元 $N_m+8$ の既約成分である。
その一般点の局所環は $K[\gamma]/(\gamma^{j+1})$ である。
従って固定した $m$ で、一般点の長さが
$1,2,\ldots,\lceil m/2\rceil$ となる相異なる成分を得る。

$F=0$ の基準点では、全完備局所環は
\[
 \boxed{
 \mathbb C[[s_1,\ldots,s_{N_m+7},\kappa,\gamma,u]]/
              (\kappa u,\gamma^{j+1}u)
 }
\]
である。大きい方の枝は次元 $N_m+9$ の $L(1,2,m)$ 成分で、
小さい方が $\mathscr E_{m,j}$ である。
小さい方の被約成分は重み付き座標変換群の稠密な一軌道を持つ。
\end{theorem}

証明では、基底 $D(c)=0$ の七軌道方向を除く二変数切片を取り、
実際の $w$ 移動の核を非被約基底上でも自由な直和因子として除く。
残る一列のSchur補完の可能な次数は四つの余核重量から制約される。
直線と二次曲線を横断させた場合の全混合核を多項式の次数降下で解き、
二つの特異点での評価を使うと、最初の非零障害が正確に $j+1$ 次であると分かる。
有限個の重みの階数検算から一般公式を外挿する証明ではない。

\subsection{Calabi--Yau量子化とDu Val中心}

今度は $j,r\ge0$、$m=2j+r\ge3$ とし、$r=0$ も含める。
導分
\[
 N=x\partial_y+2y\partial_z-q^jz^r\partial_w,\qquad
 X_\alpha=\alpha x\partial_x+(\alpha+1)y\partial_y+(\alpha+2)z\partial_z
                +(m\alpha+m+r+1)w\partial_w
\]
は $[N,X_\alpha]=N$ を満たす。
Coll--Gerstenhaber--Giaquintoの有限和
\[
 a\star b=\sum_{n\ge0}s^nN^n(a)\binom{X_\alpha}{n}(b)
\]
から得る重み $(1,1,1,m)$ の代数を $\mathcal S_{j,r}(\alpha,s)$ とする。

\begin{theorem}[Nakayama公式と中心]
\label{thm:main-quantum-v34}
$\mathcal S_{j,r}(\alpha,s)$ はnoetherian整域かつAS正則四次元代数で、
Hilbert級数は $((1-t)^3(1-t^m))^{-1}$、Gorenstein parameterは $m+3$ である。
$T=1+sN$ と置くと、Nakayama自己同型は
\[
 \boxed{\mu=T^{(m+3)\alpha+m+r+3}}
\]
となる。従って $s\ne0$ のとき通常の4-Calabi--Yau条件は
$\alpha=-(m+r+3)/(m+3)$ と同値である。
このCalabi--Yau族の中心は
\[
 \begin{array}{c|c}
 r\ge2&\mathbb C\\
 r=0&\mathbb C[q,J]\\
 r=1&\mathbb C[S,P,R]/(R^2-SP^2-\tfrac{16}{81}S^{m+3})
 \end{array}
\]
であり、最後の場合は $D_{m+4}$ 型のDu Val曲面である。
ここで $J,S,P,R$ は本文で明示する多項式であり、中心の生成性も証明する。
\end{theorem}

三変数部分代数はLecoutre--Sierraの $R(2,\alpha)$ である。
その既知のNakayama公式\cite[Theorem 4.16]{LecoutreSierra2019}とOre拡大の制限公式を用い、
最後の重み $m$ の生成元に関する剛性を証明することで、上の四変数の公式を得る。
中心の決定には、局所冪零導分の核を四生成元一関係で完全に記述し、
その零重量部分を求める。多項式不変環の生成に関する方法は
Maubach\cite{Maubach2000}の研究に位置付けられる。

$r>0$ のCalabi--Yau族は、前定理の成分の稠密な軌道を実際に量子化する。
全ての族で標準 $(m+3)$ 項列はfull strongであり、
$(0,1,m,m+1)$ の四項部分列は許容部分圏を与える。
特に $m=3$ では、二つの十八次元成分の双方が非可換VK圏の標準六項実現を持つ。
固定した非零量子代数自体の変形空間の一般点長さを、前定理で計算したという意味ではない。

二つの定理を奇数重みの最後の共鳴に適用すると、次の対応を得る。
\[
 \boxed{
 m=2\ell-1,\ j=\ell-1,\ r=1\quad(\ell\ge2):\qquad
 \text{一般点の長さ }\ell,\qquad
 \text{量子中心の型 }D_{2\ell+3}.
 }
\]
一方、次数零の量子関数斜体は全て
$\mathbb C(C)$ 上の第一Weyl斜体であり、$C=J^2/q^{m+r+1}$ である。
従って多項式中心の差は、この双有理表示だけでは検出できない。
異なる重みの射影圏は $K_0$ の階数 $m+3$ が異なるため導来同値でないが、
同じ重みの異なる共鳴の射影圏の比較は残る。

\subsection{引戻し成分に現れる別の厚みと圏の変形基底}

\begin{theorem}[引戻し成分]
\label{thm:main-v33}
全ての $m\ge3$ で、$\mathbb P(1,1,m)$ からの引戻し族の閉包は
次元 $N_m+9$ の既約成分であり、一般点の局所環は
\[
 K[a,b]/(a^{m+2},ab,b^2)
\]
となる。その長さは $m+3$ である。
$L(1,1,1,m)$ 成分との一般的な交叉では、全完備局所環は
\[
 \mathbb C[[s_1,\ldots,s_{N_m+8},a,b,u]]/(ab,b^2,a^{m+2}u)
\]
となる。この局所環はCohen--Macaulayであるが完全交叉ではない。
\end{theorem}

接線系列の一変数のArtin環と異なり、この引戻し成分の一般点には
二つの独立な接方向が残る。
全重みでの証明には、実際の座標移動核、重量による障害次数の制約、
三角形への変形と三交点での評価を用いる。
この族には発散零の微分Ore拡大による、別の明示的な4-Calabi--Yau量子化も与える。

標準例外列代数
\[
 \Lambda_m=\operatorname{End}_{X_m}
       \left(\bigoplus_{i=0}^{m+2}\mathcal O_{X_m}(i)\right)
\]
のminiversal変形基底は、非標準的な形式座標の選択の下で
$\mathscr P_m$ の原点完備化に同型である。
変数数は $m^2+6m+26$、独立な二次生成元の数は $10+\binom{m+5}{2}$ となる。
斉次錐の一般点長さは原点完備化でも保存されるので、
上の二種類の非被約構造はこの圏の変形基底にも現れる。
これは自己同型を含む変形groupoidのfine表現を主張するものではない。

\subsection{重み三の成分と境界}

重み三の問題は、Veronese--Kronecker四項切り出しを
標準六項へ補完する問題から生じる\cite{VKAmbient32}。
以下の相異なる十成分を認定する。
\begin{center}\small
\begin{tabular}{lcl}\toprule
一般的な構成&次元&一般点の性質\\\midrule
$\operatorname{PB}(\mathbb P^2)$&35&滑らか\\
有理型 $R(1,5)$&28&滑らか\\
対数型 $L(1,1,4)$&23&滑らか\\
有理型 $R(2,4)$&22&滑らか\\
有理型 $R(3,3)$&19&滑らか\\
対数型 $L(1,2,3)$&19&滑らか\\
対数型 $L(1,1,1,3)$&19&滑らか\\
$\operatorname{PB}(\mathbb P(1,1,3))$&19&長さ六\\
接線共鳴 $a=-5/2$&18&滑らか（長さ一）\\
接線共鳴 $a=-3/2$&18&長さ二\\\bottomrule
\end{tabular}
\end{center}
$R(p,q)$ は $pf\,dg-qg\,df$ の族、$L(d_1,\ldots,d_k)$ は
$(\prod f_i)\sum\lambda_i\,df_i/f_i$ の族の閉包を表す。
十成分が全てであること、三十五が最大次元であることは主張しない。

接線型の四つの非零共鳴では、基底変数を含む全形式近傍を
有理係数のSchur消去でも独立に求める。
非共点三直線の特殊点では五つの局所枝を持つ非被約環を決定し、
$D=0$ の全非零三次式の混合核は留数で記述する。
$c=0$ の閉部分の被約化は次元 $35,20,18$ の三成分からなり、
全て既知の大域成分に含まれることも示す。

\subsection{先行研究との関係と論証の範囲}

三次元のPoisson構造と反標準束で捻った可積分一形式の対応、
有理型・対数型・引戻しという構成自体は古典的である。
通常の射影三空間におけるPoisson構造と量子代数の関係はPym\cite{Pym}、
関連する変形論は\cite{Molinuevo2016}、
閉有理一形式の留数は\cite{PereiraClosedForms}を参照されたい。
CGG公式、AS正則性に関するOre拡大定理、既知の三変数Nakayama公式を
本稿の新しい構成原理として扱わない。
本稿で直接求めるのは、この重み付き族の成分、全局所環、
非線形な最後の生成元の剛性、および中心の具体的な表示である。

重み付き射影空間を標的とする引戻しの安定性\cite{GAMQV2026}は、
通常の射影空間を始域とする場合などに仮定を課しており、
その一般被約性を本稿の重み付き三次元始域へ直接適用することはできない。
葉層スキームの非被約性自体にも先行例がある
\cite[Remark 4.2]{BKM2026}。
本稿では一般点のArtin環と交叉環を決定し、整数パラメータとその長さの関係まで証明する。
一般重みの成分定理・中心の生成性と、有限の有理数検証の役割は明確に区別する。

\section{一般重みのJacobiスキームと圏の変形基底}
\label{sec:weighted-hulls-v33}

この節では $m\ge3$ とし、
\[
 R_m=\mathbb C[x,y,z,w],\quad \deg(x,y,z,w)=(1,1,1,m),\qquad
 X_m=\mathbb P_{\mathrm{stack}}(1,1,1,m)
\]
と置く。$S_d=\mathbb C[x,y,z]_d$、$N_m=\dim S_m=\binom{m+2}{2}$ と書く。
大域的二重ベクトル場のアフィン空間の中で、Jacobi式の係数が生成する
イデアルによって定まるスキームを $\mathscr P_m$ とする。
この定義ではイデアルを根基化せず、スカラー倍も座標変換も商しない。

\subsection{全重みに共通する標準形}

\begin{proposition}[一般重みの方程式]
\label{prop:general-equations-v33}
$\mathscr P_m$ は
\[
 (c,D,[F])\in S_3\oplus\mathfrak{sl}_3\oplus
       (S_{m+1}^{\oplus3}/\mathbf xS_m),\qquad \mathbf x=(x,y,z),
\]
を座標とする $d_m=m^2+6m+26$ 次元アフィン空間の閉部分スキームであり、
その定義式は
\begin{align}
 D(c)&=0,\label{eq:general-base-v33}\\
 M_m(c,D;F)&=
 \frac{m+3}{3}F(c)-\frac m3(\mathbf x\times D)\cdot F-c\operatorname{div}F=0.
 \label{eq:general-mixed-v33}
\end{align}
これらの係数は $q_m=10+\binom{m+5}{2}$ 個の一次独立な二次式である。
\end{proposition}

\begin{proof}
$E=E_x+mw\partial_w$ を重み付きEuler場とする。
$m\ge3$ なので、重み零の二重ベクトル場は
$P=\Pi+(F+wA)\wedge\partial_w$ と書け、$\Pi$ の係数は $S_2$ に属する。
$E$ を因子に持つ場を加えて $wA$ を消去すると、残る同値は
$F\mapsto F+g\mathbf x$、$g\in S_m$ だけである。
従って場の空間の次元は
\[
 18+3\binom{m+3}{2}-\binom{m+2}{2}=m^2+6m+26.
\]
$\Pi=p_1\partial_y\wedge\partial_z+p_2\partial_z\wedge\partial_x
+p_3\partial_x\wedge\partial_y$ と書くと、
$c=\mathbf x\cdot p$、$D=\operatorname{curl}p$ によって
$p=(\nabla c-\mathbf x\times D)/3$ となる。

射影上のJacobi条件は $E\wedge[P,P]=0$ である。
$w$ の係数から $D(c)=0$、定数項から
$\mathbf x\cdot q=0$ を得る。ここで
$q=F(p)+(\operatorname{Jac}F)^tp-(\operatorname{div}F)p$ は $[F,\Pi]$ の係数である。
$F$ の係数が次数 $m+1$ なので
\[
 \mathbf x\cdot q=F(c)+m\,p\cdot F-c\operatorname{div}F=M_m(c,D;F).
\]
$F=g\mathbf x$ を代入すると零になるため、式はEuler商上で定義される。

第一式の十係数の独立性は $\mathfrak{sl}_3 S_3=S_3$ から従う。
第二式で $c=0$ とし、$D$ を非対角行列単位に取ると、
$(\mathbf x\times D)_i$ の一成分を $\pm x_j^2$ とできる。
次数 $m+3\ge6$ の任意の単項式はある $x_j^2$ で割れるため、
適切な $F_i\in S_{m+1}$ によって第二式の像は $S_{m+3}$ を張る。
$F=0$ と置けば二群を分離できるので、全 $q_m$ 個が独立である。
\end{proof}

\subsection{一意な発散零代表と全Jacobi軌跡の量子化}

\begin{lemma}[重み付きEuler切断]
\label{lem:general-lift-v33}
標準体積形式 $dx\wedge dy\wedge dz\wedge dw$ に関するBV発散を $\delta$ とする。
重み零の多重ベクトル場について
\[
 \delta(E\wedge U)+E\wedge\delta U=(m+3)U.
\]
従って $L(U)=U-E\wedge\delta U/(m+3)$ はEuler商から発散零代表への一意な切断であり、
Schouten括弧を保つ。この切断はPoisson場を通常の多項式Poisson場へ持ち上げる。
\end{lemma}

\begin{proof}
BV恒等式の右辺は $(\operatorname{div}E)U+[E,U]$ であり、
重み零では後項が零、前項が $(m+3)U$ である。
$\delta^2=0$ から切断の存在を得る。
$E\wedge U=0$ かつ $\delta U=0$ なら同じ恒等式から $U=0$ なので一意である。
Euler場を因子に持つ場は重み零のSchouten代数のイデアルであり、
発散零の場も括弧で閉じる。一意性により切断は括弧を保つ。
\end{proof}

標準形に対する持ち上げは具体的に
\[
 \widehat P=\Pi+\frac1{m+3}E_x\wedge D+
 \left(F-\frac{\operatorname{div}F}{m+3}\mathbf x
               -\frac m{m+3}wD\right)\wedge\partial_w.
\]
通常の射影空間におけるこの構成は \cite[\S2.2]{Pym}、
重み付きPoisson代数の分解は \cite{TWW} と関係する。
上の証明は本稿で使う一意性と括弧の保存を直接与えている。

$\mathscr P_m$ の原点で完備化し、普遍Poisson場を上の切断で持ち上げる。
Kontsevich量子化 \cite{Kontsevich} を係数の極大イデアル位相で適用すれば、
全Jacobi軌跡上に平坦な形式結合代数を得る。
各グラフ演算は線形座標変換に同変であり、重み付き $\mathbb G_m$ 作用も線形なので、
量子化は重みを保つ。各次数成分の台は可換環の次数成分の自由な形式拡張である。

\subsection{有限代数のminiversal基底}

\[
 \Lambda_m=\operatorname{End}_{X_m}
    \left(\bigoplus_{i=0}^{m+2}\mathcal O_{X_m}(i)\right)
\]
と置く。重み付き多項式環のGorenstein parameterは $m+3$ であり、
標準列はfull strongである \cite[Corollary~2.7]{Orlov2009}。
従って $\operatorname{Perf}(\Lambda_m)\simeq\operatorname{Perf}(X_m)$ である。
また
\[
 \dim\Lambda_m=\sum_{d=0}^{m+2}(m+3-d)\dim(R_m)_d
              =\binom{m+6}{4}+15.
\]
最後の十五は $w,wS_1,wS_2$ の寄与 $3+6+6$ である。

\begin{theorem}[一般重みの二次的miniversal基底]
\label{thm:general-hull-v33}
$\Lambda_m$ の結合代数変形のminiversal基底は、非標準的な形式同型を除いて
\[
 R_{\Lambda_m}\simeq\widehat{\mathcal O}_{\mathscr P_m,0}
 =\mathbb C[[u_1,\ldots,u_{d_m}]]/J_m
\]
である。ここで $J_m$ は命題~\ref{prop:general-equations-v33}の $q_m$ 個の独立な二次式で生成される。
自己同型を含む変形groupoidのfine代表性や、座標の標準性は主張しない。
\end{theorem}

\begin{proof}
Euler列とその外積、および非負次数線束の高次コホモロジー消滅から、
$H^i(X_m,\wedge^jT_{X_m})=0$、$i>0$、$0\le j\le3$ となる。
非自明慣性成分は余次元三の点 $B\mu_m$ だけなので、
スタック版HKR \cite[Corollary~4.16]{FPSS} により
\[
 HH^2(\Lambda_m)=HH^2(X_m)=H^0(X_m,\wedge^2T_{X_m}).
\]
この恒等慣性成分の多重微分作用素によるHKR代表について、Gerstenhaber一次障害は
通常のSchouten括弧である。従って二次障害の独立な式は正確に $J_m$ の二次生成元となる。
非自明慣性成分から新しい接方向はない。

上の全Jacobi軌跡の量子化を次数差 $0,\ldots,m+2$ に評価すると、
$P/J_m$ 上の $\Lambda_m$ の全形式変形を得る。ここで $P=\mathbb C[[u_1,\ldots,u_{d_m}]]$。
その接写像はHKR同型である。
miniversal基底を $P/I$ と書き、形式的versalityによる写像
$P/I\to P/J_m$ を接写像が恒等となる座標で持ち上げると
$I\subset J_m$ かつ $I_2=(J_m)_2$ となる。
二次生成性から $J_m\cap\mathfrak m_P^3=\mathfrak m_PJ_m$ なので
$J_m=I+\mathfrak m_PJ_m$ である。有限生成加群 $J_m/I$ にNakayamaの補題を適用して
$I=J_m$ を得る。これは公表稿VIII \cite{VKAmbient32} の認定法の一般重み版でもある。
\end{proof}

\begin{remark}[三次Hochschild群と二次式の個数]
各非自明慣性成分 $B\mu_m$ の法束行列式は、全安定化群の文字 $\chi^3$ を持つ。
これは慣性成分を添字付ける個々の元の行列式だけを判定する問題ではない。
$m\ge4$ では $\chi^3$ の不変部分は零なので、twisted $HH^3$ は零であり、
\[
 \dim HH^3(\Lambda_m)=\dim(R_m)_{m+3}
 =\binom{m+5}{2}+10=q_m\qquad(m\ge4).
\]
$m=3$ では $w^2$ のuntwisted一方向と、二つのtwisted方向が加わって
$\dim HH^3(\Lambda_3)=41$、二次括弧の像は三十八次元である。
従って $m\ge4$ では一次方向間の括弧の線形包が $HH^3$ 全体を張る。
この全射性はJacobi式が正則列であることを意味しない。
\end{remark}

\subsection{原点での形式化と一般点の長さ}

\begin{lemma}[錐の成分と形式化]
\label{lem:cone-completion-length-v33}
$A$ を標数零体上の有限生成非負次数付き代数、$A_0=k$ とし、
$\widehat A$ を正次数イデアルでの完備化とする。
斉次最小素イデアル $\mathfrak p$ に対し $\mathfrak p\widehat A$ は最小素イデアルであり、
\[
 \operatorname{length}A_{\mathfrak p}
 =\operatorname{length}(\widehat A)_{\mathfrak p\widehat A}.
\]
\end{lemma}

\begin{proof}
$A/\mathfrak p$ は次数付き整域である。その完備化の非零元は最小の非零斉次項を持ち、
二つの最小項の積は零にならないため、完備化も整域である。
従って $\mathfrak p\widehat A$ は素イデアルとなる。
原点で局所化してから完備化する写像の忠実平坦性と最小素イデアルの性質から、これは最小である。
局所化した平坦局所写像
$A_{\mathfrak p}\to(\widehat A)_{\mathfrak p\widehat A}$ の閉ファイバーは
$\operatorname{Frac}(A/\mathfrak p)$ から
$\operatorname{Frac}(\widehat{A/\mathfrak p})$ への体拡大である。
長さ有限の組成列を基底変換すれば各商は後者の剰余体になるので、長さは保存される。
\end{proof}

この補題により、後で求める $\mathscr P_m$ の一般点の厚みは、
$R_{\Lambda_m}$ の対応する形式成分の一般点の厚みにもなる。
量子化を固定した $\hbar=1$ の非可換代数の全変形空間と、このminiversal基底とを
同一視する主張は含まない。また $m>3$ の四対象cornerの変形が無障害であることや、
その接空間が外側と一致することは、この節からは導かない。

% Independent all-weight proof; baseline v32 is unchanged.
\section{任意の重みで増大する一般点の非被約厚み}
\label{sec:all-weights-concurrent-v33}

この節では整数 $m\geq3$ を固定し、
$X_m=\mathbb P_{\mathrm{stack}}(1,1,1,m)$ 上の大域的 Poisson テンソルの
アフィンスキームを $\mathscr P_m$ と書く。係数体は標数零とする。
$S_d=k[x,y,z]_d$、$N_m=\dim S_m=\binom{m+2}{2}$ と置く。

方程式と記号は命題~\ref{prop:general-equations-v33}のものを用いる。
最初に、$m$ に依存しない基底切片を決定する。
% Complete base slice before imposing the mixed F equations.
\subsection{共点三直線における基底の非被約な形式切片}
\label{subsec:concurrent-base-v32}

$\mathscr Z=\{(c,D)\in S_3\oplus\mathfrak{sl}_3:D(c)=0\}$ と置く。
この基底スキーム自身に既に非被約構造が現れる。
以下の補題は混合項 $F$ の方程式をまだ課さない、基底だけの全形式的な記述である。

\begin{lemma}[共点三直線の基底切片]
\label{lem:concurrent-base-v32}
\[
 c_0=xy(x+y),\qquad D_0=(Ax+By)\partial_z
\]
とし、
\[
 U=B-2A,\qquad V=A-2B,\qquad W=A+B
\]
が全て非零とする。例えば $(A,B)=(2,5)$ はこの条件を満たす。
この点における $\mathscr Z$ の完備局所環は
\[
 \frac{\mathbb C[[u_1,\ldots,u_8,a,b]]}{(ab,b^2)}
\]
に同型である。八つの滑らかな変数のうち六つは $c_0$ の
$\mathrm{SL}_3$ 軌道方向、残る二つは $A,B$ の変動である。
被約化は滑らかな九次元であり、零でない平方零イデアル $(b)$ は
共点三直線の部分 $a=0$ に支持される。
\end{lemma}

\begin{proof}
$c_0$ の $\mathfrak{sl}_3$ による接軌道は六次元である。
その像は二変数三次式の四次元空間と
$z\partial_xc_0,z\partial_yc_0$ の二方向からなる。
従って、四次元の形式横断切片を
\[
 c=xy(x+y)+a xyz+b xz^2+\gamma yz^2+d z^3
\]
と取れる。六つの独立な軌道方向を形式的な群座標として選べば、
軌道方向とこの切片から $S_3$ への写像の接写像は同型であり、形式逆写像定理が適用できる。
$D$ も同じ $\mathrm{SL}_3$ 座標変換で移すので、$D(c)=0$ の問題は
この六次元の滑らかな因子と切片上の問題に分かれる。

切片上で $D_{zx}=A,D_{zy}=B$ と書く。
$D(c)=0$ の四つの二変数三次単項式と $x^2z,y^2z$ の係数から、
残る六つの係数は一意に
\[
\begin{aligned}
 D_{yx}=D_{xy}&=0,\\
 D_{xx}&=\frac{a(B-2A)}3,&
 D_{yy}&=\frac{a(A-2B)}3,\\
 D_{xz}&=-2B\gamma,&D_{yz}&=-2Ab,\\
 D_{zz}&=-D_{xx}-D_{yy}
\end{aligned}
\]
と決まる。代入後の残る四式は、零でない定数倍を除いて
\[
\begin{aligned}
 E_1&=Ub+V\gamma,\\
 E_2&=aUb+3Ad,\\
 E_3&=aV\gamma+3Bd,\\
 E_4&=W(ad-2b\gamma)
\end{aligned}
\]
である。

この段階で $U,V,W$ は形式局所環の単元である。まず
\[
 E_2+E_3-aE_1=3Wd
\]
から $d$ を消去でき、次に $E_2-3Ad=Uab$ から $ab=0$ を得る。
また $E_1=0$ により $\gamma=-Ub/V$ である。
最後に $E_4=0$ から $b\gamma=0$、従って $b^2=0$ が従う。
逆に $d=0,\gamma=-Ub/V,ab=b^2=0$ を満たせば四つの $E_i$ は全て零になる。
従ってこれは零点集合だけの比較でなく、局所化された関係イデアルの等式である。

残った自由変数は六つの軌道座標と $A,B,a,b$ であり、表示した完備局所環を得る。
$\mathbb C[[a,b]]/(ab,b^2)$ において $b\ne0$、$b^2=ab=0$ だから、
非被約性とその支持も明らかである。
なお、消去には $A$ や $B$ 自身の可逆性は必要なく、上記三つの単元条件で十分である。
\end{proof}

被約化の切片では $c=xy(x+y+az)$ となる。
従って $a\ne0$ は非共点の三直線、$a=0$ は共点三直線に対応する。
この補題だけから、混合項を含む全 Poisson スキームの既約成分性や
一般点の非被約厚みの次数を結論してはならない。
その判定には、この環上での混合写像をさらに解析する必要がある。


\subsection{全重みに共通する局所環}

\begin{theorem}[一般点の厚み $m+3$]
\label{thm:allm-concurrent-v33}
各整数 $m\geq3$ に対し、$(A,B)=(2,5)$ を含む非空 Zariski 開集合が存在し、
\[
 c_0=xy(x+y),\qquad D_0=(Ax+By)\partial_z,\qquad [F]=0
\]
における $\mathscr P_m$ の完備局所環は
\[
 \boxed{\quad
 \frac{k[[s_1,\ldots,s_{N_m+8},a,b,u]]}
 {(ab,b^2,a^{m+2}u)}
 \quad}
\]
と同型である。開集合は $m$ に依存してよい。
被約化の二つの枝の次元はともに $N_m+9$、その被約交叉は滑らかな
$N_m+8$ 次元である。$u\ne0$ の枝の一般点では横断局所環が
\[
 \frac{K[[a,b]]}{(a^{m+2},ab,b^2)}
\]
となり、その長さは $m+3$ である。
ここで $K$ は当該被約成分の関数体であり、係数体の選択は非標準的である。
\end{theorem}

以下の証明は任意の $m$ を一度に扱う。
有限個の重みにおける行列計算から帰納的な一般化を推測するものではない。

\subsubsection*{基底切片と二つの次数上昇}

$D(c)=0$ は $m$ に依存しない。
共点三直線の基底切片の消去により、六つの軌道座標を別にして
\begin{align*}
 c&=xy(x+y)+a xyz+b(xz^2+\gamma yz^2),\\
 D_x&=\frac{a(B-2A)}3x-2B\gamma bz,\\
 D_y&=\frac{a(A-2B)}3y-2Abz,\\
 D_z&=Ax+By+\frac{a(A+B)}3z,\qquad
 \gamma=-\frac{B-2A}{A-2B},
\end{align*}
かつ関係は正確に $ab=b^2=0$ となる。
この消去には $(B-2A)(A-2B)(A+B)\ne0$ だけを要し、$(2,5)$ で成立する。
$A,B$ の二変数と合わせて八つの滑らかな基底変数を得る。

混合写像を $M=M_0+aM_1+bM_2$ と書く。
$F_i=x^{n_1}y^{n_2}z^{n_3}$ の重量を
$n_3-\mathbf1_{i=3}$ とし、標的には $z$ 次数を与える。
このとき $M_0,M_1,M_2$ はそれぞれ重量を $0,1,2$ 上げる。
この主張は各項の次数から直接従い、係数 $A,B$ や $m$ には依存しない。

\subsubsection*{共点側の核を全重みで計算する}

$L=Ax+By$ と置き、二変数の二次ベクトル場
\[
 V=\frac13\bigl((c_{0,y}+xL)\partial_x
              -(c_{0,x}-yL)\partial_y\bigr)
\]
を取る。$\Pi_0=\partial_z\wedge V$ であり、
$V(c_0)=Lc_0$、$\operatorname{div}V=L$ である。
$F=Q+H\partial_z$ と書く。
混合式を三直線へ制限すると、次の三つの係数が非零である限り $Q$ は各直線に接する。
\[
 \kappa_1=1+\frac m3(1-B),\qquad
 \kappa_2=1+\frac m3(1+A),\qquad
 \kappa_3=1+\frac m3(1-A+B).
\]
実際、$x=0$ では $\kappa_1 y^2F_x=0$、$y=0$ では
$\kappa_2 x^2F_y=0$、$x+y=0$ では
$-\kappa_3x^2(F_x+F_y)=0$ となる。
この条件は $(2,5)$ において
\[
 (\kappa_1,\kappa_2,\kappa_3)
 =\left(1-\frac{4m}3,1+m,1+\frac{4m}3\right)
\]
なので、全ての $m\geq3$ で成立する。

三直線に接する二変数導分の自由基底は $E_2,V$ である。
ここで $\det(E_2,V)=-c_0$ であるから、これは係数 $A,B$ によらない。
Euler 同値で $E_2$ 成分を消すと、重量 $j\geq0$ の核は
\[
 Q=z^jgV,\qquad H=-\frac{z^{j+1}}{j+1}V(g),
 \qquad g\in k[x,y]_{m-1-j},\quad 0\leq j\leq m-1.
\]
なぜなら混合式の左辺は、この代表に対して
$-c_0(V(g)z^j+\partial_zH)$ となるからである。
これらは全て $-z^{j+1}g/(j+1)$ による次数 $m$ の $w$ 移動である。
重量 $-1$ では $H=f(x,y)$、$\deg f=m+1$ の $m+2$ 方向が残る。
従って
\[
 \dim\ker M_0=\frac{m(m+1)}2+(m+2)=N_m+1.
\]

二変数の次数 $m$ の $w$ 移動を明示すると、
$g_k=x^{m-k}y^k$、$e_k=x^{m+1-k}y^k$ に対し
\[
 V(g_k)=\left(\frac{m(A+1)}3-k\right)e_k
       +\left(\frac{m(B+2)}3-k\right)e_{k+1},
 \quad 0\leq k\leq m.
\]
$(2,5)$ では後者の係数が $7m/3-k>0$ である。
従って二変数移動の像は $m+1$ 次元、その余核は $e_0=x^{m+1}$ で生成される。
非二変数の移動の単射性は上の $Q=z^jgV$ の表示から分かる。
結局、全ての次数 $m$ の移動の像は $N_m$ 次元で、混合核との商は
重量 $-1$ の一点
\[
 F_{\mathrm{ex}}=(0,0,x^{m+1})
\]
で生成される。

以上から重量ブロックの大きさと階数は正確に次のようになる。
\[
\begin{array}{c|c|c|c}
 \text{重量 }r&\text{源}&\text{標的}&\operatorname{rank}M_0\\\hline
 -1&m+2&0&0\\
 0\leq r\leq m-1&2(m+2-r)&m+4-r&m+4-r\\
 m&4&4&4\\
 m+1&2&3&2\\
 m+2&0&2&0\\
 m+3&0&1&0
\end{array}
\]
源の次元は、例えば $F_x$ の $x$ で割り切れる単項式を Euler 同値で
消すことにより数えられる。余核は四次元で、重量は
$m+1,m+2,m+2,m+3$ である。

\subsubsection*{非共点側の核と非共鳴条件}

$b=0$ とし、$a$ を反転させる。
$X=x,Y=y,Z=x+y+az$ の線形座標変換を行うと、
係数ベクトルの体積変換を含めて
\[
 c'=aXYZ,\qquad
 D'=\frac a3\bigl((B-2A)X\partial_X
 +(A-2B)Y\partial_Y+(A+B)Z\partial_Z\bigr).
\]
標的の単項式 $X^{\nu_1}Y^{\nu_2}Z^{\nu_3}$、
$\sum\nu_i=m+3$ を固定する。
全ての $\nu_i$ が正なら、対応する源は Euler 同値の前で三次元、
同値の後で二次元であり、$M/a$ の行は
\[
 (\kappa_1-\nu_1,\kappa_2-\nu_2,\kappa_3-\nu_3).
\]
$\sum\kappa_i=m+3$ なので、この行は Euler 同値を殺す。
その階数が零となるのは $(\kappa_1,\kappa_2,\kappa_3)=\nu$ の場合だけである。
一つの $\nu_i$ が零で、他の二つが正なら、源は一次元で係数は $\kappa_i$ である。
二つ以上が零なら源は存在しない。

従って、$\kappa_i\ne0$ かつ $\kappa$ が正整数三つからなる上記有限集合を
避けるとき、核の次元は内部単項式の数 $N_m$ に等しい。
$(2,5)$ では $\kappa_1<0$ なので、この条件も全ての $m\geq3$ で成立する。
次数 $m$ の $w$ 移動は常に実際の核に入るから、ここではその像が核の全体である。
特に $a$ を反転すると、混合写像の階数は共点側から正確に一つ増える。

\subsubsection*{Schur 補完と三交点評価}

各重量で $M_0$ の可逆小行列を選び、対応する源の変数を消去する。
$r\leq m$ の標的行は全て消去し、$r=m+1$ では
$xy z^{m+1},y^2z^{m+1}$ の二行を消去する。
後者の二次小行列式は $(2,5)$ において
\[
 -\frac{(3-4m)(6+7m)}9\ne0
\]
である。従って、消去後に残る標的の四行は
\[
 x^2z^{m+1},\quad xz^{m+2},\quad yz^{m+2},\quad z^{m+3}
\]
と取れる。必要な有限個の小行列式を反転させた非空開集合上で、この選択は保たれる。

次数 $m$ の移動の写像
\[
 T:S_m\longrightarrow S_{m+1}^3/S_m\mathbf x,
 \qquad g\longmapsto-\,p\times\nabla g
\]
は $D(c)=0$ 上で常に混合方程式を満たす。
特殊点における階数が $N_m$ なので、$N_m$ 次小行列式が局所環の単元であり、
$T(S_m)$ は非被約基底 $k[[a,b]]/(ab,b^2)$ 上でも自由な直和因子である。
Schur 消去後にも、その像は残った源の自由加群の直和因子となる。
実際、Schur 消去で消した座標は残りの座標から一意に復元されるから、
特殊点における像の階数は変わらない。
従って $N_m$ 個の核座標と、残余の一座標 $u$ を選べる。

$b=0$ とした、まだ核座標を取り直していない Schur 写像を $R(a)$ とする。
可逆小行列の定数部は重量を保ち、その $a$ 部は重量を一つ上げる。
従って逆行列の幾何級数は重量の有限性により有限和である。
重量 $w$ の自由列から重量 $r$ の自由行への成分は、零でなければ
正確に $a^{r-w}$ の定数倍である。
特に $r=m+1$ の行では、全ての成分の $a$ 次数は $m+2$ 以下である。

$R(a)$ の最初の非零 $a$ 係数は
$\ker M_0/T(S_m)$ を通る。
これは $R(a)T_{\mathrm{free}}(a)=0$ の最初の非零係数を取り出せば分かる。
この商の重量は $-1$、余核の最小重量は $m+1$ である。
従って残余写像は $a^{m+2}$ で割り切れる。
同じ議論で $b$ の一次係数は重量を二つしか上げず、余核に届かないため零である。
$ab=b^2=0$ なので、$b$ を含む他の係数もない。

ここで先頭 $a^{m+2}$ 係数の非零性を、全ての $m$ に対して証明する。
Schur 消去の直後、残余出力は上記四行だけに支持される。
$b=0$ では $c=xy(x+y+az)$ であり、三つの交点
\[
 [0,0,1],\qquad[-a,0,1],\qquad[0,-a,1]
\]
において混合写像の像は消える。
$a\ne0$ では三角形座標の単項式表示から明らかであり、
多項式恒等式として $a=0$ にも及ぶ。
出力を
\[
 r_0x^2z^{m+1}+r_1xz^{m+2}+r_2yz^{m+2}+r_3z^{m+3}
\]
と書けば、三交点での評価から
\[
 r_3=0,\qquad a^2r_0-ar_1=0,\qquad-ar_2=0.
\]
$b=0$ の係数環は整域なので、$r_1=ar_0,r_2=r_3=0$ である。
次数上昇と残余商の重量から、$r_0$ の最初に可能な項は
$a^{m+2}$ である。一方、まだ核座標を取り直していない $R(a)$ の
$r_0$ 行は、前段により $a$ の次数 $m+2$ 以下の多項式である。
従ってこの係数が零なら $r_0$ 行は恒等的に零となる。
従って先頭係数が零なら残余写像全体が零となるが、これは $a$ を反転したとき
階数が一つ増えることに矛盾する。

従って $(2,5)$ で先頭係数は非零であり、非空開集合を縮めれば単元である。
残余の関係イデアルは正確に $(a^{m+2}u)$ となる。
八つの滑らかな基底座標と $N_m$ 個の移動座標を合わせて、定理の局所環を得る。

\subsubsection*{一般点の長さ}

被約化の二つの素イデアルは $(a,b)$ と $(b,u)$ である。
前者の枝で $u$ を反転すると、横断環は
\[
 K[a,b]/(a^{m+2},ab,b^2)
\]
となり、その $K$ 基底は $1,a,\ldots,a^{m+1},b$ である。
係数体を幾何学的な拡大へ取る必要はない。
完備局所環の係数体を選び、長さ $m+2$ の主鎖を与える $a$ と、
その最終項から独立な socle 元 $b$ を選べば、表示した $K$ 形式を得る。
これで定理が証明された。

\begin{remark}[計算による確認の役割]
付属の有理数検証コードは複数の重みで重量ブロック、Schur 消去、三交点評価、
$b$ 項の消滅を独立に確認する。
全重みの結論は、ここで証明した核の分解、三角形側の非共鳴条件、
重量と三交点評価から従うのであり、有限標本に基づく推定ではない。
\end{remark}

\section{引戻し成分、対数型成分、および交叉の厚み}
\label{sec:allm-geometry-v33}

\subsection{重み付き射影平面への引戻し}

$X_m=\mathbb P_{\mathrm{stack}}(1,1,1,m)$ から
$Y_m=\mathbb P_{\mathrm{stack}}(1,1,m)$ への一般の重み付き射影を考える。
$Y_m$ の同次座標を $(x,y,w)$ とすれば、重み $m+3$ の水平一形式は一意に
\[
 \eta=c_3\,dw-\frac m3w\,dc_3
       +\left(f_{m+1}-\frac m3wL_1\right)(x\,dy-y\,dx)
\]
と書ける。$c_3,L_1,f_{m+1}$ は二変数の、それぞれ三次、一次、$m+1$ 次の式である。
係数空間の次元は
\[
 4+2+(m+2)=m+8
\]
となる。
この表示は、$dw$ 係数を $c_3$ と置き、水平性から $w$ を含む項を解き、
残りの二変数の水平形式が $(x\,dy-y\,dx)$ の倍数になることから従う。
$m>3$ では次数だけから $dw$ 係数が二変数三次式となる。
$m=3$ では一見 $\lambda w$ も $dw$ 係数に許されるが、水平性を表す式の
$w^2$ 係数は $3\lambda$ だけなので $\lambda=0$ である。
ここで $dx,dy$ の係数の次数は $m+2<2m$ であり、$w^2$ を含めないことを用いた。

二次元の標的では可積分条件は自動的に成立する。
従って、これらの形式を $X_m$ へ引き戻し、
$\wedge^2T_{X_m}\simeq\Omega^1_{X_m}\otimes K_{X_m}^{-1}$ により
対応する Poisson テンソルを取れば、$\mathscr P_m$ の族を得る。
この族の閉包を $\mathscr C_m$ と書く。

射影の自由度は、三次元の一次形式空間の二平面を選ぶ $2$ 方向と、
$w$ に加える次数 $m$ の式を二変数式で割った $N_m-(m+1)$ 方向である。
従って総パラメータ数は
\[
 2+\{N_m-(m+1)\}+(m+8)=N_m+9.
\]
これは一般ファイバーの次元を無視した数え上げではない。
一般の $c_3$ は三つの異なる共点直線に分解するため、選んだ二平面はその三因子の
張る空間から復元できる。
二平面を固定すれば、非二変数の $w$ 移動は、前節で求めた
重量 $0,\ldots,m-1$ の混合核を一意に与える。
標的形式による残りの自由度は重量 $-1$ にある。
従って、このパラメータ表示は一般点で余分な正次元ファイバーを持たない。

\begin{corollary}[任意に大きい一般点の長さを持つ成分]
\label{cor:allm-pullback-component-v33}
各 $m\geq3$ に対し、$\mathscr C_m$ は $\mathscr P_m$ の既約成分であり、
\[
 \dim\mathscr C_m=N_m+9.
\]
その一般点の局所環の長さは $m+3$ である。
したがって、重み付き射影三次元スタック上の大域的 Poisson スキームには、
一般点の非被約な長さが任意に大きくなる成分が存在する。
\end{corollary}

\begin{proof}
パラメータ空間は既約であり、上の復元から像の次元は $N_m+9$ である。
一般の二変数三次式は線形座標変換と倍率によって $xy(x+y)$ に直せる。
従って前節の局所模型における $a=b=0$ の枝を、この引戻し族が一般点で覆う。
その枝の次元も $N_m+9$ なので、閉包は一つの既約成分である。
この枝上の一般点では残余混合座標 $u$ が単元となる。
定理~\ref{thm:allm-concurrent-v33} の横断環から長さ $m+3$ を得る。
\end{proof}

全ての次元はアフィン係数空間での次元である。
Poisson テンソルの倍率や重み付き座標変換群で割った次元ではない。
この節は全ての既約成分の分類も、成分の最大次元の決定も主張しない。

\subsection{対数型の枝と二つの primary 成分}

前節の局所模型で $b=u=0$、$a\ne0$ とすると、$c$ は三つの非共点直線の積となる。
混合項は次数 $m$ の $w$ 移動で消去でき、三角形座標では基本一形式が
\[
 XYZ w\left(\lambda_0\frac{dw}{w}
 +\lambda_1\frac{dX}{X}
 +\lambda_2\frac{dY}{Y}
 +\lambda_3\frac{dZ}{Z}\right),
 \qquad m\lambda_0+\lambda_1+\lambda_2+\lambda_3=0
\]
となる。従ってこの枝は対数型 $L(1,1,1,m)$ の閉包の局所枝である。
逆に、重み $m$ の因子の $w$ 係数が非零なら、その因子は $w+g_m$ と書けるので、
一般の対数型テンソルもこの表示に含まれる。
この対数型成分の次元も $N_m+9$ であり、一般点では被約である。

滑らかな因子を省き、$n=m+2$ と置く。局所模型の関係イデアルには
正確な primary 分解
\[
 (ab,b^2,a^n u)=(a^n,ab,b^2)\cap(b,u)
\]
がある。
第一のイデアルは $(a,b)$-primary、第二は素イデアルであり、embedded prime はない。
二つの被約枝の交叉は滑らかだが、引戻し側の primary 成分と対数型枝の
スキーム論的交叉は
\[
 k[[a,b,u]]/(a^n,b,u)\simeq k[a]/(a^n)
\]
である。従って、滑らかな被約交叉の上に長さ $m+2$ の厚みを持つ。

さらに Hilbert--Burch 行列
\[
 \begin{pmatrix}
 b&a^{n-1}u\\
 -a&0\\
 0&-b
 \end{pmatrix}
\]
の最大小行列式は、符号を除いて $ab,b^2,a^nu$ である。
関係イデアルの高さは二なので、滑らかな因子を除いた一次元局所環は
Cohen--Macaulay である。一方、三つの生成元は最小なので完全交叉ではない。
一般の横断超平面を $u=-a-\beta b$ と取ると、切断環は
\[
 k[a,b]/(a^{n+1},ab,b^2)
\]
となり、長さは $n+2=m+4$ である。
これは二枝の一般点での長さ $(m+3)+1$ とも一致する。

\begin{remark}[原点の形式変形問題との接続]
ここでの非被約な長さは、Poisson スキームの非零点で計算されたものである。
原点で完備化した形式変形基底に対する帰結には、斉次錐の既約成分とその一般点の
長さが原点完備化により保たれることを別に用いる。
その議論と外側の例外列代数の二次 miniversal hull の同定を組み合わせることで、
同じ長さ $m+3$ が圏の変形基底にも現れる。
この接続は、任意の Poisson 量子化が AS 正則であるという主張を必要としない。
\end{remark}

\section{引戻し族の明示的な4-Calabi--Yau量子化}
\label{sec:pullback-ore-v33}

本節では、引戻し成分の被約族に属するPoisson構造を具体的な座標環で量子化する。
構成の基本原理は、通常の $\mathbb P^3$ の引戻し成分に対する
Pymの微分Ore拡大による量子化\cite[\S3.5]{Pym}と共通する。
ここで必要なのは、重みに依存する係数を求め、発散零の条件と
標準例外列による実現を全ての重みについて確かめることである。

\subsection{発散零代表の分解}

$m\ge3$ とし、$x,y,z$ の重みを一、$w$ の重みを $m$ とする。
二変数三次式 $c\in\mathbb C[x,y]_3$、一次式 $L\in\mathbb C[x,y]_1$、
および $f\in\mathbb C[x,y]_{m+1}$ を任意に取る。
\[
 H_c=c_y\partial_x-c_x\partial_y,\qquad
 E_{xy}=x\partial_x+y\partial_y,\qquad
 V=\frac13(H_c+LE_{xy})
\]
と置けば、$D=L\partial_z$、$F=f\partial_z$ に対応する代表は
\[
 P=\partial_z\wedge V+f\partial_z\wedge\partial_w
\]
である。$\Omega=dx\wedge dy\wedge dz\wedge dw$ および
$E=x\partial_x+y\partial_y+z\partial_z+mw\partial_w$ に対し、
\[
 \iota_E\iota_P\Omega
 =c\,dw-\frac m3 w\,dc+
       \left(f-\frac m3wL\right)(x\,dy-y\,dx).
\]
右辺は $(x,y,w)$ だけに依存する重み $m+3$ の基本1形式である。
したがって、これは $\mathbb P(1,1,m)$ からの引戻し族の表示である。
基本1形式を $f-wL_{\mathrm{form}}$ という係数で書く規約では、
$L_{\mathrm{form}}=mL/3$ となることに注意する。

\begin{lemma}[引戻し族のunimodular lift]
\label{lem:pullback-unimodular-v33}
次の重み一次の導分を $\mathbb C[x,y,w]$ 上に定める。
\begin{align*}
 \delta_m(x)&=\frac{c_y}{3}+\frac{m}{3(m+3)}xL,\\
 \delta_m(y)&=-\frac{c_x}{3}+\frac{m}{3(m+3)}yL,\\
 \delta_m(w)&=f-\frac{m}{m+3}Lw.
\end{align*}
$P$ の発散零代表は
\[
 \widehat P=\partial_z\wedge\delta_m
\]
である。特に $\operatorname{div}\delta_m=0$ かつ
$[\partial_z,\delta_m]=0$ である。
\end{lemma}

\begin{proof}
重みの和は $m+3$ であり、$\operatorname{div}(f\partial_z)=0$ なので、
Euler場による発散零代表の公式は
\[
 \widehat P=\Pi+\frac1{m+3}E_x\wedge D+
 \left(F-\frac m{m+3}wD\right)\wedge\partial_w
\]
となる。$D=L\partial_z$ を代入し、
$E_x\wedge D=-\partial_z\wedge LE_{xy}$ とすると表示を得る。
また $\operatorname{div}H_c=0$、$\operatorname{div}(LE_{xy})=3L$ より
\[
 \operatorname{div}\delta_m
 =\frac{m}{m+3}L-\frac{m}{m+3}L=0.
\]
各係数は $z$ に依存しないので、二つの導分は可換する。
したがって $[\widehat P,\widehat P]=0$ も直接従う。
\end{proof}

\subsection{座標環と全次数の積}

\begin{theorem}[全重みに対する微分Ore族]
\label{thm:pullback-ore-v33}
上の任意の $c,L,f$ と $\hbar\in\mathbb C$ に対して、
\[
 S_m(c,L,f;\hbar)=\mathbb C[x,y,w][z;\hbar\delta_m]
\]
を取る。すなわち $x,y,w$ は互いに可換であり、残りの関係は
\begin{align*}
 [z,x]&=\hbar\left(\frac{c_y}{3}+\frac{m}{3(m+3)}xL\right),\\
 [z,y]&=\hbar\left(-\frac{c_x}{3}+\frac{m}{3(m+3)}yL\right),\\
 [z,w]&=\hbar\left(f-\frac{m}{m+3}Lw\right)
\end{align*}
である。この代数は以下の性質を持つ。
\begin{enumerate}[label=\textup{(\roman*)}]
\item $x^iy^jw^kz^\ell$ が基底をなし、
\[
 H_{S_m}(t)=\frac1{(1-t)^3(1-t^m)}.
\]
さらに $c,L,f,\hbar$ を係数変数とした普遍族も、この基底を持つ自由加群である。
\item $S_m$ はnoetherian整域、global dimension四、GK次元四のAS正則代数である。
\item $S_m$ は通常の非次数付きの意味で4-Calabi--Yauであり、
そのNakayama自己同型は恒等である。Gorenstein parameterは $m+3$ である。
\item $\hbar$ を形式変数とした族は、$\widehat P$ を半古典極限とする
重み付き次数を保つ平坦形式量子化である。
\end{enumerate}
\end{theorem}

\begin{proof}
$\delta_m$ は可換多項式環の導分であるから、微分Ore拡大の定義により
$1,z,z^2,\ldots$ が左 $\mathbb C[x,y,w]$-基底である。
同じ構成は係数の普遍多項式環上で行えるので、族についても基底は共通である。
$\delta_m$ の重み次数が一であることから、関係は重み付き斉次であり、
表示したHilbert級数を得る。微分Ore拡大はnoetherian性と整域性を保つ。

多項式環 $\mathbb C[x,y,w]$ は三次元Calabi--Yauであり、
Liu--Wang--WuのOre拡大定理\cite[Theorem 3.3]{LWW}より、$S_m$ は
twisted Calabi--Yau次元四である。
微分Ore拡大に対するLiu--Maの公式\cite[Theorem 4.2]{LiuMa2020}は
\[
 \mu_{S_m}|_{\mathbb C[x,y,w]}=\operatorname{id},\qquad
 \mu_{S_m}(z)=z+\hbar\operatorname{div}\delta_m
\]
を与える。補題により発散は零なので、$\mu_{S_m}=\operatorname{id}$ である。
連結次数付きであることからAS正則性が従う
\cite[Lemma 1.3]{LWW}。有限GK次元はHilbert級数から四と分かる。

Gorenstein parameterを $\ell$ とすると、有限次数付き自由分解と
AS--Gorenstein双対から
$H_{S_m}(t^{-1})=t^\ell H_{S_m}(t)$ を得る。
左辺は $t^{m+3}H_{S_m}(t)$ なので $\ell=m+3$ である。
最後に一次交換子は表示した $\partial_z\wedge\delta_m$ と一致し、
全次数の平坦性は共通のOre基底による。
\end{proof}

この量子化では、星積も閉じた式で書ける。
通常の可換多項式を正規順序の基底と同一視すると
\[
 a\star b=\sum_{r\ge0}\frac{\hbar^r}{r!}
              (\partial_z^r a)(\delta_m^r b)
 \qquad(a,b\in\mathbb C[x,y,z,w]).
\]
各対 $(a,b)$ について和は $\deg_z a$ で打ち切られる。
したがって、これは形式級数だけでなく多項式的な族の積を与える。
可換する導分 $\partial_z,\delta_m$ に対し、三重積の両方の括弧付けは
\[
 \mu^{(3)}\exp\!\left(\hbar(
 \partial_z^{(1)}\delta_m^{(2)}+
 \partial_z^{(1)}\delta_m^{(3)}+
 \partial_z^{(2)}\delta_m^{(3)})\right)
\]
であるので、結合性も直接確かめられる。

\begin{corollary}[標準例外列と四対象の許容成分]
\label{cor:pullback-exceptional-v33}
$S=S_m(c,L,f;\hbar)$ に対し
\[
 (\mathcal O_S,\mathcal O_S(1),\ldots,\mathcal O_S(m+2))
\]
は $D^b(\operatorname{qgr}S)$ のfull strong exceptional collectionである。
特に
\[
 B_m(S)=\End\bigl(\mathcal O_S\oplus\mathcal O_S(1)
             \oplus\mathcal O_S(m)\oplus\mathcal O_S(m+1)\bigr)
       \simeq S_{(0,1,m,m+1)}
\]
に対し、$\Perf(B_m(S))$ はその許容部分圏である。
この有限代数の次元は $2m^2+6m+20$ であり、$m=3$ では五十六である。
\end{corollary}

\begin{proof}
定理のnoetherian性、有限global dimension、Gorenstein parameterに
Orlovの結果\cite[Corollary 2.7]{Orlov2009}を適用する。
strong性によるMorita同値と例外部分列の許容性から四対象の主張を得る。
次元については、四つの恒等射と六つの上三角成分のHilbert関数を足せば
\[
 4+3+\left(\binom{m+2}{2}+1\right)
 +\left(\binom{m+3}{2}+3\right)
 +\binom{m+1}{2}+\left(\binom{m+2}{2}+1\right)+3
 =2m^2+6m+20
\]
となる。残る $m-1$ 個の例外対象をmutationすれば、この四対象が生成する
圏を一つの成分にまとめた半直交分解を得る。
\end{proof}

\begin{remark}[量子化の範囲]
代数の構成と定理の\textup{(i)--(iv)}自体は $m\ge1$ でも有効である。
本稿のPoisson成分に関する一般点の局所環の主張は $m\ge3$ を対象とする。
一次形式の二平面と重み $m$ の座標の移動を用いて適合座標を取れば、
引戻し族の一般の点は上の表示に入り、従って実際のAS正則座標環による量子化を持つ。
これは被約引戻し族における存在の主張であり、全ての量子化の分類ではない。
また、Poissonスキームの一般点のArtin環と、$\hbar=1$ に特殊化した
AS正則代数の変形空間のArtin環を同一視する主張は含めない。
\end{remark}

\section{接線共鳴の全重み系列と任意の一般点長さ}
\label{sec:tangent-resonance-allm-v34}

整数 $m\geq3$ を固定し、$X_m=\mathbb P_{\mathrm{stack}}(1,1,1,m)$、
$S_d=\mathbb C[x,y,z]_d$、$n_m=\binom{m+2}{2}$ とする。
Poisson スキーム $\mathscr P_m$ の座標と方程式を
\[
 (c,D,[F])\in S_3\oplus\mathfrak{sl}_3\oplus W_m,
 \qquad W_m=S_{m+1}^{\oplus3}/S_m\mathbf x,
\]
\begin{equation}
 D(c)=0,\qquad
 M_m(c,D;F)=\frac{m+3}{3}F(c)
 -\frac m3(\mathbf x\times D)\cdot F-c\operatorname{div}F=0
 \label{eq:resonance-mixed-allm-v34}
\end{equation}
と書く。ここで $\mathbf x=(x,y,z)$、$\dim W_m=(m+2)(m+4)$ であり、
方程式のイデアルは根基化しない。

\begin{theorem}[接線共鳴系列の全形式近傍]
\label{thm:tangent-resonance-allm-v34}
$j\geq0$、$r=m-2j>0$ とし、
\[
 q=xz-y^2,\qquad c_*=xq,\qquad N=x\partial_y+2y\partial_z,
 \qquad a_j=-\frac{4m+3-6j}{2m},
\]
\[
 D_j=a_jN,\qquad
 F_j=q^j(0,\tfrac12yz^r,z^{r+1})
\]
と置く。次が成り立つ。
\begin{enumerate}
\item $(c_*,D_j)$ の $\operatorname{SL}_3$ 軌道上で全混合核を動かした族の閉包
 $\mathscr E_{m,j}$ は、$\mathscr P_m$ の既約成分である。
 その被約成分の次元は $n_m+8$、一般点の長さは $j+1$ であり、
 被約成分には重み付き座標変換群の稠密な一軌道がある。
\item $P_j=(c_*,D_j,0)$ における全完備局所環は
 \begin{equation}
 \boxed{\quad
 \widehat{\mathcal O}_{\mathscr P_m,P_j}
 \simeq
 \frac{\mathbb C[[s_1,\ldots,s_{n_m+7},\kappa,\gamma,u]]}
 {(\kappa u,\gamma^{j+1}u)}.
 \quad}
 \label{eq:tangent-resonance-origin-v34}
 \end{equation}
 $u=0$ の枝は $n_m+9$ 次元の対数型成分 $L(1,2,m)$ であり、
 他方は $\mathscr E_{m,j}$ の厚みを記述する。
\item $Q_j=(c_*,D_j,[F_j])$ では
 \begin{equation}
 \widehat{\mathcal O}_{\mathscr P_m,Q_j}
 \simeq
 \frac{\mathbb C[[t_1,\ldots,t_{n_m+8},\gamma]]}{(\gamma^{j+1})}.
 \label{eq:tangent-resonance-nonzero-v34}
 \end{equation}
 従って接空間の次元は $j=0$ なら $n_m+8$、$j\geq1$ なら $n_m+9$ である。
\end{enumerate}
固定した $m$ に対し、これらは相異なる $\lceil m/2\rceil$ 個の成分であり、
一般点の長さは $1,2,\ldots,\lceil m/2\rceil$ を尽くす。
\end{theorem}

この主張は有限個の $m,j$ における階数からの外挿ではない。
以下の証明では、一つの二重対角行列、二変数多項式の次数降下、
四つの余核重量、および二つの特異点での評価によって、全ての整数範囲を扱う。
仮定 $r>0$ は本質的に用いる。$r=0$ は本定理に含めない。

\subsection{基底スキームの二変数切片}

\begin{lemma}
\label{lem:resonance-base-slice-v34}
$a\ne0$ とする。$\mathscr Z=\{(c,D):D(c)=0\}$ は $(xq,aN)$ で滑らかな九次元であり、
七つの軌道方向を除く形式座標には
\begin{equation}
 q_\gamma=xz-y^2+\gamma z^2,\qquad
 c_\gamma=xq_\gamma,\qquad
 D_{a,\gamma}=a\bigl((x+2\gamma z)\partial_y+2y\partial_z\bigr)
 \label{eq:resonance-base-slice-v34}
\end{equation}
を取れる。
\end{lemma}

\begin{proof}
$x^2z-xy^2$ の $\mathfrak{sl}_3$ 接軌道は
\[
 x^3,\ x^2y,\ x^2z,\ xy^2,\ xyz,\ y^3,\ 2xz^2-y^2z
\]
の七次元空間で、安定化 Lie 代数は $\mathbb C N$ である。
$N$ は $aN$ と可換なので、組の軌道も七次元となる。
三次形式の横断方向を
$\dot c=Rz^3+Syz^2+Txz^2$、$\dot D=(h_{ij})\mathbf x$、
$h_{33}=-h_{11}-h_{22}$ とする。
線形化 $aN(\dot c)+\dot D(c_*)=0$ の $yz^2,xz^2,y^2z$ 係数は
\[
 6aR,\qquad aS+2h_{13},\qquad4aS-h_{13}
\]
なので $R=S=h_{13}=0$ である。残る係数から
\[
 h_{12}=h_{11}=h_{22}=h_{31}=0,\qquad
 h_{23}=2aT,\qquad h_{32}=2h_{21}
\]
を得る。横断接空間は正確に
$\dot c=Txz^2$、$\dot D=\lambda N+2aTz\partial_y$ の二次元である。
式~\eqref{eq:resonance-base-slice-v34} は恒等的に $D(c)=0$ を満たすので、
七軌道座標と $a,\gamma$ の写像は九次元接空間への同型を与える。
局所次元と接空間次元が共に九であることから正則性が従い、形式逆写像定理を適用できる。
七軌道座標は群の局所的な代数的切片で選べるので、この写像は対象点の近傍でエタールである。
\end{proof}

\subsection{混合核の共鳴と実際の Hamiltonian 直和因子}

以下、
\[
 T=2x\partial_x+y\partial_y,\quad
 L=y\partial_y+2z\partial_z,\quad
 T+L=2E_{\mathbf x},\quad[T,N]=N,
 \qquad\beta=m+\tfrac32-j
\]
と置く。

\begin{lemma}[共鳴混合核]
\label{lem:resonance-kernel-v34}
$(c_*,D_j)$ における混合核の次元は $n_m+1$ である。
次数 $m$ の $w$ 移動の像は $n_m$ 次元であり、その商は $[F_j]$ で生成される。
\end{lemma}

\begin{proof}
混合式の $x=0$ への制限は $r y^2F_x$ である。
$r>0$ なので $F_x$ は $x$ で割れ、Euler 同値を一意に用いて $F_x=0$ とできる。
$F=(0,G,H)$ とすると、混合式は
\[
 q(G_y+H_z)=\beta(xH-2yG)
\]
である。$q$ は既約、$\beta\ne0$ なので $xH-2yG=2qB$ と書ける。
$x=0$ での制限から $G-yB$ が $x$ で割れ、従って一意に
\[
 G=xA+yB,\qquad H=2yA+2zB,\qquad A,B\in S_m
\]
と書ける。この表示で混合条件は正確に
\begin{equation}
 N(A)=(T-2j)B
 \label{eq:resonance-operator-v34}
\end{equation}
となる。

$T$ 固有値 $k$ の部分を $S_m[k]$ と書く。
$k\ne2j$ では $B_k$ が $A_{k-1}$ から一意に決まる。
問題となるブロックの基底は
\[
 A_i=x^iy^{2j-1-2i}z^{r+1+i}\quad(0\leq i\leq j-1),
 \qquad
 B_i=x^iy^{2j-2i}z^{r+i}\quad(0\leq i\leq j)
\]
であり、直接計算により
\begin{equation}
 N(A_i)=2(r+1+i)B_i+(2j-1-2i)B_{i+1}.
 \label{eq:resonance-bidiagonal-v34}
\end{equation}
最初の $j$ 行の対角係数は全て非零なので、この $N$ は単射である。
従って $A_{2j-1}=0$ が必要十分で、$A$ の自由度は $n_m-j$、
$B_{2j}$ の自由度は $j+1$ である。混合核の次元は $n_m+1$ となる。
$j=0$ では空の源から一次元の標的への写像と解釈する。

この二重対角写像の余核は $Q=q^jz^r$ の類で生成される。
非零性を明示するため、$B_i$ に対する線形汎関数の値を
\[
 \ell(B_i)=(-1)^i\prod_{h=0}^{i-1}
       \frac{2(r+1+h)}{2j-1-2h}
\]
と置く。式~\eqref{eq:resonance-bidiagonal-v34} の像では $\ell$ は零となる。
一方、$Q=\sum_{i=0}^j(-1)^{j-i}\binom ji B_i$ なので
\[
 \ell(Q)=(-1)^j\sum_{i=0}^j\binom ji
      \prod_{h=0}^{i-1}\frac{2(r+1+h)}{2j-1-2h}\ne0.
\]
最後の和の全ての項は正の有理数である。
従って、ここで用いる非零性は一般点での階数という仮定を含まない。

次に次数 $m$ の $w$ 移動 $g$ を考える。
$p=(\nabla c_*-\mathbf x\times D_j)/3$ とすると、その混合項は
$-p\times\nabla g$ であり、第一成分を Euler 同値で消した表示は
\[
 (G,H)=\bigl(qg_z-(\beta-1)xg,
             -qg_y-2(\beta-1)yg\bigr).
\]
対応する $(A,B)$ は
\begin{equation}
 A=-\tfrac12(T+1-2j)g,\qquad B=-\tfrac12N(g).
 \label{eq:resonance-hamiltonian-v34}
\end{equation}
第一式の核は $S_m[2j-1]$ だけであり、そこでは第二式が単射である。
従ってこの移動の像は正確に $n_m$ 次元となる。
他の全ての $A$ 方向は第一式で尽くされ、残る $B_{2j}$ の商は上で求めた一次元余核である。
その代表 $A=0,B=Q/2$ が $F_j$ となる。
\end{proof}

\subsection{接触を解いたときに余分な核が消えること}

\begin{lemma}[横断する直線・二次曲線での全解]
\label{lem:resonance-transverse-kernel-v34}
$a=a_j$ を固定し、$\gamma\ne0$ とする。
式~\eqref{eq:resonance-base-slice-v34} の混合核は、全て次数 $m$ の $w$ 移動から来る。
その次元は $n_m$ である。
\end{lemma}

\begin{proof}
$x=0$ での制限は $r(y^2-\gamma z^2)F_x$ なので、再び $F_x=0$ とできる。
$x=1$ へ非斉次化し、$Q=z-y^2+\gamma z^2$、$V=(G,H)$ と書く。
$G,H$ の次数は高々 $m+1$ であり、混合式は
\begin{equation}
 Q\operatorname{div}V=\beta V(Q)
 \label{eq:resonance-affine-log-v34}
\end{equation}
となる。$b_0=\beta-1$ として
\[
 K_{b_0}(g)=(Qg_z-b_0Q_zg,\ -Qg_y+b_0Q_yg)
\]
と置けば、任意の多項式 $g$ について
$Q\operatorname{div}K_{b_0}(g)=\beta K_{b_0}(g)(Q)$ が恒等的に成り立つ。

$\gamma\ne0$ なので $Q$ の最高次二次部分は非退化であり、線形変数変換の後に $uv$ とできる。
非零解 $V$ の最高次数を $d$ とする。
次数 $d+1$ の式の境界係数から $G_d$ の $v^d$ 係数と $H_d$ の $u^d$ 係数は零となる。
内部の単項式 $u^av^b$、$a+b=d+1$、$a,b\geq1$ については、
$G_d$ の $u^av^{b-1}$ 係数と $H_d$ の $u^{a-1}v^b$ 係数に対し
\[
 (a-\beta)G_{a,b-1}+(b-\beta)H_{a-1,b}=0
\]
が課される。$\beta$ は半整数なので、二つの係数は共に非零である。
従ってこの二係数の解空間は
$(b-\beta,\beta-a)$ で生成される一次元であり、これは
$K_{b_0}(u^{a-1}v^{b-1})$ の最高次部分である。

よって $d\geq1$ なら、次数 $d-1$ の斉次多項式 $g_{d-1}$ を一意に選んで、
$V-K_{b_0}(g_{d-1})$ の次数を下げられる。
$d=0$ の解は非退化二次部分の式から零だけである。
次数を順に下げれば、全解は $V=K_{b_0}(g)$、$\deg g\leq m$ と書ける。
同じ最高次計算から $K_{b_0}$ はこの多項式空間上で単射である。
$g$ を次数 $m$ に斉次化すると、これは正確に次数 $m$ の $w$ 移動の式となる。
従って全混合核の次元は $n_m$ である。
\end{proof}

\subsection{四つの余核重量と Schur 残差の全形}

$H_0=x\partial_x-z\partial_z$ の重量を用いる。
$\gamma=0$ では混合写像の重量は $+1$、$\gamma$ の係数である混合写像の重量は $-1$ である。
$F_j$ の重量は $-r=2j-m$ である。

\begin{lemma}[余核の重量]
\label{lem:resonance-cokernel-weights-v34}
$(c_*,D_j)$ の混合余核は四次元であり、その重量は
\[
 -m-3,\qquad -m-2,\qquad -m-1,\qquad 1-r
\]
で、それぞれの多重度は一である。
\end{lemma}

\begin{proof}
$h_d(t)=\operatorname{ch}_{H_0}S_d$ とすると
\[
 h_d(t)=t^{-d}\frac{(1-t^{d+1})(1-t^{d+2})}{(1-t)(1-t^2)},
 \qquad
 \operatorname{ch}W_m=(t^{-1}+1+t)h_{m+1}(t)-h_m(t).
\]
Hamiltonian 写像の重量は $+1$ なので、前補題で求めた核の指標は
$t h_m(t)+t^{-r}$ となる。
混合写像の重量も $+1$ であることから、余核の指標は
\[
 h_{m+3}-(1+t+t^2)h_{m+1}+(t+t^2)h_m+t^{1-r}
 =t^{-m-3}(1+t+t^2)+t^{1-r}
\]
となる。最後の等式は表示した有理式への代入による恒等式である。
\end{proof}

基底を $\mathbb C[[a-a_j,\gamma]]$ とし、特殊点で重量を保つ混合写像の可逆小行列を固定する。
その源の部分空間を $U$、標的の選択した座標空間を $V$ と書く。
特殊点で
\[
 W_m=U\oplus\operatorname{Ham}(S_m)\oplus\mathbb C F_j
\]
となるように選ぶ。
基底の全域で実際の Hamiltonian 列を用い、$U$ と $F_j$ は固定する。
この連結した列の行列式は特殊点で非零なので、局所環上で可逆である。
従って Hamiltonian 像は非被約な基底上でも自由な直和因子であり、
それを除いた後に $F_j$ 以外の未知列が残ることはない。

選択した正方小行列は
$P_0(a)+\gamma P_1(a)$ と書け、$P_0(a)$ は局所環上で可逆である。
$M=M_m(c_\gamma,D_{a,\gamma};-)$、$P=P_0(a)+\gamma P_1(a)$ と略記すると、
補正した源は具体的に
\[
 \widehat F_j=F_j-U P^{-1}\operatorname{pr}_V M(F_j)
\]
である。ここで $U$ は選んだ列の包含行列も表す。
Schur 残差は実際の混合像 $M(\widehat F_j)$ であり、その $V$ 座標は全て零である。
$P_0(a)^{-1}P_1(a)$ は源の重量を二つ下げるため冪零である。
従って逆行列は $\gamma$ の有限多項式であり、
$F_j$ の消去後の列において、重量 $k$ の標的係数に現れ得る $\gamma$ の次数は正確に
\[
 \ell=\frac{1-r-k}{2}\in\mathbb Z_{\geq0}
\]
に限られる。

最低三重量のうち、$-m-3,-m-2$ の標的はそれぞれ
$z^{m+3},yz^{m+2}$ の一次元である。
重量 $-m-1$ の標的は
$xz^{m+2},y^2z^{m+1}$ の二次元だが、対応する最低重量の源は
$z^{m+1}\partial_x$ のみであり、
\begin{equation}
 M_m(c_*,D_j;z^{m+1}\partial_x)
 =\bigl((j+\tfrac32)xz+r y^2\bigr)z^{m+1}.
 \label{eq:resonance-lowest-column-v34}
\end{equation}
従って $xz^{m+2}$ を可逆小行列の行、$y^2z^{m+1}$ を自由行として選べる。
重量 $1-r$ でも一つの自由行を選ぶ。
四つの自由行上の Schur 残差は、適切な順で正確に
\begin{equation}
 f(a),\qquad \gamma^{j+1}g(a),\qquad0,\qquad\gamma^{j+2}h(a)
 \label{eq:resonance-schur-weight-shape-v34}
\end{equation}
となる。ここで $f,g,h$ は $a_j$ で正則な有理関数である。
第三成分が零なのは対応する指数が $j+3/2$ で整数でないためである。
この議論は最初の可能次数だけでなく、各成分に他の高次項が一切ないことも示す。

\subsection{先頭係数の非零性と全局所環}

まず $f(a_j)=0$ であり、その微分は非零である。
実際、固定した $F_j$ に対する $a$ 微分は
\[
 \Omega_a=-\frac m3xq^{j+1}z^r.
\]
これが混合像に入ると仮定する。
右辺が $x$ で割れるため $F_x=0$ とでき、さらに $xq$ で割れるため
前と同じ表示 $G=xA+yB,H=2yA+2zB$ が可能である。
すると混合式の像は
\[
 xq\{(T-2j)B-N(A)\}
\]
となる。しかし重量 $2j$ の余核で $q^jz^r$ は非零と既に証明した。
従って $\Omega_a$ の余核類は非零であり、$f'(a_j)\ne0$ となる。
特に $f(a)$ は $a-a_j$ の単元倍である。

次に $a=a_j$ に固定する。このとき $f=0$ なので、残る標的多項式は
\[
 \gamma^{j+1}g(a_j)y^2z^{m+1}
 +\gamma^{j+2}h(a_j)z^{m+3}
\]
である。$\gamma\ne0$ では三次式 $xq_\gamma$ の二つの特異点は
$[0,\sqrt\gamma,1]$、$[0,-\sqrt\gamma,1]$ である。
ここでは $c=dc=0$ であり、$D_{a_j,\gamma}$ は径方向のベクトルに比例する。
従って、任意の源 $F$ に対する混合像がこの二点で零になる。
上の残差を評価すると
\[
 h(a_j)=-g(a_j)
\]
を得る。

もし $g(a_j)=0$ なら、固定した $a=a_j$ 上の全 Schur 残差が恒等的に零となる。
その場合 $\gamma\ne0$ の混合核にも Hamiltonian 像に加えて一方向が残り、
補題~\ref{lem:resonance-transverse-kernel-v34} の核次元 $n_m$ に矛盾する。
従って $g(a_j)\ne0$ である。
以上により、四つの Schur 係数のイデアルは正確に
\begin{equation}
 (f(a),\gamma^{j+1}g(a),\gamma^{j+2}h(a))
 =(a-a_j,\gamma^{j+1}).
 \label{eq:resonance-schur-ideal-v34}
\end{equation}

残った源の座標を $u$ とすると、全関係は式~\eqref{eq:resonance-schur-ideal-v34} に $u$ を掛けたものである。
七つの軌道変数と $n_m$ 個の実際の $w$ 移動を戻し、$\kappa=a-a_j$ とすれば
式~\eqref{eq:tangent-resonance-origin-v34} を得る。
全ての消去は基底変数だけに依存し、$u$ に関しては線形である。
従って同じ変換を $u=1$ の点でも用いることができ、
式~\eqref{eq:tangent-resonance-nonzero-v34} が従う。

\subsection{大域成分・primary 交叉・一般点の長さ}

\begin{proof}[定理~\ref{thm:tangent-resonance-allm-v34} の残りの証明]
\[
 (\kappa u,\gamma^{j+1}u)
 =(u)\cap(\kappa,\gamma^{j+1})
\]
は primary 分解である。
二つの被約枝の次元は $n_m+9,n_m+8$ であり、被約交叉は滑らかな $n_m+7$ 次元である。
primary 成分同士の交叉は、この滑らかな空間上の
$\mathbb C[[\gamma]]/(\gamma^{j+1})$ という長さ $j+1$ の厚みを持つ。

$(c_*,D_j)$ の軌道は既約七次元であり、その上の混合核は階数 $n_m+1$ のベクトル束である。
従ってその全空間は既約 $n_m+8$ 次元で、上の小さい方の局所枝を覆う。
その閉包 $\mathscr E_{m,j}$ は全 Poisson スキームの既約成分である。
混合核は実際の $w$ 移動と $uF_j$ の直和なので、$u\ne0$ では移動を消し、
$w$ の非零倍率によって $u$ を一にできる。
従って被約成分のこの部分全体は一つの座標変換軌道である。

一般点の長さは式~\eqref{eq:tangent-resonance-nonzero-v34} から $j+1$ と分かる。
より直接には、基底切片はエタールであり、混合行列に対する消去は正則な有理係数の可逆操作である。
従って同じ局所模型はエタール局所的に成立する。
小さい方の枝の一般点で $u$ を可逆にし、滑らかな変数の分数体を $K$ とすれば、
局所 Artin 環は $K[\gamma]/(\gamma^{j+1})$ となる。
ここでは完備化だけから長さを推定していない。

$u=0$ の枝では全混合項を $w$ 移動で消せる。
$\gamma\ne0$ なら直線と二次曲線は横断的に交わり、対応する水平一形式は
\[
 xq_\gamma w\left\{\frac{dw}{w}
 -\frac m3(1+2a)\frac{dx}{x}
 +\frac m3(a-1)\frac{dq_\gamma}{q_\gamma}\right\}.
\]
従ってこの枝は対数型 $L(1,2,m)$ に含まれる。
対数型を与える既約なパラメータ空間の次元は
\[
 \dim\mathbb P(S_1)+\dim\mathbb P(S_2)
 +\dim\mathbb P(S_m\oplus\mathbb Cw)+2
 =2+5+n_m+2=n_m+9
\]
であるから、その像の閉包の次元は高々 $n_m+9$ である。
一方、上で得た同じ次元の局所枝を含むため、次元は正確に $n_m+9$ となる。
従ってこの既約な閉包が大きい方の成分である。

許される整数は $j=0,\ldots,\lfloor(m-1)/2\rfloor$ である。
対応する一般点の長さ $j+1$ は互いに異なり、一般点の長さは既約成分に固有なので、
成分 $\mathscr E_{m,j}$ も互いに異なる。
これで全ての主張が証明された。
\end{proof}

\subsection{発散零の持ち上げとアフィン量子化への接続}

\begin{proposition}
\label{prop:resonance-affine-lift-v34}
$E=x\partial_x+y\partial_y+z\partial_z+mw\partial_w$ とし、
\[
 H'_r=x\partial_x-z\partial_z-(r+1)w\partial_w,\qquad
 H_{m,r}=\frac12\left(H'_r+\frac r{m+3}E\right),
 \qquad \mathcal N_j=N-q^jz^r\partial_w
\]
と置く。点 $(c_*,D_j,[F_j])$ の一意な発散零持ち上げは
\[
 \widehat P_{m,j}=H_{m,r}\wedge\mathcal N_j
\]
である。また
\[
 [H_{m,r},\mathcal N_j]=\tfrac12\mathcal N_j,\qquad
 \operatorname{div}H_{m,r}=-\tfrac12,\qquad
 \operatorname{div}\mathcal N_j=0.
\]
従ってこの点は、後のアフィン量子化で
$X_\alpha=-2H_{m,r}$、$\alpha=-(m+r+3)/(m+3)$ とするパラメータに対応する。
\end{proposition}

\begin{proof}
$q$ は $H'_r$ 不変であり、$\mathcal N_j$ の $H'_r$ 重量は一、$E$ 重量は零である。
従って Lie 括弧と二つの発散の式が従う。
標準形からの持ち上げは
\[
 \Pi+\frac1{m+3}E_{\mathbf x}\wedge D_j+
 \left(F_j-\frac{\operatorname{div}F_j}{m+3}\mathbf x
                   -\frac m{m+3}wD_j\right)\wedge\partial_w,
 \qquad p=\frac{\nabla c_*-\mathbf x\times D_j}{3}
\]
である。
$\operatorname{div}F_j=(m+3/2-j)q^jz^r$ を代入すると、
この六係数は $H_{m,r}\wedge\mathcal N_j$ の六係数と一致する。
$X_\alpha$ の四つの固有値は
$\alpha,\alpha+1,\alpha+2,m\alpha+m+r+1$ となり、最後の主張も従う。
\end{proof}

\begin{remark}[重み三および除外した端点]
$m=3$ では $j=0,1$ となり、一般点で滑らかな十八次元成分と、
一般点で長さ二の十八次元成分を得る。
後者は独立な有理係数消去による局所環 $(\kappa u,\gamma^2u)$ と一致する。
実際、独立計算で用いる $F_-=(0,y^3z/2,y^2z^2)$ との間には
$F_{3,1}+5F_-=-\operatorname{Ham}(yz^2)$ という Euler 商上の恒等式があり、
移動と非零倍率によって同じ軌道を表す。
$m$ が偶数の端点 $j=m/2$ では $r=0$ となり、
$x=0$ で $F_x$ を消す議論が成立しないため、この定理の結論を延長してはならない。
一方、この端点の Poisson 場と量子化自体は定義できる。
$J=xw+q^j y$ とし、$\Omega=dx\wedge dy\wedge dz\wedge dw$、
$\alpha=\iota_E\iota_{H_{m,0}\wedge\mathcal N_j}\Omega$ とすれば、
\[
 2q\,dJ-(m+1)J\,dq=2\alpha.
\]
従って端点は通常の有理型 $R(2,m+1)$ に属する。
この式に共通因子はない。実際、左辺の $dw$ 係数は $2xq$ であり、
$x=0$ では $dx$ 係数が $2qw-q^jyz$ として非零、
$q=0$ の一般点では $-(m+1)J\,dq$ が非零だからである。
ここから $r=0$ における成分性や全形式局所環は主張しない。
また本定理は、全 Poisson スキームの既約成分の完全分類を主張しない。
\end{remark}

\section{非線形アフィン量子化とその中心}
\subsection{非線形アフィン族のNakayama自己同型とCalabi--Yau条件}
\label{subsec:affine-nakayama-v34}

CGG積に現れる三変数部分代数はLecoutre--Sierraの代数
$R(2,\alpha)$である。この既知のNakayama公式と、最後の生成元に関する
剛性を組み合わせることで、非線形な四変数族のNakayama自己同型を
全パラメータについて決定できる。以下の $\alpha$ は一次導分の固有値であり、
三次形式の層を記述するパラメータとは区別する。

$j,r\ge0$ を整数、$m=2j+r\ge3$ とし、
$\deg(x,y,z,w)=(1,1,1,m)$ とする。記号
\[
 q=xz-y^2,\qquad N_0=x\partial_y+2y\partial_z
\]
を用い、次の二つの多項式のいずれかを選ぶ。
\[
 P^{(q)}_{j,r}=q^jz^r,\qquad
 P^{(y)}_{j,r}=y^{2j}z^r.
\]
選んだ多項式を $P$ と書き、
\begin{align*}
 N&=N_0-P\partial_w,\\
 X_\alpha&=\alpha x\partial_x+(\alpha+1)y\partial_y
       +(\alpha+2)z\partial_z+(m\alpha+m+r+1)w\partial_w
\end{align*}
と置く。$[N,X_\alpha]=N$ である。$s\in\mathbb C$ に対して
\[
 a\star b=\sum_{\ell\ge0}s^\ell N^\ell(a)
                  \binom{X_\alpha}{\ell}(b)
\]
で定まる代数を $\mathcal S_{j,r}^{(\epsilon)}(\alpha,s)$ と書く。
ここで $\epsilon=q,y$ は $P$ の選択を表す。補助次数
$\deg'(x,y,z,w)=(0,1,2,m+r+1)$ に関して $N$ は次数を一つ下げるので、
積は各多項式対について有限和である。
この積の結合性は $[N,X_\alpha]=N$ に対する
Coll--Gerstenhaber--Giaquintoの普遍変形公式による。
この規約での公式と射影的な適用についてはPym\cite[\S3.6, \S4]{Pym}を参照する。

\begin{theorem}[全パラメータのNakayama公式]
\label{thm:affine-nakayama-v34}
上の二つの族は、全ての $\alpha,s\in\mathbb C$ について
noetherian整域かつglobal dimension四のAS正則代数である。
Hilbert級数とGorenstein parameterは
\[
 \frac1{(1-t)^3(1-t^m)},\qquad m+3
\]
である。$T=1+sN$ と置くと、そのNakayama自己同型は正確に
\[
 \boxed{\ \mu=T^{\kappa},\qquad
       \kappa=(m+3)\alpha+m+r+3=\operatorname{div}X_\alpha-1\ }
\]
である。従って $s\ne0$ のとき、通常の非次数付きの意味で
$4$-Calabi--Yauであるための必要十分条件は
\[
 \boxed{\ \alpha=-\frac{m+r+3}{m+3}\ }.
\]
$s=0$ では可換多項式環なので常に $4$-Calabi--Yauである。
\end{theorem}

この公式は $T^\kappa$ を有限二項和として解釈する。特に
\begin{align*}
 \mu(x)&=x,\\
 \mu(y)&=y+\kappa s x,\\
 \mu(z)&=z+2\kappa s y+\kappa(\kappa-1)s^2x,\\
 \mu(w)&=w-\sum_{\ell\ge1}\binom\kappa\ell s^\ell N_0^{\ell-1}(P).
\end{align*}
最後の和も有限であり、形式変数の有限次数切断ではない。

証明に用いる剛性を先に述べる。
$A_3\subset\mathcal S_{j,r}^{(\epsilon)}(\alpha,s)$ を
$x,y,z$ が生成する部分代数とする。

\begin{lemma}[最後の生成元の剛性]
\label{lem:affine-last-generator-v34}
$s\ne0$ とする。$A_3$ の各元を固定する次数付き自己同型は恒等である。
この主張には $\alpha=0,-1$ も含まれる。
\end{lemma}

\begin{proof}
$d=m+r+1=2(j+r)+1$ と置く。以下の多項式を用いる。
\begin{align*}
 J^{(q)}&=x^{r+1}w+
 q^j\sum_{k=0}^{r}\binom rk\frac{q^{r-k}y^{2k+1}}{2k+1},\\
 J^{(y)}&=x^{r+1}w+
 \sum_{k=0}^{r}\binom rk\frac{q^{r-k}y^{2j+2k+1}}{2j+2k+1}.
\end{align*}
選んだ $P$ に対応するものを $J=x^{r+1}w+J_0$ とする。
二項定理から
\[
 N(J)=0,\qquad X_\alpha(J)=d(\alpha+1)J,\qquad \deg J=d
\]
である。例えば $q$ 型では
$N_0(J_0)=xq^j(q+y^2)^r=x^{r+1}P$ となり、
$y$ 型でも $N_0(J_0)=xy^{2j}(q+y^2)^r$ となる。
さらに
\[
 J_0\bmod x=c_{j,r}^{(\epsilon)}y^d
\]
であり、その係数は
\begin{align*}
 c_{j,r}^{(q)}&=(-1)^{j+r}\int_0^1(1-t^2)^r\,dt,\\
 c_{j,r}^{(y)}&=(-1)^r\int_0^1t^{2j}(1-t^2)^r\,dt
\end{align*}
である。どちらも非零の有理数である。

$\psi$ を補題の自己同型とする。重み $m$ の空間は
$(A_3)_m\oplus\mathbb Cw$ なので
\[
 \psi(w)=uw+b,\qquad u\in\mathbb C^\times,\quad b\in(A_3)_m
\]
と書ける。$g=\psi(J)-uJ$ と置くと、$g\in(A_3)_d$ である。
$N(x)=N(q)=N(J)=0$ なので、$J$ は $x,q$ と可換であり、従って $g$ も
この二つの元と可換である。一方、$p=x,q$ に対して
\[
 h\star p=T^{\lambda_p}(h)p,\qquad p\star h=ph,
 \qquad \lambda_x=\alpha,\quad\lambda_q=2(\alpha+1).
\]
右辺の積は通常の可換積である。
$\lambda_x$ と $\lambda_q$ は同時には零にならない。
局所冪零な $N$ に対して
$T^\lambda-1=s\lambda N(1+\text{$N$の正次数項})$ であることから、
非零の固有値を持つ方の元との可換性は $N(g)=0$ を与える。

また $N(y)=x$、$N^2(y)=0$ より
\[
 [y,J]_\star=s\,d(\alpha+1)xJ.
\]
これに $\psi$ を適用して $u$ 倍の元の等式を引くと
$[y,g]_\star=s\,d(\alpha+1)xg$ を得る。
$N(g)=0$ なので $X_\alpha(g)=d(\alpha+1)g$ が従う。
$A_3$ 上では
\[
 X_\alpha=(\alpha+1)(x\partial_x+y\partial_y+z\partial_z)
              -(x\partial_x-z\partial_z)
\]
である。従って $g$ は次数 $d$ の多項式として
$(x\partial_x-z\partial_z)g=0$ を満たす。
$n=j+r$ とすると
\[
 g=\sum_{i=0}^{n}g_i x^i y^{2n+1-2i}z^i.
\]
この空間から $N_0$ を作用させ、標的の基底を
$x^{i+1}y^{2n-2i}z^i$ ($0\le i\le n$) と取ると、
行列は対角成分 $2n+1-2i$、直上成分 $2(i+1)$ の上二重対角行列である。
その行列式は $(2n+1)!!\ne0$ なので、$N(g)=N_0(g)=0$ から $g=0$ である。

最後に、$x$ は $N$ 不変なので $x^{r+1}\star w=x^{r+1}w$ であり、
$g=0$ は通常の多項式として
\[
 x^{r+1}b+(1-u)J_0=0
\]
となる。$x=0$ と置き $c_{j,r}^{(\epsilon)}\ne0$ を使うと $u=1$、
次いで $b=0$ を得る。
\end{proof}

\begin{proof}[定理の証明]
$T^t$ は全ての $t\in\mathbb C$ に対して星積の自己同型である。
これは $[N,X_\alpha]=N$ による捩れたLeibniz則から $T$ の自己同型性を得て、
各入力上で有限になる二項式を $t$ に関して補間すれば従う。
変数を $x,y,z,w$ の順に付け加えると、各 $N(v)$ は先行変数に属する。
従って星積は反復Ore拡大となり、最後の段は
\[
 \mathcal S_{j,r}^{(\epsilon)}(\alpha,s)
       =A_3[w;\tau,\delta],\qquad
 \tau=T^{-(m\alpha+m+r+1)}|_{A_3}
\]
と書ける。各段の先行変数の部分代数を $A_{<v}$、$X_\alpha(v)=\chi_vv$ と書くと
\[
 a\star v=T^{\chi_v}(a)v,\qquad
 \delta_v(a):=v\star a-T^{-\chi_v}(a)\star v\in A_{<v}.
\]
後者は $N^k(v)\in A_{<v}$ ($k\ge1$) と
$T^{-\chi_v}(a)\star v=av$ から従う。結合性により $\delta_v$ は
$T^{-\chi_v}$-導分である。
$v$ の星積冪は通常の $v^n$ に関してmonicな三角変換をなし、
左係数には可逆な作用 $T^{n\chi_v}$ が現れる。
従って各段は左Ore基底を持ち、表示した反復Ore拡大は全代数に一致する。
これとLiu--Wang--Wuの定理
\cite[Theorem 3.3 and Lemma 1.3]{LWW}から、
noetherian性、整域性およびAS正則性が従う。
Hilbert級数は元の多項式環のものと同じであり、その相反性から
Gorenstein parameterは $m+3$ である。

$s\ne0$ とする。変数を $x,y/s,z/(2s^2)$ と取り直すと、
$A_3$ はLecoutre--Sierraの $R(2,\alpha)$ と同じCGG積を持つ。
両者のNakayamaの規約は
$\operatorname{Ext}^d_{A^e}(A,A^e)\simeq {}_1A_\mu$ で一致する。
その全パラメータ公式\cite[Theorem 4.16]{LecoutreSierra2019}から
\[
 \mu_{A_3}=T^{3\alpha+2}
\]
を得る。従ってOre拡大の制限公式\cite[Theorem 3.3]{LWW}は
\[
 \mu|_{A_3}=\tau^{-1}\mu_{A_3}
       =T^{m\alpha+m+r+1+3\alpha+2}=T^\kappa
\]
を与える。
連結次数付きAS正則整域のNakayama自己同型は次数を保つ。
その可逆元は非零定数だけなので、内自己同型による選択の曖昧さもない。
従って $T^{-\kappa}\mu$ は補題の仮定を満たし、恒等である。
これで $\mu=T^\kappa$ が全ての生成元上で証明された。

さらに $T^\kappa(y)=y+\kappa sx$ なので、$s\ne0$ のとき
$\mu=\mathrm{id}$ と $\kappa=0$ は同値である。
内自己同型は恒等に限るから、これが通常のCalabi--Yau条件である。
$s=0$ の場合は可換多項式環として直接従う。
\end{proof}

\begin{corollary}[次数付き自己同型群]
\label{cor:affine-automorphisms-v34}
$s\ne0$ のとき、全ての $j,r,\alpha$ について
\[
 \operatorname{Aut}_{\mathrm{gr}}
     \mathcal S_{j,r}^{(\epsilon)}(\alpha,s)
       =\{\xi_\lambda T^t:\lambda\in\mathbb C^\times,
                                      \ t\in\mathbb C\}
       \simeq\mathbb C^\times\times(\mathbb C,+),
\]
ただし $\xi_\lambda(x,y,z,w)=(\lambda x,\lambda y,\lambda z,\lambda^m w)$
である。
\end{corollary}

\begin{proof}
$m\ge3$ なので $A_3$ は次数一の空間が生成する内在的な部分代数である。
従って制限は定義され、剛性補題により単射である。
Lecoutre--Sierra\cite[Theorem 4.10]{LecoutreSierra2019}による
$\operatorname{Aut}_{\mathrm{gr}}R(2,\alpha)$ の分類は全ての $\alpha$ に適用でき、
その全自己同型はスカラー作用と $T^t$ の積である。
両作用は上の四変数代数にも延長し、互いに可換である。
従って制限は全射でもある。
\end{proof}

\begin{corollary}[全アフィン族の標準例外列]
\label{cor:affine-exceptional-alljr-v34}
$S=\mathcal S_{j,r}^{(\epsilon)}(\alpha,s)$ とすると、
\[
 (\mathcal O_S,\mathcal O_S(1),\ldots,\mathcal O_S(m+2))
\]
は $D^b(\operatorname{qgr}S)$ のfull strong exceptional collectionである。
従って $\Perf(S_{(0,1,m,m+1)})$ はその許容部分圏であり、
$m=3$ では標準六項列による非可換VK四項列の半直交実現を得る。
この主張はCalabi--Yauパラメータ以外の $\alpha$ にも成り立つ。
\end{corollary}

\begin{proof}
定理のnoetherian性、有限global dimension、Gorenstein parameterに
Orlov\cite[Corollary 2.7]{Orlov2009}を適用する。
四対象の生成する圏の許容性は例外部分列から従う。
\end{proof}

\begin{corollary}[発散零のアフィン量子化]
\label{cor:affine-unimodular-cy-v34}
\[
 H'_r=x\partial_x-z\partial_z-(r+1)w\partial_w,
 \qquad H_{m,r}=\frac12\left(H'_r+\frac r{m+3}E\right)
\]
と置く。ただし $E=x\partial_x+y\partial_y+z\partial_z+mw\partial_w$ である。
$X_\alpha=-2H_{m,r}$ となる固有値
$\alpha=-(m+r+3)/(m+3)$ と $s=\hbar/2$ を取れば、
上の有限CGG積は $H_{m,r}\wedge N$ を量子化する通常の
$4$-Calabi--Yau代数である。
特に $(j,r)=(0,m)$ の族を $S_{\mathrm{aff},m}(\hbar)$ と書く。
この族は通常の4-Calabi--Yauであり、$(j,r)=(1,1)$ の $y$ 型は
$N=x\partial_y+2y\partial_z-y^2z\partial_w$ の場合を与える。
\end{corollary}

\begin{proposition}[Zhang twistと非零量子化係数]
\label{prop:affine-zhang-v34}
$j,r,\epsilon,s$ を固定すると、$\alpha$ の変化は次数付き自己同型
$T^\beta$ による左Zhang twistである。従って全ての $\alpha$ で
次数付き加群圏および $\operatorname{qgr}$ 圏は同値である。
また、固定した $j,r,\epsilon,\alpha$ に対する全ての $s\ne0$ は
次数付き同型な代数を与える。

一方、$m$ を固定し、異なる許容値 $r$ に対してCalabi--Yau値の
$\alpha=-(m+r+3)/(m+3)$ を取ると、その非零量子化は
次数付き代数として相互に非同型である。
\end{proposition}

\begin{proof}
$X_{\alpha+\beta}=X_\alpha+\beta E$、$[E,N]=0$ である。
$b$ が重み付き次数 $\ell$ の斉次元なら、有限二項和のVandermonde恒等式から
\[
 a\star_{\alpha+\beta}b
       =T^{\beta\ell}(a)\star_\alpha b
\]
を得る。これは左Zhang twistの定義そのものであり、通常次数の場合の
Lecoutre--Sierra\cite[Theorem 4.2]{LecoutreSierra2019}と同じ機構である。
左twistは逆自己同型による右twistと次数付き同型である。
Zhangの同値\cite[Theorem 1.1]{Zhang1996}は有限次元加群を保存するので、
$\operatorname{qgr}$ にも降下する。

補助次数によるスカラー作用
\[
 \psi_\lambda(x,y,z,w)=(x,\lambda y,\lambda^2z,\lambda^d w),
 \qquad d=m+r+1
\]
は $\psi_\lambda N\psi_\lambda^{-1}=\lambda^{-1}N$ を満たし、
$X_\alpha$ と可換である。従って $s$ の代数から $s/\lambda$ の代数への
次数付き同型を与える。

最後に、次数一生成部分代数は $R(2,\alpha)$ である。
Lecoutre--Sierra\cite[Theorem 5.1]{LecoutreSierra2019}によれば、
これらが同型となるのは $\alpha$ が等しいときに限る。
固定した $m$ に対して異なる $r$ は異なるCalabi--Yau値を与えるので、
四変数代数も次数付き同型にはならない。
この最後の主張は、異なる $r$ の $\operatorname{qgr}$ 圏の非同値までは意味しない。
\end{proof}

この結論は、一次Poissonテンソルの発散だけから推論したものではない。
Nakayama自己同型の全次数の公式を具体的に求めた結果である。
量子化・中心の定理は $r=0$ も含む一方、先の接線成分定理は $r>0$ に対するものである。

\subsection{四生成元の不変環と中心に現れるDu Val曲面}
\label{subsec:affine-polynomial-center-v34}

前節のCalabi--Yau族について、斉次座標環自身の中心を決定する。
三角型の単項式導分の核を四生成元で記述する一般的な方法は
Maubach\cite{Maubach2000}による。本節では必要な生成元・関係を直接証明し、
その零重量部分を求める。この最後の計算で、奇数重みの端の族に
$D_{m+4}$ 型のDu Val曲面が現れる。

以下では $j,r\ge0$、$m=2j+r\ge3$、$d=m+r+1=2(j+r)+1$ とし、
まず $q$ 型の導分
\[
 N=N_0-q^jz^r\partial_w,\qquad N_0=x\partial_y+2y\partial_z,
 \qquad q=xz-y^2
\]
を扱う。記号を
\[
 c_r=(-1)^r\frac{(2r)!!}{(2r+1)!!},\qquad
 J=x^{r+1}w+q^j\sum_{k=0}^r\binom rk
                   \frac{q^{r-k}y^{2k+1}}{2k+1},\qquad
 \Delta=\frac{J^2+c_r^2q^d}{x^{r+1}}
\]
と定める。$0!!=(-1)!!=1$ の規約を用いる。

\begin{proposition}[多項式不変環の全体]
\label{prop:affine-invariant-ring-v34}
$A=\mathbb C[x,y,z,w]$ とすると、$\Delta$ は多項式であり、
\[
 \boxed{
 A^N=\mathbb C[x,q,J,\Delta]
 \simeq\frac{\mathbb C[X,Q,U,V]}
 {\bigl(X^{r+1}V-U^2-c_r^2Q^d\bigr)}.
 }
\]
四つの生成元の重み付き次数は、それぞれ $1,2,d,2m+r+1$ である。
また $x=0$ での値は
\begin{equation}
 \Delta\bmod x
 =2(-1)^jc_r y^d w+\frac{c_r}{r+1}y^{2m}z^{r+1}.
 \label{eq:Delta-specialization-v34}
\end{equation}
\end{proposition}

\begin{proof}
$J_r^0=\sum_{k=0}^r\binom rk q^{r-k}y^{2k+1}/(2k+1)$ と置く。
形式的に $u=xz/y^2$ と書けば
\[
 J_r^0=y^{2r+1}P_r(u),\qquad
 P_r(u)=\sum_{i=0}^r\binom ri c_{r-i}u^i.
\]
係数恒等式
\[
 \binom ri c_{r-i}
 =c_r(-1)^i\binom{r+1/2}{i}\quad(0\le i\le r),\qquad
 c_r(-1)^{r+1}\binom{r+1/2}{r+1}=-\frac1{2(r+1)}
\]
から、$P_r$ は $c_r(1-u)^{r+1/2}$ の $r$ 次までの打切りであり、
\[
 P_r(u)^2-c_r^2(1-u)^{2r+1}
 =\frac{c_r}{r+1}u^{r+1}+O(u^{r+2}).
\]
従って $(J_r^0)^2+c_r^2q^{2r+1}$ は $x^{r+1}$ で割り切れる。
これは多項式の係数恒等式であって、$y$ の可逆性を仮定する結論ではない。
実際、分子は $x,y,z$ の多項式であり、その $x$ 次数 $0,\ldots,r$ の係数が消える。

$J=x^{r+1}w+q^jJ_r^0$ と $d=2j+2r+1$ を代入すれば、
\[
 \Delta=x^{r+1}w^2+2q^jJ_r^0w
       +q^{2j}\frac{(J_r^0)^2+c_r^2q^{2r+1}}{x^{r+1}}
\]
は多項式であり、式\eqref{eq:Delta-specialization-v34}も従う。
$N$ は $x,q,J$ を殺すので $N(\Delta)=0$ である。

表示環を $B$ と書き、$X,Q,U,V$ を $x,q,J,\Delta$ へ送る。
表示式は $V$ に関して原始的な一次式なので $B$ は整域である。
$x$ を反転すると
\[
 A[x^{-1}]=\mathbb C[x,x^{-1},q,J,t],\qquad t=y/x,\qquad N=\partial_t.
\]
従って $B\to A$ は単射で、
$B[x^{-1}]=A[x^{-1}]^N$ である。

さらに $B/xB\to A/xA$ も単射である。
実際、その源は
$\mathbb C[Q,U,V]/(U^2+c_r^2Q^d)$ であり、
\[
 Q\longmapsto-y^2,\qquad U\longmapsto(-1)^jc_ry^d
\]
は、$d$ が奇数であるため既約なcuspのパラメータ表示を与える。
式\eqref{eq:Delta-specialization-v34}の $V$ の像は $w$ に関する非零の一次式で、
$\mathbb C(y)$ 上超越的である。従って追加の関係はない。
よって $B\cap xA=xB$ が成り立つ。
$f\in A^N\subset B[x^{-1}]$ を $f=x^{-e}b$ と書き、$e$ を最小に取る。
$e>0$ なら $b\in B\cap xA=xB$ となって矛盾するので、$f\in B$ である。
\end{proof}

この不変量は、接線系列を特殊な有理pencilとしても記述する。
$\Omega=dx\wedge dy\wedge dz\wedge dw$、
$\omega=\iota_E\iota_{H_{m,r}\wedge N}\Omega$ と置くと
\[
 2q\,\mathrm dJ-(m+r+1)J\,\mathrm dq=2x^r\omega
\]
である。左辺は $E,H'_r,N$ を全て殺すので商体上で $\omega$ の倍数であり、
$dw$ 係数を比較すると比例係数は $2x^r$ と決まる。
従って $r>0$ の場合、共通因子 $x^r$ の除去が重み $m+3$ の一形式を与える。
$r=0$ では通常の有理型 $R(2,m+1)$ に戻る。

\begin{theorem}[Calabi--Yau族の多項式中心]
\label{thm:affine-polynomial-center-v34}
$\alpha=-(m+r+3)/(m+3)$、$s\ne0$ とする。
Poissonテンソル $H_{m,r}\wedge N$ の多項式中心と、
$\mathcal S_{j,r}^{(q)}(\alpha,s)$ の量子中心は、同じ部分環として
\[
 Z_{\mathrm{Pois}}(A)=Z\bigl(\mathcal S_{j,r}^{(q)}(\alpha,s)\bigr)
 =\ker N\cap\ker H_{m,r}
\]
であり、次の三つの場合に完全に分かれる。
\[
 \boxed{
 \begin{array}{c|c}
 r\ge2&\mathbb C\\[2pt]
 r=0\ (m\text{ は偶数})&\mathbb C[q,J]\\[2pt]
 r=1\ (m\text{ は奇数})&
 \mathbb C[S,P,R]/(R^2-SP^2-\tfrac{16}{81}S^{m+3})
 \end{array}
 }
\]
$r=1$ の場合の生成元は
\begin{equation}
 S=q^2\Delta,\qquad
 P=2xJ\Delta^{(m+3)/2},\qquad
 R=q(x^2\Delta+J^2)\Delta^{(m+3)/2}.
 \label{eq:ADE-central-generators-v34}
\end{equation}
それらの次数は $2(m+3),(m+2)(m+3),(m+3)^2$ である。
従ってこの中心のスペクトルは $D_{m+4}$ 型のDu Val曲面である。
\end{theorem}

\begin{proof}
$H_{m,r}(x)\ne0$、$N(x)=0$、$N(y)=x$ なので、二つの導分は一般点で一次独立である。
従ってPoisson中心は $\ker H_{m,r}\cap\ker N$ である。

量子中心について、$X=X_\alpha=-2H_{m,r}$、$T=1+sN$ とする。
任意の $a\in A$ に対して
\[
 x\star a=xa,\qquad a\star x=T^\alpha(a)x.
\]
$\alpha\ne0$ であり、局所冪零性から
$\ker(T^\alpha-1)=\ker N$ である。
具体的には、各有限次元の $N$-巡回空間上で
$T^\alpha-1=N U(N)$、$U(0)=s\alpha\ne0$ となり、$U(N)$ は可逆である。
従って中央元 $a$ は $N(a)=0$ を満たす。
このとき
\[
 [y,a]_\star=sxX(a)
\]
なので $X(a)=0$ である。逆に $N(a)=X(a)=0$ なら、星積の全高次項が消え、
$a$ は中央元である。従って二つの中心は一致する。
また $A^N$ の元同士の星積は通常の可換積そのものである。

前命題の四生成元の $2(m+3)H_{m,r}$ に関する重量は
\[
 x:m+r+3,\quad q:2r,\quad J:rd,\quad
 \Delta:G=(r-1)m+(r-3)(r+1).
\]
$r\ge3$ では全て正であり、$r=2$ でも $m\ge4$ より $G=m-3>0$ である。
従って $r\ge2$ の零重量部分は定数だけとなる。

$r=0$ では $x,\Delta$ の $H$ 重量は $1/2,-1/2$、$q,J$ は零である。
零重量単項式では $x$ と $\Delta$ の指数が等しく、
$x\Delta=J^2+q^{m+1}$ により消去できる。
従って中心は代数的に独立な二元 $q,J$ が生成する多項式環である。

$r=1$ とし、$d=m+2$、$c=c_1=-2/3$ と書く。
不変環と整数重量は
\[
 B=\frac{\mathbb C[x,q,J,\Delta]}{(J^2-x^2\Delta+c^2q^d)},\qquad
 (\mathrm{wt}(x),\mathrm{wt}(q),\mathrm{wt}(J),\mathrm{wt}(\Delta))
 =(d+2,2,d,-4).
\]
$B$ は $\mathbb C[x,q,\Delta]$ 上で基底 $1,J$ を持つ自由加群である。
まず $J$ を含まない零重量単項式を考える。
$x$ の指数は偶数で、$x^{2h}q^b\Delta^e$ について
$h$ と $b$ の偶奇は一致しなければならない。
従ってこれらの単項式は
\[
 S=q^2\Delta,\qquad
 V=x^2q\Delta^{(d+3)/2},\qquad
 W=x^4\Delta^{d+2}
\]
で生成される。$J$ を一つ含む零重量単項式では $x$ の指数が奇数となり、
$xJ\Delta^{(d+1)/2}=P/2$ を取り出せば、残りは前の場合に帰着する。
よって $B_0$ は $S,V,W,P$ で生成される。
関係式から
\[
 P^2=4W-4c^2VS^{(d-1)/2},\qquad V^2=SW
\]
なので $W$ を消去できる。
$R=2V-c^2S^{(d+1)/2}$ と置くと、これは
式\eqref{eq:ADE-central-generators-v34}の $R$ に一致し、
\[
 R^2=SP^2+c^4S^{d+1}
\]
を得る。$S,V$ は多項式環 $\mathbb C[x,q,\Delta]$ 内の
独立な指数ベクトルを持つ単項式なので代数的に独立である。
上の関係から $V$ は $\mathbb C(S,P)$ 上代数的であり、従って $S,P$ も代数的に独立である。
表示多項式は
$\mathbb C(S,P)$ 上で $R$ に関して既約である
（$S(P^2+c^4S^d)$ の $S=0$ での付値は一）。
従ってこの関係以外に関係はなく、表示環が中心全体である。
$c^4=16/81$、$d+1=m+3$ を代入すれば主張の表示を得る。
最後に $P,R$ を $c^2$ で割れば係数が一となり、
$R^2-SP^2-S^{m+3}=0$ は通常の $D_{m+4}$ の方程式である
\cite[Table 1]{ReidDuVal}。
\end{proof}

\subsection{二つの単項式表示の等変同値}

$y$ 型にも前定理を適用するため、座標変換を明示する。
この変換は単なる有理変換ではなく、多項式の次数付き自己同型である。

\begin{proposition}[二つの表示は同じアフィン作用を与える]
\label{prop:affine-q-y-conjugacy-v34}
\[
 c_{j,r}=(-1)^r\frac{(2r)!!(2j-1)!!}{(2j+2r+1)!!},\qquad
 \lambda_{j,r}=(-1)^j\frac{c_{j,r}}{c_r}\ne0
\]
とし、前節の二つの $J$ の $w$ を含まない項を $J_0^{(y)},J_0^{(q)}$ と書く。
すると
\[
 b_{j,r}=\frac{J_0^{(y)}-\lambda_{j,r}J_0^{(q)}}{x^{r+1}}
 \in\mathbb C[x,y,z]_m
\]
である。$y$ 型の座標で
\[
 w'=\lambda_{j,r}^{-1}(w+b_{j,r})
\]
と置けば $N(w')=-q^jz^r$ であり、全ての $\alpha$ について
$X_\alpha(w')=(m\alpha+d)w'$ となる。
従って両者のPoisson代数およびCGG量子化は次数付き同型である。
\end{proposition}

\begin{proof}
$u=xz/y^2$ と形式的に置くと
\[
 J_0^{(y)}=y^d\sum_{i=0}^r\binom ri c_{j,r-i}u^i.
\]
二重階乗の係数恒等式
\[
 \binom ri c_{j,r-i}
 =c_{j,r}(-1)^i\binom{j+r+1/2}{i}\quad(0\le i\le r)
\]
より、この和は $c_{j,r}(1-u)^{j+r+1/2}$ の $r$ 次打切りである。
一方、$J_0^{(q)}=q^jJ_r^0$ の同じ表示は
$y^d(u-1)^jP_r(u)$ なので、$\lambda_{j,r}$ を掛ければ
$u$ 次数 $0,\ldots,r$ の係数が一致する。
従って分子の $x$ 次数 $0,\ldots,r$ の係数は消え、$b_{j,r}$ は多項式である。
次数は $d-(r+1)=m$ となる。

$N_0(J_0^{(y)})=x^{r+1}y^{2j}z^r$、
$N_0(J_0^{(q)})=x^{r+1}q^jz^r$ なので
\[
 N_0(b_{j,r})=y^{2j}z^r-\lambda_{j,r}q^jz^r.
\]
これが $N(w')$ の式を与える。
また二つの $J_0$ の $X_\alpha$ 固有値はともに $d(\alpha+1)$ なので、
$b_{j,r}$ の固有値は $d(\alpha+1)-(r+1)\alpha=m\alpha+d$ である。
従って $X_\alpha$ も保存され、普遍変形式から星積も保存される。
\end{proof}

\begin{example}[重み三の二つのアフィン族]
$(j,r)=(0,3)$ の中心は定数環である。
これに対し $(j,r)=(1,1)$ では
\[
 \lambda_{1,1}=-\frac15,\qquad b_{1,1}=\frac15yz^2,\qquad
 w'=-5w-yz^2
\]
によって $N(w)=-y^2z$ を $N(w')=-qz$ に変換できる。
その通常の $4$-Calabi--Yau量子化の中心は
\[
 \mathbb C[S,P,R]/(R^2-SP^2-\tfrac{16}{81}S^6),
\]
すなわち $D_7$ 型である。
従って両者は、たとえ関数斜体の中心だけでは区別できなくても、
斉次座標代数の中心によって区別できる。
これはそれらの $\operatorname{qgr}$ 圏が同値でないことまでを主張するものではない。
\end{example}

\subsection{全共鳴族に共通する量子関数斜体}
\label{subsec:affine-quantum-field-v34}

多項式中心の三つの場合は、次数零の関数斜体を取ると区別できなくなる。
$Q_{\mathrm{gr}}(S)$ は非零斉次元を可逆にした次数付き商環を表し、
$Q_{\mathrm{gr}}(S)_0$ をその次数零の除法環とする。
以下では $m=2j+r\ge3$、$d=m+r+1$、$s\ne0$ とする。
前節の等変同値により $q$ 型を用いてよい。

\begin{proposition}[共通のWeyl斜体]
\label{prop:affine-quantum-field-v34}
全ての $\alpha\in\mathbb C$ に対して、
$S=\mathcal S_{j,r}^{(q)}(\alpha,s)$ の次数零の関数斜体は
\[
 \boxed{
 Q_{\mathrm{gr}}(S)_0
 \simeq\operatorname{Frac}\bigl(\mathbb C(C)[\eta][\partial;d/d\eta]\bigr),
 \qquad C=J^2/q^d,
 \quad[\partial,\eta]=1.
 }
\]
その中心は正確に $\mathbb C(C)$ であり、PIではない。
対応する重み付き射影Poisson関数体の中心も $\mathbb C(C)$ である。
\end{proposition}

\begin{proof}
$N$ 不変な元 $x,q,J$ は、星積に関しても互いに可換である。
$F=\mathbb C(x,q,J)$、$t=y/x$ と置くと、
有理座標 $x,q,J,t$ に関して $N=\partial_t$ となる。
$F[t]$ 上では星積の全ての和が有限であり、
\[
 [t,f]_\star=sX_\alpha(f)\quad(f\in F)
\]
となる。星積冪と通常の $t^n$ の三角変換により、これは微分Ore拡大
$F[t;sX_\alpha|_F]$ である。

この表示は元の代数と同じ除法閉包を持つ。
実際、正規元 $x$ の局所化の中で $t=x^{-1}\star y$ が成り立つ。
逆に $y=xt$、$z=(q+x^2t^2)/x$ と $J$ の式から、
元の四生成元が $F[t]$ の中に復元される。
従って有理関数一般に無限の星積公式を用いる必要はない。

次数零の元を
\[
 u=q/x^2,\qquad v=J/x^d,\qquad t=x^{-1}\star y
\]
と取ると、$X_\alpha(u)=2u$、$X_\alpha(v)=dv$ であり、
\[
 [t,u]=2su,\qquad[t,v]=dsv,\qquad[u,v]=0.
\]
さらに $[t,x]=s\alpha x$ なので、
\[
 D_0=\operatorname{Frac}\bigl(\mathbb C(u,v)[t;s(2u\partial_u+dv\partial_v)]\bigr)
\]
に対して
\[
 Q_{\mathrm{gr}}(S)=D_0[x,x^{-1};\sigma],\qquad
 \sigma(t)=t-s\alpha,\quad\sigma(u)=u,\quad\sigma(v)=v
\]
となる。元の生成元は次数一または次数 $m$ の単項式と $D_0$ の元の積として復元でき、
逆方向には $u,v,t$ は元の斉次分数である。よって次数零部分は正確に $D_0$ である。

$d=2\ell+1$、$\ell=j+r$ と書き、
\[
 C=v^2/u^d,\qquad\eta=v/u^\ell
\]
と置く。$u=\eta^2/C$、$v=\eta^d/C^\ell$ なので
$\mathbb C(u,v)=\mathbb C(C,\eta)$ である。
$[t,C]=0$、$[t,\eta]=s\eta$ より、
$\partial=\eta^{-1}t/s$ とすればWeylの関係を得る。

中心の計算には微分Ore除法環の形式的な $\partial^{-1}$ 級数表示を用いられる。
$\eta$ と可換する級数の最高次数は
$[\partial^n,\eta]=n\partial^{n-1}$ により零である。
次数零の係数を引いて負次数にも同じ議論を適用すると、
$\eta$ の中心化体は $\mathbb C(C)(\eta)$ である。
さらに $\partial$ と可換する条件は $\eta$ 微分が零という条件なので、中心は
$\mathbb C(C)$ である。中心上無限次元の除法環はPIではない
\cite[Part~II, Theorem~50]{Kaplansky1972}。

Poisson側でも同じ有理座標を用い、$N=\partial_t$ と
$X_\alpha|_{\mathbb C(u,v)}=2u\partial_u+dv\partial_v$ の共通不変体を取ればよい。
$(C,\eta)$ 表示で後者は $\eta\partial_\eta$ となるので、中心は同じ $\mathbb C(C)$ である。
\end{proof}

\begin{corollary}[関数斜体と射影圏の区別]
全ての $m,j,r,\alpha$ と非零の $s$ に対し、上の次数零の関数斜体は互いに同型である。
一方、異なる $m$ に対する $D^b(\operatorname{qgr}S)$ は導来同値でない。
\end{corollary}
\begin{proof}
前半は前命題による。標準full exceptional collectionの長さが $m+3$ なので、
$K_0(D^b(\operatorname{qgr}S))\simeq\mathbb Z^{m+3}$ である。
この階数は導来同値で保存される。
\end{proof}

特に、Calabi--Yau値における $r\ge2$ の定数中心、$r=0$ の平面中心、$r=1$ のDu Val中心という
斉次座標環の差は、次数零の関数斜体の同型型には残らない。
同じ重みの異なる $j$ の射影圏が同値かどうかは、ここでは決定していない。

\section{重み三への特殊化と幾何的辞書}
\label{sec:poisson-coordinates-v32}
以下の成分分類では $m=3$ とし、$X=X_3$、$\mathscr P=\mathscr P_3$ と書く。
命題~\ref{prop:general-equations-v33}は、五十三変数の標準形
\[
 c\in S_3,\quad D\in\mathfrak{sl}_3,\quad [F]\in S_4^{\oplus3}/S_3\mathbf x
\]
と三十八個の独立な二次式を与える。
\begin{proposition}[重み三の方程式]
\label{prop:weighted-poisson-equations-v31}
全Jacobi式は
\begin{align}
 D(c)&=0,\label{eq:weighted-base-jacobi-v31}\\
 2F(c)-(\mathbf x\times D)\cdot F-c\operatorname{div}F&=0.
 \label{eq:weighted-mixed-jacobi-v31}
\end{align}
係数ベクトルは $p=(\nabla c-\mathbf x\times D)/3$ である。
\end{proposition}
証明は一般重みの命題で与えた。
$S_2^{\oplus3}\simeq S_3\oplus\mathfrak{sl}_3$ という三変数二次Poisson構造の分解は
\cite{LiuXu1992,DoninMakarLimanov1998}と関係する。

\section{三つの基礎的な成分}
\label{sec:basic-components-v32}
\subsection{滑らかな三次曲線に対応する19次元成分}

\begin{theorem}[滑らかな三次曲線上の標準形]
\label{thm:smooth-cubic-component-v31}
$c=0\subset\mathbb P^2$ が滑らかであるとする。
式 \eqref{eq:weighted-base-jacobi-v31}--\eqref{eq:weighted-mixed-jacobi-v31}
の解はちょうど
\[
 D=0,\qquad [F]=[\nabla c\times\nabla g],
 \quad g\in S_3/\mathbb Cc
\]
である。したがってこの開部分は、滑らかな三次形式の $10$ 次元開集合上の
階数 $9$ のベクトル束であり、滑らかな既約 $19$ 次元スキームである。
その閉包は Poisson 構造のスキームの既約成分をなす。

さらに、Euler 同値による代表の変更の後、重み付き座標変換
$w'=w+3g$ により Poisson 構造は
\[
 \Pi_c=\frac13\left(
 c_1\partial_2\wedge\partial_3+
 c_2\partial_3\wedge\partial_1+
 c_3\partial_1\wedge\partial_2\right),
 \qquad\{w',x_i\}=0
\]
という形になる。
\end{theorem}

\begin{proof}
滑らかさから $c_1,c_2,c_3$ は多項式環の二次正則列をなす。
その Koszul 複体は一次係数の syzygy を持たないので、
$D(c)=\sum_iD_i c_i=0$ から $D=0$ となる。

残る式は $2F(c)=c\operatorname{div}F$ である。
\[
 F_0=F-\frac{\operatorname{div}F}{6}x
\]
と置くと $F_0(c)=0$ かつ $\operatorname{div}F_0=0$ である。
Koszul 完全性から $F_0=\nabla c\times A$、$A\in S_2^{\oplus3}$ と書ける。
その発散は
\[
 \operatorname{div}(\nabla c\times A)
 =-\nabla c\cdot\operatorname{curl}A
\]
である。$\operatorname{curl}A$ は線形なので、同じ一次 syzygy の消滅から
$\operatorname{curl}A=0$ となる。
斉次 Poincar\'e の補題により
$A=\nabla g$、$g=(x\cdot A)/3\in S_3$ と書ける。

もし $\nabla c\times\nabla g=hx$ なら、$\nabla c$ との内積を取り
$3hc=0$ を得るので $h=0$ である。
再び Koszul 完全性により $\nabla g=\lambda\nabla c$、
したがって $g=\lambda c$ である。これで核が $\mathbb Cc$ と分かる。
各線形写像の階数は滑らかな三次形式の開集合上で一定なので、
この記述はスキームとしての階数 $9$ のベクトル束の同型を与える。
この開部分が滑らかで既約であることから、その閉包は一つの既約成分となる。

最後に、採用した括弧の符号では
\[
 \{x_i,g\}_{\Pi_c}
 =-\frac13(\nabla c\times\nabla g)_i
\]
なので、$w'=w+3g$ は混合項を消す。
\end{proof}

この定理は重み付き座標変換を許した Poisson 構造の標準化である。
特殊ファイバーの標識を保つ形式変形の同値を主張しているわけではない。
また $c$ が特異な場合には結論が変わる。
例えば $c=3x_1x_2x_3$、$D=0$ では、$F$ の解空間の Euler 商は
$11$ 次元であるのに対し、三次式による $w$ の平行移動の像は $9$ 次元である。
この境界には中央変数への標準化で消えない混合方向が残る。


% Additional self-contained subsection for Paper VIII.
\subsection{別の 19 次元成分：一般の対数型構造}
\label{subsec:logcanonical-component-v31}

滑らかな三次曲線から来る成分が Poisson 構造の空間を尽くすわけではない。
本節では別の成分を明示する。以下は全成分の分類ではない。

線形場 $D\in\mathfrak{sl}_3$ が正則半単純で、
\[
 \dim\ker(D:S_3\longrightarrow S_3)=1
\]
であると仮定する。固有座標を選ぶと
\[
 D=\delta_1x_1\partial_1+\delta_2x_2\partial_2+
     \delta_3x_3\partial_3,
 \quad \delta_1+\delta_2+\delta_3=0,
 \quad \ker(D|_{S_3})=\mathbb Cx_1x_2x_3.
\]
$c\neq0$ とすると、ある $t\neq0$ により $c=tx_1x_2x_3$ と書ける。
対応する $p$ は
\[
 p=(\alpha x_2x_3,\beta x_3x_1,\gamma x_1x_2),
 \qquad t=\alpha+\beta+\gamma
\]
という形であり、
\[
 (\delta_1,\delta_2,\delta_3)
 =(\gamma-\beta,\alpha-\gamma,\beta-\alpha)
\]
となる。

さらに、有限個の開条件
\begin{align*}
 &3(\alpha,\beta,\gamma)-tm\neq0
 &&(m\in\mathbb Z_{\geq0}^3,\ |m|=3),\\
 &3\alpha+t\neq0,\qquad3\beta+t\neq0,\qquad3\gamma+t\neq0
\end{align*}
を課す。これらは固有座標の並べ替えに関して不変である。
条件は空でない。例えば $(\alpha,\beta,\gamma)=(2,3,5)$ のとき
$D$ の固有値は $(2,-3,1)$ であり、三次単項式の固有値は
$m=(1,1,1)$ の場合に限って零になる。

\begin{theorem}[対数型の開成分]
\label{thm:logcanonical-component-v31}
上の開条件の下で、混合 Poisson 方程式の解の Euler 商は $10$ 次元である。
これらの解はすべて三次式による重み付き座標変換 $w\mapsto w+g$ で消去できる。
したがって、この開部分は滑らかで既約な $19$ 次元スキームであり、
その閉包は Poisson 構造のスキームの一つの既約成分となる。
この成分は定理~\ref{thm:smooth-cubic-component-v31} の成分とは異なる。
\end{theorem}

\begin{proof}
Laurent 単項式を用い、四次ベクトル場の各文字成分を
\[
 F_i=x_ix^m v_i,\qquad |m|=3
\]
と書く。$m\geq0$ の場合は $v\in\mathbb C^3$ が任意である。
それ以外の場合は、ちょうど一つの成分が $m_i=-1$ であり、
残り二成分は非負で和が $4$、かつ $v_j=0$ ($j\neq i$) でなければならない。
混合方程式の左辺は、この文字成分上で
\[
 x^{m+(1,1,1)}
 \bigl(3(\alpha,\beta,\gamma)-tm\bigr)\cdot v
\]
となる。異なる $m$ は異なる六次単項式を与えるため、方程式は文字ごとに分離する。

$m\geq0$ の $10$ 個の文字について、Euler 商は
$\mathbb C^3/\mathbb C(1,1,1)$ である。
係数ベクトルの成分の和が零なので方程式は商上で定まり、
開条件によりその核は各文字で一次元である。
$m_i=-1$ の文字は各 $i$ について $5$ 個ある。
その係数は $3\alpha_i+t$ なので、これら $15$ 個の成分はすべて零になる。
したがって全体の核は $10$ 次元であり、混合写像の階数は $25$ である。

$g=x^m$ による $w$ の平行移動の変化は
\[
 \{x,g\}_\Pi
 =x^m\bigl(x_i v_i\bigr)_{i=1}^3,
 \qquad v=-(\alpha,\beta,\gamma)\times m
\]
である。この $v$ が Euler 方向 $\lambda(1,1,1)$ に入ると仮定する。
$(\alpha,\beta,\gamma)$ との内積を取ると $\lambda t=0$、したがって $\lambda=0$。
すると $(\alpha,\beta,\gamma)$ と $m$ が比例し、
$|m|=3$ から $3(\alpha,\beta,\gamma)=tm$ となって開条件に反する。
よって各文字の一次元の核は平行移動の像と一致する。

正則半単純かつ上の核次元が $1$ である $D$ は
$\mathfrak{sl}_3$ の非空な開集合をなす。
そこでは $\ker(D|_{S_3})$ は線束であるから、
$D(c)=0$、$c\neq0$ を満たす $(D,c)$ の空間は滑らかで既約な $9$ 次元空間となる。
追加の有限個の非消滅条件はその非空な開集合を定める。
混合写像はそこでは一定階数 $25$ を持つので、
その Euler 商の核は階数 $10$ のベクトル束である。
これにより次元 $19$ とスキームとしての滑らかさ・既約性が従う。
滑らかな開部分の閉包は一つの既約成分となる。

滑らかな三次曲線の成分の閉包は閉条件 $D=0$ の中に含まれる。
本節の開部分では $D\neq0$ なので、両成分は異なる。
\end{proof}

\begin{remark}[三角形の境界で残る二方向]
$p=(x_2x_3,x_3x_1,x_1x_2)$ の場合、$c=3x_1x_2x_3$、$D=0$ である。
上の文字計算では $m=(1,1,1)$ の方程式が恒等的に零になり、
この文字の核は Euler 商で二次元となる。
一方、$x_1x_2x_3$ は Casimir なので、この文字では三次式による平行移動の像は零である。
残る二方向は、例えば
\[
 F^{(1)}=(x_1^2x_2x_3,0,0),\qquad
 F^{(2)}=(0,x_1x_2^2x_3,0)
\]
の Euler 商で表される。
全体の混合核は $11$ 次元、平行移動の像は $9$ 次元である。
したがって、一般点における中央変数への標準化はこの境界では失敗する。
\end{remark}

% The larger component: exact proof, no computer rank assumption.
\subsection{35 次元の線形成分と非等次元性}
\label{subsec:vertical-component-v31}

$\Pi=0$、すなわち $c=D=0$ とすると、二つの Poisson 方程式はともに恒等的に零になる。
したがって
\[
 \mathcal L=\{\Pi=0\}
 \simeq S_4^{\oplus3}/xS_3\simeq\mathbb A^{35}
\]
は Poisson 構造のスキームに含まれる。
この包含は、単なる高次元の部分族にとどまらない。

\begin{theorem}[垂直な線形成分]
\label{thm:vertical-component-v31}
$\mathcal L$ は Poisson 構造のスキームの既約成分であり、
その一般点では同スキームは滑らかである。
$\mathcal L$ 上の全 $35$ 次元族は、可換多項式環の次数を保つ Ore 拡大により同時に量子化できる。
\end{theorem}

\begin{proof}
$F^*=(x_1^4,x_2^4,x_3^4)$ を取る。
この点で混合方程式を $p\in S_2^{\oplus3}$ の方向に線形化した写像は
\[
 L_{F^*}(p)
 =F^*(x\cdot p)+3p\cdot F^*-(x\cdot p)\operatorname{div}F^*
 \in S_6
\]
である。この写像が単射であることを係数比較で示す。
$\{i,j,k\}=\{1,2,3\}$ として
\[
 p_i=a_i x_i^2+
      \sum_{j\neq i}b_{ij}x_ix_j+
      \sum_{j\neq i}d_{ij}x_j^2+
      e_i x_jx_k
\]
と書く。
$x_i^6$ の係数は $2a_i$ なので、まず $a_i=0$ である。
$x_i^2x_j^4$ の係数は $-3b_{ij}$ なので $b_{ij}=0$ となる。
$x_ix_j^5$ の係数は $-2d_{ij}+b_{ji}$ なので $d_{ij}=0$ となる。
最後に $x_i^4x_jx_k$ の係数は $-3(e_j+e_k)$ であるから、
$e_1=e_2=e_3=0$ である。
よって $L_{F^*}$ は階数 $18$ を持つ。
この $18$ 個の係数を行とし、上の順に係数を列とした小行列の行列式は
\[
 2^3(-3)^6(-2)^6\det
 \begin{pmatrix}0&-3&-3\\-3&0&-3\\-3&-3&0\end{pmatrix}
 =-20155392\neq0
\]
である。

混合方程式は $p$ に関して線形で、その係数は $F$ に関して線形である。
したがって上の小行列式が非零である $F$ の開近傍では、
混合方程式そのものがスキームとして $p=0$ を強制する。
その開集合上で Poisson 構造のスキームは $\mathcal L$ と一致する。
よって $\mathcal L$ は一般点で滑らかな既約成分である。

量子化には Euler 商の線形切断を一つ選ぶ。
例えば $F_1$ に含まれる $x_1$ で割り切れる単項式をすべて消す切断を取れる。
$F=\sum_iF_i\partial_i$ は可換多項式環 $\mathbb C[x_1,x_2,x_3]$ の導分であるから、
\[
 S_\hbar=\mathbb C[x_1,x_2,x_3][[\hbar]]
 [w;\,-\hbar F],\qquad
 wa=aw-\hbar F(a)
\]
は Ore 拡大である。特に $[x_i,w]=\hbar F_i$ であり、
順序付き単項式 $x_1^{a_1}x_2^{a_2}x_3^{a_3}w^b$ が形式基底をなす。
各 $F_i$ は四次なので関係は重み $(1,1,1,3)$ を保つ。
普遍 $F_{\mathrm{univ}}$ の $35$ 個の係数を完備局所基底の極大イデアルに取り、
式の $\hbar F$ を $F_{\mathrm{univ}}$ に置き換えれば、
$\mathbb C[[u_1,\ldots,u_{35}]]$ 上の平坦な形式変形が得られる。
形式変形の HKR 規約に応じた一様なパラメータの定数倍は、この結論を変えない。
\end{proof}


\section{可積分一形式による成分の幾何的表示}
\label{sec:geometric-dictionary-v32}

上の座標に対応する一形式は
\[
 \alpha=c\,dw-w\,dc+w(x\times D)\cdot dx+(F\times x)\cdot dx.
\]
重み六の水平一形式の空間と $H^0(X,\wedge^2T_X)$ の対応は線形同型である。
これは滑らかな三次元スタック上の
$\wedge^2T_X\simeq\Omega_X^1\otimes K_X^{-1}$ による。
局所体積形式で確認すれば、Jacobi方程式とFrobenius方程式は同じ係数イデアルを与える。
したがって、以下の幾何的表示はスキームの方程式と両立する。

\subsection{最初の三成分}

滑らかな三次曲線型では、三次式による $w$ の移動後に
$\alpha=c\,dw-w\,dc$ となる。
一般の重み三の二つの斉次式は、線形結合によって一つを $c(x,y,z)$、他方を $w+g_3(x,y,z)$
と選べる。従ってその成分は $R(3,3)$ である。
重み三の空間は十一次元なので、二次元部分空間のGrassmann多様体と非零倍率から
$\dim\operatorname{Gr}(2,11)+1=19$ という次元も読める。

一般対数型 $p=(\alpha_1yz,\alpha_2zx,\alpha_3xy)$、$t=\sum\alpha_i$ では、
混合項を消すと
\[
 \alpha=xyzw\left(t\frac{dw}{w}-3\alpha_1\frac{dx}{x}
                   -3\alpha_2\frac{dy}{y}-3\alpha_3\frac{dz}{z}\right).
\]
従ってこの成分は $L(1,1,1,3)$ である。
一般の一次三因子は座標に取れ、重み三の因子の $w$ 係数が非零である開集合で
それを $w+g_3$ に取れるから、逆の包含も一般点で成立する。

垂直型 $p=0$ では、$\alpha=(F\times x)\cdot dx$ は $w$ に依存しない。
これは射影 $X\dashrightarrow\mathbb P^2$ による
$H^0(\mathbb P^2,\Omega^1(6))$ の引戻し全体である。
平面上では一形式のFrobenius条件は自動的に成立し、この空間の次元は
$3h^0(\mathcal O(5))-h^0(\mathcal O(6))=63-28=35$ である。
成分性は定理\ref{thm:vertical-component-v31}の横方向の単射性により証明済みである。

\subsection{次数と同値関係}

これらの表示には重み付き座標変換を使うが、成分の次元をその群で割ってはいない。
有理型・対数型の一般点をスカラー倍で割れば、通常の葉層のパラメータ空間になる。
本稿では零Poisson構造を含むアフィン錐とその冪零元を保つ。
また、特殊ファイバーの標識を固定した結合代数の形式変形の同値を、
この座標変換による標準形と取り違えない。

% Insert after the weighted Poisson equations; all statements use their normalization.
\section{直線と二次曲線の層、二つの有理型成分とその交叉}
\label{sec:lineconic-v32}

この節では $S=\mathbb C[x,y,z]$ とする。
三次形式が直線と非退化二次形式の積であり、直線が二次曲線に横断的に交わる層を調べる。
$D(c)=0$ の形式近傍を先に決定することで、混合項の階数の跳びを
Poisson スキーム全体の既約成分へ結び付ける。

\subsection{正則半単純な基底の形式切片}

\[
 c_0=x(x^2+yz),\qquad D_0=y\partial_y-z\partial_z,
 \qquad c=c_0,\quad D=aD_0
\]
と置く。この正規化で $a$ は $c$ と $D$ の相対的な係数である。
座標変換で $c$ の determinant twist を無視して $a$ を定義しているわけではない。
以下では $\mathrm{SL}_3$ による座標変換と、$c,D$ の共通の非零係数を明示的に分ける。

\begin{lemma}[直線・二次曲線型の九次元基底]
\label{lem:lineconic-base-v32}
$a_0\ne0$ とし、$(c,D)=(c_0,a_0D_0)$ を考える。
スキーム $\{D(c)=0\}\subset S_3\oplus\mathfrak{sl}_3$ はこの点の近傍で滑らかで、次元は九である。
形式座標を選ぶと、八つの滑らかな変数と一変数 $a$ に分かれ、
混合写像は単元と可逆な基底変換を除いて
\[
 M_a([F])=2F(c_0)-(x\times aD_0)\cdot F-c_0\operatorname{div}F
\]
に一致する。
\end{lemma}

\begin{proof}
$D$ は三つの相異なる固有値を持つ。固有値を順序付け、局所的に固有線を選ぶと、
六つの共役軌道方向と固有値 $(u,v,-u-v)$ を座標に取れる。
これは形式的、あるいは適当な \'etale 近傍で行う。
$v=a_0,u=0$ の点では、三次単項式のうち固有値が零になるものは $x^3,xyz$ の二つだけである。
他の八つの係数は $D(c)=0$ の式で単元に掛けられるので消去できる。
残る式は
\[
 c=A x^3+Bxyz,\qquad 3uA=0.
\]
ここでは $A,B$ はともに単元だから $u=0$ であり、
六つの共役軌道方向、$v,A,B$ が九つの自由な座標となる。

さらに対角的な $\mathrm{SL}_3$ 座標変換と形式的な三乗根を使い $A=B$ とできる。
$c=Bc_0$、$D=vD_0$ と書けば $a=v/B$ であり、混合写像は $BM_a$ となる。
この変換に使った $A/B$、共通係数 $B$ と六つの共役軌道方向が、上記の八つの変数を与える。
\end{proof}

\subsection{混合写像の全階数跳躍}

$F$ の Euler 商の代表として、$F_x$ のうち $x$ で割り切れる単項式を零とする。
これにより $35$ 個の単項式基底を得る。
$(x,y,z)\mapsto(x,ty,t^{-1}z)$ の文字で混合写像を分解する。
$F_i=x^{n_1}y^{n_2}z^{n_3}$ の文字は $n_2-n_3-\operatorname{wt}(x_i)$ であり、
出力の六次単項式の文字は $n_2-n_3$ である。

\begin{lemma}[直線・二次曲線型の混合行列]
\label{lem:lineconic-matrix-v32}
$M_a:S_4^{\oplus3}/xS_3\to S_6$ の一般階数は $25$ であり、階数が下がる値は
\[
 a=2,\qquad a=-1
\]
だけである。各点における階数はそれぞれ $15,21$ である。
$a=2$ における局所的な Smith 標準形には十五個の単元と十個の $a-2$ があり、
$a=-1$ では二十一個の単元と四個の $a+1$ がある。
いずれも恒常的な核は十次元である。
\end{lemma}

\begin{proof}
各文字ブロックの最大非零小行列式の最大公約数を、定数倍を除いて表に示す。
\[
\begin{array}{c|c|c|c}
\text{文字}&\text{行数}\times\text{列数}&\text{一般階数}&\text{最大小行列式の gcd}\\\hline
\pm5&1\times1&1&a-2\\
\pm4&2\times2&2&(a-2)(a+1)\\
\pm3&2\times3&2&a-2\\
\pm2&3\times4&3&(a-2)(a+1)\\
\pm1&3\times5&3&a-2\\
0&4\times5&3&1
\end{array}
\]
零文字の四次小行列式は全て零である。
この表は、上記の単項式基底で $M_a$ を展開して得られる高々四行五列の行列から直接求まる。
付属の \texttt{verify\_lineconic\_components\_v32.py} は
$\mathbb Q[a]$ 上で各ブロックの全最大小行列式を展開し、その gcd を計算する。
したがってこれは有限個の $a$ を試す計算ではなく、全パラメータを扱う多項式の恒等式である。

各非零文字ブロックは $a=2$ で階数を一つ落とし、
$a=-1$ では文字 $\pm4,\pm2$ の四ブロックだけが階数を一つ落とす。
最大小行列式の gcd の各零点の位数も一なので、局所的な非単元の基本因子は全て一次である。
一般階数の和は $25$ であり、主張の標準形が従う。
\end{proof}

\begin{theorem}[二つの共鳴点における完全な形式局所環]
\label{thm:lineconic-local-v32}
Poisson スキームの点 $(c,D,F)=(c_0,a_0D_0,0)$ を考える。
$a_0=2$ のとき、その完備局所環は
\[
 \frac{\mathbb C[[b_1,\ldots,b_8,z_1,\ldots,z_{10},t,v_1,\ldots,v_{10}]]}
 {(tv_1,\ldots,tv_{10})}
\]
に同型である。
$a_0=-1$ のときは、同じ式で $v_1,\ldots,v_{10}$ を $v_1,\ldots,v_4$ に置き換えた環となる。
従って前者は次元 $28$ と $19$、後者は次元 $22$ と $19$ の二つの滑らかな枝を持ち、
いずれもその交叉は滑らかな十八次元である。両局所環は被約である。
\end{theorem}

\begin{proof}
補題 \ref{lem:lineconic-base-v32} によって、基底側の九変数以外の変数を消去する。
$t=a-a_0$ とし、補題 \ref{lem:lineconic-matrix-v32} の Smith 標準形を
$\mathbb C[[t]]$ 上で用いる。単元の基本因子に対応する $F$ の変数を消去すると、
恒常的な核の十変数 $z_i$ と、基本因子 $t$ に対応する $v_i$ が残る。
混合方程式は $F$ に線形なので、これ以外の項や方程式は生じない。

形式冪級数環において
\[
 (tv_1,\ldots,tv_r)=(t)\cap(v_1,\ldots,v_r)
\]
であり、右辺は二つの素イデアルの交叉である。
これが二枝、被約性、交叉および各次元を与える。
\end{proof}

この定理は $F=0$ の点における局所環の記述である。
$v_i$ のいずれかが非零である点では $t=0$ が強制され、
次元 $28$ または $22$ の枝が Poisson スキーム全体の開近傍をなす。
これは新しい既約成分を認定するための局所的な根拠となる。

\subsection{有理型と対数型の幾何}

$\Omega=dx\wedge dy\wedge dz\wedge dw$ とし、
$P=\Pi+F\wedge\partial_w$ に対応する重み六の水平一形式を
\[
 \alpha=\iota_E\iota_P\Omega
 =c\,dw-w\,dc+w(x\times D)\cdot dx+(F\times x)\cdot dx
\]
とする。三次元では $P$ の Poisson 条件と $\alpha\wedge d\alpha=0$ は同値である。
有理型・対数型の一形式そのものは古典的な構成であり、通常の射影空間の場合は
\cite{Molinuevo2016} とその参考文献を参照されたい。
ここでは重み $(1,1,1,3)$ と全 Poisson スキームの形式切片に対して成分性を証明する。

$c=\ell q$、$\deg\ell=1$、$\deg q=2$ と書き、$q$ は非退化、$\ell=0$ は $q=0$ に横断的とする。
正規化 $\ell=x,q=x^2+yz$ の下で
\[
 D=-a\,\nabla\ell\times\nabla q=aD_0
\]
である。ベクトル三重積を用いると混合方程式は
\begin{equation}
 (2+2a)qF(\ell)+(2-a)\ell F(q)-\ell q\operatorname{div}F=0
 \label{eq:lineconic-factorized-v32}
\end{equation}
となる。

\begin{proposition}[二つの有理型成分の標準形]
\label{prop:lineconic-rational-v32}
上の直線・二次曲線の層において次が成り立つ。
\begin{enumerate}
\item $a=2$ では、Euler 商の全ての解は
\[
 [F]=[\nabla\ell\times\nabla H],\qquad H\in S_5/\mathbb C\ell^5
\]
である。対応する一形式は
\[
 \alpha=\ell\,d(qw+H)-5(qw+H)d\ell
\]
となる。この族の閉包は二十八次元の既約成分 $R(1,5)$ である。
\item $a=-1$ では、Euler 商の全ての解は
\[
 [F]=[\nabla q\times\nabla H],\qquad H\in S_4/\mathbb Cq^2
\]
である。対応する一形式は
\[
 \alpha=q\,d(\ell w+2H)-2(\ell w+2H)dq
\]
となる。この族の閉包は二十二次元の既約成分 $R(2,4)$ である。
\item 一般の $a$ では、三次式による $w$ の移動の後に
\[
 \alpha=\ell qw\left\{\frac{dw}{w}-(1+2a)\frac{d\ell}{\ell}
                          +(a-1)\frac{dq}{q}\right\}
\]
と書ける。この族の閉包は十九次元の既約成分 $L(1,2,3)$ である。
\end{enumerate}
三つの成分は相異なる。
\end{proposition}

\begin{proof}
$a=2$ のとき式 \eqref{eq:lineconic-factorized-v32} は
$6F(\ell)=\ell\operatorname{div}F$ となる。
$\ell=x$ とすれば $F_x$ は $x$ で割り切れるので、Euler 同値によって $F_x=0$ にできる。
すると $\partial_yF_y+\partial_zF_z=0$ である。
$\mathbb C[x]$ 上の二変数の斉次 Poincar\'e 補題により、
$F_y=-\partial_zH,F_z=\partial_yH$、$H\in S_5$ と書ける。
核は $\mathbb Cx^5$ であるから、混合解の次元は $21-1=20$ となる。

$a=-1$ のときは $3F(q)=q\operatorname{div}F$ となる。
発散零の代表 $F_0=F-x\operatorname{div}F/6$ を取ると、$F_0(q)=0$ である。
$q$ の偏微分は独立な一次形式だから Koszul 完全性により
$F_0=\nabla q\times A$、$A\in S_3^{\oplus3}$ と書ける。
発散零性から $\nabla q\cdot\operatorname{curl}A=0$ なので
$\operatorname{curl}A=\nabla q\times B$、$B\in S_1^{\oplus3}$ と書ける。
もう一度発散を取ると $\nabla q\cdot\operatorname{curl}B=0$ となるが、
$\operatorname{curl}B$ は定数であるため零である。
従って $B=\nabla b$、$b\in S_2$ と書け、
$\operatorname{curl}(A-q\nabla b)=0$ から $A-q\nabla b=\nabla H_0$ を得る。
結局 $F_0=\nabla q\times\nabla(H_0+qb)$ である。
核が $\mathbb Cq^2$ であることも、Koszul 完全性と Euler の式から従う。
したがって混合解は $15-1=14$ 次元である。

両場合の一形式は
\[
 (\nabla f\times\nabla H)\times x
 =(\deg f)f\nabla H-(\deg H)H\nabla f
\]
を代入して得る。
一般の重み五の多項式は $qw+H$ と書け、一般の重み四の多項式は $\ell w+H$ と書ける。
従ってこれらは、それぞれ次数 $(1,5)$、$(2,4)$ の一般の有理型一形式を与える。

一般の $a$ について、$w$ の三次式による移動は混合核の中に十次元の像を持つ。
例えば $a=1/4$ で、$p$ の三係数は共通因子を持たない。
三次式 $g$ の移動が Euler 方向なら、$p\times\nabla g=hx$ と $p$ の内積を取り
$hc=0$、従って $h=0$ を得る。三係数の最大公約因子が $1$ であることから
$\nabla g=\kappa p$、$\kappa\in\mathbb C$ となり、curl を取ると
$\kappa D=0$ だから $\kappa=0$ である。したがって移動の核は零であり、その階数は十である。
この階数条件は開条件である。
補題 \ref{lem:lineconic-matrix-v32} によって混合核も十次元だから、
一般点の全混合項をこの移動で消せる。残る一形式は表示した対数型となる。

直線・横断的非退化二次曲線の積の層は既約八次元である。
$a$ を固定した二つの混合核はそれぞれ階数二十、十四のベクトル束であり、
$a$ も動かす一般層では九次元基底上の階数十のベクトル束となる。
これらの族の次元は従って $28,22,19$ である。
定理 \ref{thm:lineconic-local-v32} により、それぞれの族には Poisson スキーム全体の
一つの枝を占める開集合がある。よって各閉包は既約成分である。
最初の二つは次元が異なり、最後の成分は一般の $a$ を持つので、三つは相異なる。
\end{proof}

\begin{remark}[独立した滑らかな点の証明書]
付属検証器は、同じ三成分を別途全 Jacobi 行列で照合する。
$c=c_0$ として
\[
 (a,F)=\bigl(2,(0,0,y^4)\bigr),\quad
 \bigl(-1,(-y^4,0,2xy^3)\bigr),\quad
 \bigl(1/4,0\bigr)
\]
における三十八行五十三列の Jacobi 行列の階数は、それぞれ $25,31,34$ である。
対応する正方小行列式として
\[
 61917364224,\qquad -1410554953728,\qquad
 -\frac{5791625339653125}{562949953421312}
\]
を得る。行・列の指定は JSON 証明書に含める。
成分性の証明は有限個の接空間計算だけに依存せず、上記の形式切片と全パラメータの行列計算による。
\end{remark}

\section{接線型の第二共鳴と一般点で二重な成分}
\label{sec:tangent-double-v34}

この節では重みを $(1,1,1,3)$ に固定する。
Poisson スキームを $\mathcal P$ とし、その座標と混合式を
\[
 (c,D,[F])\in S_3\oplus\mathfrak{sl}_3\oplus(S_4^{\oplus3}/xS_3),
 \qquad M(c,D;F)=2F(c)-(x\times D)\cdot F-c\operatorname{div}F
\]
と書く。方程式は $D(c)=M(c,D;F)=0$ であり、イデアルは根基化しない。
前節で残した接線型の第二の共鳴 $a=-3/2$ は、
既知成分の内部の核の増大だけではなく、別の非被約成分を与える。

\begin{theorem}[一般点で二重な十八次元成分]
\label{thm:tangent-double-v34}
\[
 q=xz-y^2,\quad c_0=xq,\quad N_0=x\partial_y+2y\partial_z,\quad
 D_-=-\tfrac32N_0,\quad F_-=(0,\tfrac12y^3z,y^2z^2)
\]
と置く。組 $(c_0,D_-)$ の $\operatorname{SL}_3$ 軌道上で全混合核を動かした
十八次元族の閉包 $\mathcal E_{-3/2}$ は、$\mathcal P$ の既約成分である。
その一般点の局所環は
\[
 K[\epsilon]/(\epsilon^2)
\]
であり、一般点の長さは二である。被約成分は重み付き座標変換群の稠密軌道を持つ。

より正確には、$P_-=(c_0,D_-,0)$ における全完備局所環は
\begin{equation}
 \widehat{\mathcal O}_{\mathcal P,P_-}
 \simeq
 \mathbb C[[s_1,\ldots,s_{17},k,b,v]]/(kv,b^2v).
 \label{eq:tangent-double-origin-v34}
\end{equation}
$v=0$ の枝は十九次元成分 $L(1,2,3)$ であり、もう一方は
$\mathcal E_{-3/2}$ の二重構造である。
点 $Q_-=(c_0,D_-,F_-)$ では
\begin{equation}
 \widehat{\mathcal O}_{\mathcal P,Q_-}
 \simeq\mathbb C[[s_1,\ldots,s_{18},b]]/(b^2).
 \label{eq:tangent-double-nonzero-v34}
\end{equation}
従ってこの成分は、既知の九成分と相異なる。
重み三の Poisson スキームには少なくとも十個の既約成分がある。
\end{theorem}

この定理での長さ二は、混合行列の一点での階数だけから推測したものではない。
以下の消去では、二つの基底変数を含む全形式近傍の式を求める。

\subsection{非零の接線型に共通する基底の形式切断}

\begin{lemma}
\label{lem:tangent-base-slice-v34}
$A\ne0$ を固定する。基底 $D(c)=0$ の $(c_0,AN_0)$ における形式近傍は、
七個の滑らかな軌道変数と、次の二変数族の積である。
\begin{equation}
 q_b=xz-y^2+\frac{b}{2A}z^2,\qquad
 c=u xq_b,\qquad D=(0,Ax+bz,2Ay),\qquad (u,b)=(1,0).
 \label{eq:tangent-base-slice-v34}
\end{equation}
特に基底はこの点で滑らかな九次元スキームである。
\end{lemma}

\begin{proof}
$D$ は正則冪零行列なので、共役作用の六変数を除き、形式的に
\[
 D=(tz,Ax+bz,2Ay)
\]
と書ける。実際、その二つの行列不変量は
$\operatorname{tr}(D^2)=4Ab$、$\det D=2A^2t$ であり、
正則行列の共役軌道への横断切断を与える。

三次係数を単項式 $x^iy^jz^{3-i-j}$ の順に並べる。
$x^2z,x^3$ の係数を $u,r$ として残し、他の八係数を消去する。
使用する八次小行列式は $b=t=0$ で $-1152A^8$ なので単元である。
消去後の三次式を $4A^2$ 倍すると
\begin{align*}
 4A^2c={}&4A^2(ux^2z-uxy^2+rx^3)+2Abu xz^2
       +2t^2r z^3+t(Au-6br)yz^2\\
       &+8Atr y^3-12Atr xyz.
\end{align*}
これに $D$ を作用させた残差は
\[
 3bt(Au-2br)z^3+9At(Au-2br)xz^2
\]
のみである。$Au-2br$ は対象の点で単元なので、残る条件は正確に $t=0$ となる。
従って $c=uxq_b+rx^3$ である。

$D$ の中心化代数の場 $D^2-\operatorname{tr}(D^2)I/3$ は、
$b=0$ で $c$ の $x^3$ 係数を $2A^2u$ だけ動かす。
この係数は単元なので、中心化群の一変数によって $r$ を形式的に零にできる。
以上で七軌道変数と式 \eqref{eq:tangent-base-slice-v34} を得る。
独立の確認として、$D(c)$ の九次小行列式は $2592A^8$、
七軌道方向と $\partial_u,\partial_b$ の接写像の九次小行列式は $-6$ となる。
従ってこの切断は被約点だけの標準形ではなく、全形式近傍を記述している。
\end{proof}

$b$ は幾何的には直線と二次曲線の接触を解く変数である。
$b=0$ では $x=0$ と $q_b=0$ は接し、$b\ne0$ では二点で横断的に交わる。
また $u$ は三次式と線形場との相対的な係数を記録する。

\subsection{二重構造を決める正確な混合方程式}

補題の $A=-3/2$ を取る。二次場の係数は
\[
 p=\frac{\nabla c-x\times D}{3}
\]
であり、三次式による $w$ 移動の混合項
$p\times\nabla g$、$g\in S_3$ は、全 $(u,b)$ 上で実際の核を成す。
その十方向は $(1,0)$ で独立なので、形式近傍で直和因子として除ける。
残る二十五列に対し、二十四の単元行・列を消去する。
最後の一方向には次の有理係数の四次場を取れる。
\begin{align}
 \widehat F_1={}&0,\nonumber\\
 \widehat F_2={}&\frac{2u}{u+3}y^3z
 -\frac{(u-1)(4u+3)}{2(u+1)(u+3)}xyz^2
 -\frac{b(7u^2+6u-9)}{3(u+1)(u+3)(5u-3)}yz^3,\label{eq:tangent-double-Fhat-v34}\\
 \widehat F_3={}&y^2z^2+\frac{bu(9-5u)}{3(u+3)(5u-3)}z^4.\nonumber
\end{align}
全ての分母は $(u,b)=(1,0)$ で単元であり、同点で $\widehat F=F_-$ となる。
直接の係数比較により、正確な恒等式
\begin{align}
 M(c,D;\widehat F)
 ={}&\frac{2b^2u(u-3)(4u+3)}{9(u+1)(u+3)(5u-3)}xz^5\nonumber\\
 &+\frac{u(u-1)(4u+3)}{2(u+1)(u+3)}x^3z^3
 \label{eq:tangent-double-Schur-v34}
\end{align}
を得る。

この式が Schur 消去の全残差であることも、小行列式で保証できる。
Euler 商の基底は成分順、単項式順とし、第一成分の $x$ で割れる単項式を除く。
零始まりで、単元列は
\[
 \{0,1,\ldots,20,24,31,34\}
\]
であり、単元行は $S_6$ の単項式順の
\[
 \{2,3,4,5,6,8,9,10,11,12,13,14,15,16,17,19,20,21,22,23,24,25,26,27\}
\]
である。この二十四次小行列式は $-136193423625$。
十個の $w$ 移動、これらの二十四列、$\widehat F$ の三十五次基底変換の行列式は
$-16875$ である。式 \eqref{eq:tangent-double-Schur-v34} の二行は上の単元行に含まれず、
その二行以外の係数は恒等的に零である。
従ってこの基底変換と行消去は形式近傍で可逆であり、隠れた高次方程式は残らない。

$\widehat F$ の係数を $v$、$k=u-1$ とする。
$b^2$ と $k$ の係数は原点でそれぞれ $-7/36,7/16$ なので単元である。
従って混合条件は正確に
\[
 b^2v=kv=0
\]
へ変形できる。七軌道変数と十個の $w$ 移動を戻せば
式 \eqref{eq:tangent-double-origin-v34} を得る。
ここで全ての基底変換と単元小行列式は $(u,b)$ だけに依存し、
消去は $F$ に関して線形である。従って同じ操作は基底の形式近傍と全混合空間の積上で成立し、
$v=0$ だけでなく $v=1$ の近傍にも適用できる。

\begin{proof}[定理 \ref{thm:tangent-double-v34} の残りの証明]
$v\ne0$ の点では $k=0,b^2=0$ となるので、
式 \eqref{eq:tangent-double-nonzero-v34} が従う。
$(c_0,D_-)$ の軌道は七次元で、混合核の階数は十一。
従ってその核束の全体は既約十八次元族である。
上の局所環から、その閉包は全スキームの一つの既約成分であり、一般点の長さが二と分かる。

混合核は十個の $w$ 移動と $F_-$ で生成される。
残余係数が非零なら、移動によって十方向を除き、$w$ の非零定数倍で残余係数を一にできる。
従って被約成分には稠密な座標変換軌道がある。
従来の八成分は次元が十九以上であり、既知の接線型十八次元成分
$\mathcal E_{-5/2}$ は一般点で滑らかである。
いずれも一般点で二重な $\mathcal E_{-3/2}$ と一致しない。

最後に $v=0$ の枝を同定する。$b\ne0$ では直線・二次曲線が横断し、
$u$ は $-3/2$ 近傍の非共鳴な比を記録するため、混合核は全て $w$ 移動から来る。
これは十九次元対数型成分 $L(1,2,3)$ の稠密部分である。
よってその閉包が式 \eqref{eq:tangent-double-origin-v34} の $v=0$ の枝である。
\end{proof}

さらに
\[
 (kv,b^2v)=(v)\cap(k,b^2)
\]
は primary 分解である。二枝の被約交叉は滑らかな十七次元で、primary 成分同士の交叉は
その上の $\mathbb C[[b]]/(b^2)$ という二重構造を持つ。
また $Q_-$ で全 Jacobi 行列の階数は三十四、接空間次元は十九であり、
三十四次小行列式 $27799801630335/64$ が独立の確認を与える。

\paragraph{二次で止まる接触解消方向.}
$u=1$ における式 \eqref{eq:tangent-double-Fhat-v34} は単純化し、
\[
 c_b=x(xz-y^2-bz^2/3),\quad
 D_b=(0,-3x/2+bz,-3y),\quad
 F_b=(0,y^3z/2-byz^3/12,y^2z^2+bz^4/6)
\]
となる。この族は正確に
\[
 D_b(c_b)=0,\qquad M(c_b,D_b;F_b)=-\frac7{36}b^2xz^5
\]
を満たす。従って接触を解く一次方向は存在するが、この点では二次で止まる。
この障害を消して被約化してしまうと、新成分の長さ二は見えなくなる。

\subsection{他の接線共鳴との比較}

同じ基底の切断は、他の非零 $A$ にも適用できる。
$P_A=(c_0,AN_0,0)$ での全形式局所環を、次の表にまとめる。
$k=u-1$ であり、$s_1,\ldots,s_{17}$ は七軌道変数と十移動変数である。
\[
\begin{array}{c|c|c}
 A&\text{滑らかな十七変数を除いた局所環}&\text{通る大域成分}\\\hline
 2&\mathbb C[[b,k,v_1,\ldots,v_{10}]]/(kv_1,\ldots,kv_{10})
   &L(1,2,3),\ R(1,5)\\
 -1&\mathbb C[[b,k,v_1,\ldots,v_4]]/(kv_1,\ldots,kv_4)
   &L(1,2,3),\ R(2,4)\\
 -5/2&\mathbb C[[b,k,v]]/(kv,bv)&L(1,2,3),\ \mathcal E_{-5/2}\\
 -3/2&\mathbb C[[b,k,v]]/(kv,b^2v)&L(1,2,3),\ \mathcal E_{-3/2}
\end{array}
\]

$A=-5/2$ の場合には
\[
 \widehat F=(0,4uyz^3/(3u+5),z^4)
\]
を取れば、全残差は
\[
 M(c,D;\widehat F)
 =\frac{b(u-5)(4u+5)}{5(3u+5)}xz^5
  -\frac{5(u-1)(4u+5)}{2(3u+5)}x^2z^4
\]
となる。二つの係数は $b,k$ の単元倍なので、こちらの十八次元成分は被約である。

$A=2,-1$ については、それぞれ
\[
 \nabla x\times\nabla S_5,\qquad \nabla q_b\times\nabla S_4
\]
が $u=1$ 上の全混合核を与える。
十移動方向を除く残余次元は十、四である。
これらの残余方向を全形式基底へ延長したとき、混合像は全て $u-1$ で割れる。
その商を単元像で割った行列は、原点でそれぞれ階数十、四を持つ。
従って追加の基本因子は全て一位の $k$ であり、表の被約な局所環を得る。
この階数にはそれぞれ二十五次小行列式
$-8321693751705600$、$-5015306502144$ を用いる。

$A\ne0,2,-1,-5/2,-3/2$ では、既に求めた全パラメータの重みブロック表により
混合行列の階数は二十五である。基底切断の一般点では十個の独立な移動方向があるので、
二十六次小行列式は恒等的に消える。一方、対象点の二十五次小行列式は単元である。
従って混合行列の階数はこの基底近傍で一定で、核は階数十のベクトル束となる。
その一般点は $L(1,2,3)$ に属すので、$P_A$ は同成分上の滑らかな十九次元の点である。
これは有限個の非共鳴値を試す議論ではない。

この説明では $A=1$ を含めている。この値では $p=(q,0,0)$ なので
$p\times\nabla x^3=0$ となり、実際の移動の像は九次元である。
不足する核方向は $x^3N_0$ であり、
\[
 p_a\times\nabla x^3=(1-a)x^3N_0
\]
によって近傍の移動方向の有限な極限として回収できる。
従って移動の階数と混合核の次元をこの点で混同してはならない。

$A=0$ では基底補題は用いないが、固定した $(c_0,0)$ 上の混合核は全て
$L(1,2,3)$ の閉包に含まれることを明示できる。
この核は九移動方向と $c_0N_0$ で生成され、
\[
 p_a\times\nabla c_0=-a c_0N_0,\qquad
 p_a=(\nabla c_0-a x\times N_0)/3
\]
である。よって $a\ne0$ での移動 $g=-c_0/a$ の混合項は $c_0N_0$ であり、
$a\to0$ の有限な極限が残る方向を覆う。
ここでは $D=0$ における全局所環の分類までは主張しない。

\subsection{二次元 Lie 作用と有理第一積分}

\begin{proposition}
点 $Q_-$ の発散零持ち上げは
\[
 \widehat P_-=H\wedge N,\qquad
 H=\tfrac1{12}(7x\partial_x+y\partial_y-5z\partial_z-9w\partial_w),
 \qquad N=x\partial_y+2y\partial_z-y^2z\partial_w
\]
である。$[H,N]=N/2$、$\operatorname{div}H=-1/2$、$\operatorname{div}N=0$ を満たす。
また
\[
 J=x^2w+\frac13xy^3z-\frac2{15}y^5
   =x^2w+\frac13q y^3+\frac15y^5
\]
に対し $N(q)=N(J)=0$、$H(q)=q/6$、$H(J)=5J/12$ である。
対応する水平一形式 $\alpha$ は
\[
 2q\,dJ-5J\,dq=2x\alpha
\]
を満たし、$J^2/q^5$ が有理第一積分となる。
\end{proposition}

\begin{proof}
持ち上げの六係数、Lie 括弧、発散、および一形式の四係数を直接比較すればよい。
$N$ は補助次数 $(0,1,2,5)$ を一つ下げるため局所冪零である。
\end{proof}

これは次数 $(2,5)$ の特殊 pencil から共通因子 $x$ を除いた次数六の一形式であり、
通常の $R(2,5)$ と同一視しない。
重み付き座標変更 $w'=-5w-yz^2$ は $H$ と可換であり、
$N(w')=-qz$ へ変える。この表示は、同じ成分の量子化と不変式環を調べる際にも使える。

\paragraph{厳密検証と射程.}
\texttt{verify\_tangent\_resonance\_v34.py} は Python 標準ライブラリのみで動く。
分母を払った全係数恒等式、基底変換と単元部分の小行列式、四共鳴の局所式、
全 Jacobi 行列、Lie 表示と第一積分を有理数上で確認する。
通常実行では入力や保存済み証明書を書き換えない。
本節は接線型の非零線形場の形式近傍を決定するものであり、
全ての非被約三次式の層や全 Poisson 成分の完全分類を主張するものではない。

% Insert after the generic log-canonical component of Paper VIII.
% Requires the v31 notation and the bibliography item LWW.
\section{対数型の共鳴壁と23次元成分}
\label{subsec:resonance23-v32}

一般対数型の非消滅条件を外すと、四次項が五方向同時に自由になる壁が現れる。
その壁は、既知の19次元成分に接する別の23次元成分を作る。
以下ではPoisson構造の53変数のスキームを $\mathscr P$ と書く。

固有座標を $(x,y,z)$ とし、
\[
 p=(\alpha yz,\beta zx,\gamma xy),\qquad
 t=\alpha+\beta+\gamma,\qquad c=txyz
\]
とする。$D=\operatorname{curl}p$ の固有値は
\[
 (\delta_1,\delta_2,\delta_3)
 =(\gamma-\beta,\alpha-\gamma,\beta-\alpha)
\]
である。次の開条件を課す。
\begin{enumerate}
\item $D$ は正則半単純で、$\ker(D|_{S_3})=\mathbb Cxyz$、$t\neq0$。
\item $3(\alpha,\beta,\gamma)-tm\neq0$ がすべての
      $m\in\mathbb Z_{\geq0}^3$、$|m|=3$ に対して成り立つ。
\item $a_2=3\beta+t\neq0$、$a_3=3\gamma+t\neq0$。
\end{enumerate}
この開集合の中で $a=a_1=3\alpha+t$ を考える。
壁 $a=0$ は空でなく、例えば $(\alpha,\beta,\gamma)=(-1,0,4)$ は
$t=3$、$D$ の固有値 $(4,-5,1)$ を与え、上の条件を満たす。

\begin{theorem}[23次元共鳴成分と19次元成分との交叉]
\label{thm:resonance23-v32}
上の開集合において、$a=0$ 上の混合方程式の核はEuler商で15次元である。
三次式による $w$ の移動を行うと、すべての解は一意に
\[
 [F]=[(f(y,z),0,0)],\qquad f\in\mathbb C[y,z]_4
\]
という五次元の残余を持つ標準形になる。
ここで一意性は、混合項を三次式による移動で割った商についての主張である。

壁上の残余 $f=0$ の任意の点で、$\mathscr P$ の完備局所環は
\[
 \mathbb C[[u_1,\ldots,u_{18},a,z_0,\ldots,z_4]]
       /(az_0,\ldots,az_4)
\]
と同型である。したがって、この近傍には既知の19次元対数型成分と
新しい23次元成分の二つだけがあり、両者は滑らかな18次元部分で交わる。
この局所スキームは被約である。

固有座標の置換で得られる三つの共鳴壁の像は、全体として一つの
既約23次元成分 $\mathscr C_{114}$ をなす。
その一般点は $a=0$、$f\neq0$ であり、$\mathscr P$ の滑らかな点である。
\end{theorem}

\begin{proof}
正則半単純線形場の固有値を順序付ける有限\'etale被覆を取り、
固有線束を局所的に自明化する。条件1により
$D(c)=0$ は $c=txyz$ を強制し、その基底は滑らかな9次元である。
これは固有線の配置6次元、跡零固有値2次元、$t$ の1次元からなる。

四次ベクトル場の文字表示 $F_i=x_i x^m v_i$ を用いると、混合方程式は
\[
 \bigl(3(\alpha,\beta,\gamma)-tm\bigr)\cdot v=0
\]
である。非負の10文字ではEuler商は二次元で、条件2によって方程式の階数は1。
従って各文字につき一つの自由変数が残り、計10次元となる。
負の文字は $m_i=-1$ の場合であり、各 $i$ について五つある。
条件3により $i=2,3$ の10変数は零になる。
$i=1$ の五変数はちょうど $F_1\in\mathbb C[y,z]_4$ の係数であり、
その方程式は同じ係数 $a$ を持つ
\[
 az_0=az_1=\cdots=az_4=0
\]
である。異なる文字は異なる六次単項式へ写るため、追加の結合条件はない。

$a$ は9次元基底の滑らかな局所座標の一つである。
残る基底8変数と非負文字の10自由変数を $u_1,\ldots,u_{18}$ と取れば、
表示した完備局所環を得る。イデアルの等式
\[
 (az_0,\ldots,az_4)=(a)\cap(z_0,\ldots,z_4)
\]
により、被約性と二つの枝、交叉の次元がすべて従う。
$a\neq0$ では $z_i=0$ であるから、19次元の枝は既知の一般対数型成分である。
$z_i$ のいずれかが非零なら $a=0$ が強制され、その近傍は滑らかな23次元になる。

非負文字の一自由度は、一般対数型と同じ計算により三次式による $w$ の移動で
すべて消去できる。ここでは $t\neq0$ と条件2を使う。
移動の像には負の文字がないため、残余はちょうど二変数四次式の五次元である。

最後に大域的な既約性を示す。壁は
\[
 (\alpha,\beta,\gamma)=(\alpha,\beta,-4\alpha-\beta)
\]
という二次元平面の非空な開集合である。その上の15次元の混合核はベクトル束である。
これを $\operatorname{GL}_3$ で動かした像は既約であり、一般のファイバーは
固有基底の対角変更に由来する3次元のものである。従って像の次元は
\[
 9+2+15-3=23
\]
となる。固有座標の置換も $\operatorname{GL}_3$ の作用に含まれるので、
三つの壁は同じ既約閉包に入る。
上で得た滑らかな23次元点を含むことから、この閉包は
$\mathscr P$ の一つの既約成分である。
\end{proof}

\begin{remark}[検証と量化]
基準点 $(\alpha,\beta,\gamma)=(-1,0,4)$ では混合写像の階数は20、
Euler商の核は15次元、三次移動の像は10次元である。
$F=0$ における全38方程式のJacobian階数は29、接空間は24次元であり、
$F=(y^4,0,0)$ では階数30、接空間は23次元となる。
これらは有理数上で検証できるが、成分性と局所環の表示の根拠は上の文字ごとの証明である。
被約性の結論はこの共鳴開集合に限られ、全Poisson錐の被約性を意味しない。
\end{remark}

\subsection{重み付き対数型 $L(1,1,4)$ との同定}
\label{subsec:log114-v32}

Poisson構造から得られる基本1形式を
\[
 \eta=\iota_E\iota_P
       (dx\wedge dy\wedge dz\wedge dw)
      =c\,dw-3w(p_1dx+p_2dy+p_3dz)-(x\times F)\cdot dx
\]
とする。Euler場を因子に持つ二重ベクトル場はこの式に寄与しないので、
最後の表示には $w$ に依存しない代表 $P$ を用いてよい。

共鳴標準形 $F=(f(y,z),0,0)$ と $t=-3\alpha$ に対し、
\[
 Q=xw+g(y,z),\qquad
 (\lambda_y,\lambda_z,\lambda_Q)=(-3\beta,-3\gamma,t)
\]
と置くと、
\[
 \eta=yzQ\left(\lambda_y\frac{dy}{y}+
                   \lambda_z\frac{dz}{z}+
                   \lambda_Q\frac{dQ}{Q}\right)
\]
となる必要十分条件は
\[
 f=3\beta g-t\,y\partial_y g
\]
である。重み付きEuler条件は
$\lambda_y+\lambda_z+4\lambda_Q=0$ であり、これは $3\alpha+t=0$ と同じである。
二変数四次単項式 $y^kz^{4-k}$ 上の係数は $3\beta-tk$ なので、
これら五つが非零となる一般点では $g$ が一意に存在する。
例えば $(\alpha,\beta,\gamma)=(-1,-2,6)$ はすべての開条件を満たす。

逆に、一般の次数 $(1,1,4)$ の三因子を持つ重み付き対数型1形式を取る。
一次因子を $y,z$ とし、四次因子の $w$ の係数がこれらと独立である開集合で
その係数を $x$ と選ぶ。四次因子は $xw+H_4(x,y,z)$ と書ける。
$H_4$ の $x$ で割り切れる部分を三次式による $w$ の移動で消すと
$Q=xw+g(y,z)$ となり、上の表示に帰着する。
従って $\mathscr C_{114}$ は、この重み付き対数型 $L(1,1,4)$ の閉包に等しい。
この同定は一般点の幾何的な記述であり、共鳴の追加壁で同じ標準化が一意に
成立するという主張ではない。


% May be used in the article or in research notes; this is a complete D=0 classification.
\section{三次形式に付随する混合項と留数}
\label{sec:cubic-residues-v32}

$D=0$ の混合方程式は、三次形式が特異であっても閉有理一形式によって完全に解ける。
以下で $S=\mathbb C[x,y,z]$ とし、$0\ne c\in S_3$ とする。

\begin{theorem}[全ての非零三次形式に対する混合核]
\label{thm:cubic-residues-v32}
$c=\prod_{i=1}^r f_i^{m_i}$ を相異なる既約斉次形式への分解とし、
\[
 c_{\mathrm{red}}=\prod_{i=1}^r f_i,\qquad
 d=\deg c_{\mathrm{red}},\qquad h=\frac{c^2}{c_{\mathrm{red}}}\in S_{6-d}
\]
と置く。Euler 商における方程式
\[
 2F(c)-c\operatorname{div}F=0
\]
の全解は、閉有理一形式
\begin{equation}
 \gamma_F:=\frac{(F\times x)\cdot dx}{c^2}
 =d\!\left(\frac{p}{h}\right)+\sum_{i=1}^r\lambda_i\frac{df_i}{f_i},
 \qquad
 p\in S_{6-d}/\mathbb Ch,\quad
 \sum_i(\deg f_i)\lambda_i=0
 \label{eq:cubic-residue-normal-v32}
\end{equation}
と一対一に対応する。従って混合核の次元は
\[
 \dim\ker M_{c,0}=\dim S_{6-d}+r-2.
\]
三次式による $w$ の移動の像は常に九次元である。
\end{theorem}

\begin{proof}
$\beta=(F\times x)\cdot dx$ と置く。
ベクトル解析の恒等式と Euler の式から
\[
 \operatorname{curl}(F\times x)=6F-x\operatorname{div}F,
 \qquad
 \nabla c\times(F\times x)=3cF-xF(c)
\]
を得る。従って
\[
 d(\beta/c^2)=0
 \quad\Longleftrightarrow\quad
 2F(c)-c\operatorname{div}F=0.
\]
また $\beta$ は係数が五次で $\iota_{E_x}\beta=0$ を満たす一形式である。
逆にこの条件を満たす任意の多項式一形式は、Koszul 完全性により $\beta=(F\times x)\cdot dx$
と書け、$F$ は Euler 同値を除いて一意である。
発散零の代表は $F=\operatorname{curl}\beta/6$ で与えられる。

$\gamma_F$ は $\mathbb P^2$ 上の閉有理一形式であり、$f_i=0$ に沿う極の位数は高々 $2m_i$ である。
閉有理一形式の留数分解と $H^1(\mathbb P^2,\mathbb C)=0$ により、
\[
 \gamma_F=df+\sum_i\lambda_i\frac{df_i}{f_i}
\]
と書ける。ここで留数は定数であり、射影的水平性から
$\sum_i(\deg f_i)\lambda_i=0$ である。
留数分解および第二種の閉形式が $H^1$ の類を定めることについては
\cite[\S5, Corollary~5.6]{PereiraClosedForms} を参照されたい。
本例では第二種の部分の周期が全て消えるため、原始関数は大域的有理関数となる。

$df$ の極の位数が高々 $2m_i$ であることから、$f$ の極の位数は高々 $2m_i-1$ である。
他の極は持ち得ない。従って $f=p/h$、$p\in S_{6-d}$ と書け、$p$ は $\mathbb Ch$ を除いて一意である。
逆に式 \eqref{eq:cubic-residue-normal-v32} の右辺に $c^2$ を掛けると、
極が消え、係数五次で水平な多項式一形式となる。
これを $F\times x$ として戻せば混合方程式の解が得られる。
留数と、定数を除いた原始関数がそれぞれ一意であることから直和分解と次元公式が従う。

$F=\nabla c\times\nabla g$、$g\in S_3$ に対して
\[
 \gamma_F=3d(g/c)
\]
である。その核は $g\in\mathbb Cc$ なので、この像は非被約な $c$ の場合も常に九次元である。
\end{proof}

\begin{corollary}[被約性と混合方向の増大]
$c$ が被約なら $\dim\ker M_{c,0}=r+8$ であり、
$w$ の三次移動で消えない方向は既約因子数より一つ少ない $r-1$ 次元である。
非被約の場合も含めると次の表になる。
\[
\begin{array}{l|c|c|c}
 c\text{ の型}&(d,r)&\dim\ker M_{c,0}&\text{移動後に残る次元}\\\hline
 \text{既約（滑らか、節点、尖点）}&(3,1)&9&0\\
 \text{相異なる直線と非退化二次曲線}&(3,2)&10&1\\
 \text{相異なる三直線}&(3,3)&11&2\\
 \ell^2m\quad(\ell\not\parallel m)&(2,2)&15&6\\
 \ell^3&(1,1)&20&11
\end{array}
\]
直線と二次曲線の接触の有無、三直線の共点性は、この混合核の次元を変えない。
\end{corollary}

この定理は $D=0$ のファイバーを完全に記述するが、各層が全 Poisson スキームの
既約成分をなすことまでは述べない。成分性や非被約性を判定するには、
$c$ と $D$ を同時に動かした方程式を調べる必要がある。

% Section for the ambient-completion paper.  No external packages beyond the
% inherited mathematical preamble are required.
\section{三角形の形式近傍と非被約構造}
\label{sec:triangle-local-v32}

本節では $\kk$ を標数零の代数閉体とし、$\mathscr P$ を
\[
 D(c)=0,\qquad 2F(c)-(x\times D)\cdot F-c\operatorname{div}F=0
\]
で定義される五十三次元空間内の Poisson 構造のスキームとする。
三角形の点
\[
 p_\triangle=(x_2x_3,x_3x_1,x_1x_2),\qquad
 c_\triangle=3x_1x_2x_3,\quad D=0,\quad [F]=0
\]
の近傍を、被約化せずに記述する。

\begin{theorem}[三角形の完備局所環]
\label{thm:triangle-local-v32}
$\mathscr P$ の上の点における完備局所環は、非標準的に
\[
 \kk[[z_1,\ldots,z_{16}]]\widehat\otimes_\kk \mathcal T
\]
と同型である。ただし
\[
 \mathcal T=
 \frac{\kk[[a_1,a_2,a_3,\delta_1,\delta_2,\delta_3,
                   \rho_1,\rho_2,\rho_3]]}
 {\left(
  \sum_i\delta_i,\ \sum_i\rho_i,\
  a_i\delta_i,\ a_i^2\rho_i\ (i=1,2,3),\
  \delta_1\rho_2-\delta_2\rho_1
 \right)}.
\]
従って、この点における局所接空間の次元は $23$ である。
\end{theorem}

\begin{proof}
証明を、三次形式の slice、線形場の消去、混合写像の消去に分ける。

\smallskip\noindent
\textbf{三次形式の slice.}
$c=x\cdot p$ は線形座標変換に関して三次形式の行列式捻りとして変換する。
従って三角形における線形座標軌道は六次元であり、
その接空間は六つの $x_i^2x_j$ ($i\ne j$) で張られる。
Poisson テンソル全体の非零スカラー倍を加えると、
$x_1x_2x_3$ の係数も独立に動かせる。
この七方向と三つの純三次式を合わせた接写像は同型である。
形式逆写像定理により、七つの滑らかな変数を分離した後、
\[
 c=3xyz+a x^3+b y^3+c' z^3
 \tag{\(\ast\)}\label{eq:triangle-cubic-slice-v32}
\]
と置いてよい。以下、$x=x_1,y=x_2,z=x_3$、
$(a,b,c')=(a_1,a_2,a_3)$ とする。
この正規化は単に通常の三次形式の $GL_3$ 軌道を使うものではなく、
Poisson テンソルのスカラー倍一方向も使っている。

\smallskip\noindent
\textbf{線形場の消去.}
$D(c)=0$ の六つの混合三次係数は、$D$ の六つの非対角成分に関する
線形方程式であり、$a=b=c'=0$ では係数行列が可逆である。
対角成分はこれらの六係数に寄与しないので、非対角成分はすべて零となる。
従って
\[
 D=d x\partial_x+e y\partial_y-(d+e)z\partial_z,
\]
かつ残る基底方程式は
\[
 ad=0,\qquad be=0,\qquad c'(d+e)=0.
 \tag{\(\ast\ast\)}\label{eq:triangle-base-v32}
\]
すなわち $(\delta_1,\delta_2,\delta_3)=(d,e,-d-e)$ である。
以下ではこの基底環上で混合方程式を計算する。

\smallskip\noindent
\textbf{混合写像の文字分解.}
Euler 同値 $F\sim F+gx$ を用いて、$F_1$ に $x$ で割り切れる単項式が
現れない代表を選ぶ。混合写像の源は三十五次元、標的は $S_6$ の二十八次元である。
行列式一の対角三乗根の作用に関して、三つの文字ブロックに分かれる。
それぞれの源・標的の次元と三角形での階数は
\[
 (11,10,6),\qquad(12,9,9),\qquad(12,9,9)
\]
である。後二つは全射であるから、形式近傍でも十八個の変数を消去でき、
合計六つの自由変数が残る。
これらの階数は単項式ごとにも読める。
三角形で $F_i=x_ix^m v_i$ と書くと、混合係数は
$3(1-m)\cdot v$ である。$m\geq0$、$|m|=3$、$m\ne(1,1,1)$ の文字は
Euler 商上で階数一、負指数を一つ持つ文字の係数は $6$ である。

残る十一変数の代表を
\begin{align*}
 F_1&=A y^2z^2,\\
 F_2&=B x^3y+C y^4+K yz^3+u xy^2z+E x^2z^2,\\
 F_3&=H_1x^3z+H_2y^3z+H_3z^4+v xyz^2+G x^2y^2
\end{align*}
と置く。ここで $E,G$ は係数の記号であり、Euler 場とは無関係である。
次の五つの元はすべて原点で単元である。
\[
 L=3-2d-e,\quad T=6+2d+e,\quad U=6+d+2e,
 \quad V=6-2d-e,\quad W=6+d-e.
\]
混合式の $x^4yz,xy^4z,xyz^4,y^3z^3,x^3z^3,x^3y^3$ の六係数から、
順に
\begin{align*}
 B&=\frac{2a(u+v)-(d-e+3)H_1}{L},\\
 C&=\frac{4bu+(d-e+3)H_2-2bv}{T},\\
 K&=\frac{2c'u-4c'v-(d-e-6)H_3}{L},\\
 A&=\frac{4c'C-5bK-5c'H_2+4bH_3}{U},\\
 E&=\frac{c'B+aK-5c'H_1+4aH_3}{V},\\
 G&=\frac{-5bB+4aC+bH_1+aH_2}{W}
\end{align*}
を得る。この消去は可逆であり、$u,v,H_1,H_2,H_3$ が残る。
残る $x^6,y^6,z^6,x^2y^2z^2$ の四係数は、
\eqref{eq:triangle-base-v32} を法として
\[
 -\frac{2a^2(u+v)}{L},\qquad
 \frac{4b^2(2u-v)}{T},\qquad
 -\frac{2(c')^2(u-2v)}{L},\qquad
 -(2d+e)u+(d-e)v
\]
となる。最後の等式を直接確認するには
$aA+bE+c'G=0$ を示せばよい。
この左辺の各項には $ab,ac',bc'$ のいずれかが含まれる。
基底方程式から、これら三つの積はすべて $d,e$ によって殺されるので、
分母も含めて $d=e=0$ と置いて計算できる。
上の六式を代入すると $abc'$ の係数と $H_i$ の係数はそれぞれ相殺する。

ここで
\[
 \rho_1=-u-v,\qquad\rho_2=2u-v,\qquad\rho_3=-u+2v
\]
と置く。これは $\sum\rho_i=0$ の平面への線形同型である。
四つの残余方程式は、単元を除いて
\[
 a_i^2\rho_i=0\quad(i=1,2,3),\qquad
 \delta_1\rho_2-\delta_2\rho_1=0
\]
に一致する。$H_1,H_2,H_3$ は自由である。
自由変数は先の七方向、非自明文字からの六方向、最後の三方向、合計十六となる。
線形関係二つを消した $\mathcal T$ は七変数と二次以上の関係だけを持つので、
接次元は $16+7=23$ である。
\end{proof}

\begin{corollary}[Jacobi イデアルの非被約性]
\label{cor:triangle-nonreduced-v32}
$\mathscr P$ とその原点完備化は非被約である。
\end{corollary}

\begin{proof}
$\mathcal T$ の元 $a_1\rho_1$ は非零であり、
\[
 (a_1\rho_1)^2=\rho_1(a_1^2\rho_1)=0
\]
を満たす。非零性は最低次数から分かる。
線形関係を消した後、定義イデアルの二次部分はすべて $\delta_1,\delta_2$ を含むため、
$a_1\rho_1$ を含まない。
別の確認として、$\kk[\epsilon]/(\epsilon^3)$ に
\[
 a_1=\rho_1=\epsilon,\qquad
 a_2=a_3=\delta_1=\delta_2=\rho_2=0
\]
と代入すると全関係は成立するが、$a_1\rho_1=\epsilon^2\ne0$ である。

三角形での非被約な完備局所環だけから任意の環について結論するのではなく、
ここでは上の消去が単元の逆元とエタール座標変換で行われたことを用いる。
すなわち代数的なエタール局所模型にも同じ平方零元があり、
エタール射の被約性の保存から $\mathscr P$ 自体が非被約となる。
あるいは有限型 $\kk$-代数の優秀性を用いてもよい。

Jacobi イデアルは正次数の斉次イデアルである。
その商環に非零の冪零元があれば、冪零根基の斉次性から非零の斉次冪零元を取れる。
この元は斉次極大イデアルでの局所化でも消えず、原点完備化でも消えない。
従って原点の完備局所環も非被約である。
別稿\cite{VKAmbient32}の同定を用いると、六項側のminiversal hullの非被約性も従う。
\end{proof}

\begin{theorem}[被約局所枝と線・二次曲線成分]
\label{thm:triangle-branches-v32}
三角形点における $\mathscr P$ の被約完備局所スキームには、
ちょうど五つの既約枝があり、いずれも次元十九である。
十六個の滑らかな変数を除くと、それらの素イデアルは
\begin{align*}
 \mathfrak p_{\rm cub}&=(\delta_1,\delta_2,\rho_1,\rho_2),\\
 \mathfrak p_{\rm log}&=(a_1,a_2,a_3,
                  \delta_1\rho_2-\delta_2\rho_1),\\
 \mathfrak p_i&=(a_j,a_k,\delta_i,\rho_i)
       \qquad(\{i,j,k\}=\{1,2,3\})
\end{align*}
である。
最初の枝は滑らかな三次曲線成分、二番目は一般対数型成分の局所枝である。
残る三枝は同一の大域的既約成分の局所枝であり、
その一般点の三次曲線は直線と滑らかな二次曲線の和である。
この大域的成分は次元十九であり、前二者と異なる。
\end{theorem}

\begin{proof}
表示した五つのイデアルはすべて素で、定義イデアルを含む。
対数型の素性は二行二列の行列式錐の既約性から従う。
各枝の次元は三であり、互いに包含関係を持たない。

定義イデアルを含む任意の素イデアル $\mathfrak q$ を取る。
二つ以上の $a_i$ が $\mathfrak q$ に属さなければ、対応する
$\delta_i,\rho_i$ は $\mathfrak q$ に属し、和が零であることから残りも属する。
従って $\mathfrak p_{\rm cub}\subset\mathfrak q$ である。
ちょうど一つの $a_i$ が属さなければ、
$a_j,a_k,\delta_i,\rho_i\in\mathfrak q$ となり、
$\mathfrak p_i\subset\mathfrak q$ である。
全ての $a_i$ が属すれば $\mathfrak p_{\rm log}\subset\mathfrak q$ である。
これで五枝の完全性が示される。

$\mathfrak p_i$ の一般点では $a_i\ne0$、$\delta_j=-\delta_k\ne0$ であり、
例えば $i=1$ なら
\[
 c=x(ax^2+3yz),\qquad D=d(y\partial_y-z\partial_z),\quad a d\ne0
\]
という型になる。二次因子は滑らかで、直線とは二点で横断的に交わる。
その点では $a_j,a_k$ は線形に零へ消去され、残る三変数は自由であるから、
この枝はスキームとして一般点で滑らかである。
大域的には、正則半単純で固有値が $0,d,-d$ ($d\ne0$) である線形場の族が七次元、
その三次不変式の空間が
$\langle x^3,xyz\rangle$ の二次元であり、基底は既約な九次元族をなす。
ここでの一般混合核は十次元なので、その上の解の閉包は十九次元の既約成分となる。
三つの局所枝は座標の巡回置換で交換される。
線形座標変換の群は連結なので、大域的既約成分の有限集合への作用は自明であり、
これらは同じ大域的成分に属する。

滑らかな三次曲線成分は $D=0$ の中に含まれるため、この成分とは異なる。
一般対数型成分では $c$ は三つの線形因子の積であり、その性質は閉条件である。
上の一般の直線・二次曲線型はその閉集合に入らないので、対数型とも異なる。
\end{proof}

\begin{corollary}[被約局所枝の交叉]
\label{cor:triangle-intersections-v32}
上の五枝を $C_{\rm cub},C_{\rm log},C_i$ と書く。
相異なる二枝のスキーム論的交叉はすべて滑らかで、その次元は
\[
 \begin{array}{c|c}
 \text{交叉}&\text{次元}\\\hline
 C_{\rm cub}\cap C_{\rm log}&16\\
 C_{\rm cub}\cap C_i&17\\
 C_{\rm log}\cap C_i&18\\
 C_i\cap C_j\quad(i\ne j)&16
 \end{array}
\]
である。対数型枝 $C_{\rm log}$ の特異部分は
$C_{\rm cub}\cap C_{\rm log}$ に一致する。
\end{corollary}

\begin{proof}
表示した素イデアルの和を取ればよい。
十六変数を除くと、最初と最後の交叉は原点、二番目は $a_i$ の軸、
三番目は $\delta_i=\rho_i=0$ の二次元平面である。
対数型の三次元行列式錐は原点のみで特異である。
\end{proof}

\begin{remark}[二つの残余方向の意味]
三角形を固定し、$a_i=\delta_i=0$ とすると、$\rho_1+\rho_2+\rho_3=0$ の
二次元平面全体が残る。この平面は対数型枝の $\delta=0$ 部分であり、
三角形だけに支えられた新しい既約成分ではない。
$\delta$ と $\rho$ の行列式が零という条件が、
一般の対数型からこの境界への接続を記述する。
一方、$a_i^2\rho_i=0$ の二乗は被約化すると失われる。
従って既約成分の分類だけでは、補完空間の無限小構造を捉えきれない。
\end{remark}

\begin{remark}[局所成分数と大域成分数]
上の五つは一つの三角形点での完備局所枝の個数であり、
大域的な五つの成分を意味しない。
特に直線・二次曲線型成分は、この点で三つの枝を持つ。
本節は、この近傍に現れない他の大域成分を排除するものでもない。
\end{remark}

\section{$c=0$ の層の完全な被約分類}
\label{sec:c-zero-v33}

この節では重みを $(1,1,1,3)$ に固定する。Poisson スキームを
$\mathcal P$ とし、$c=0$ で切り出した閉部分スキームを
$\mathcal Z$ と書く。以下は $\mathcal Z_{\mathrm{red}}$ の分類であり、
$\mathcal Z$ の冪零構造を除去してよいという主張ではない。

\begin{theorem}[零三次形式の層]
\label{thm:c-zero-v33}
$\mathcal Z_{\mathrm{red}}$ は、次元 $35,20,18$ のちょうど三つの
既約成分を持つ。これらはそれぞれ、次の三族の閉包である。
\begin{enumerate}
\item $D=0$、$[F]\in S_4^{\oplus3}/xS_3$ は任意。
\item $D\ne0$ の最小多項式の次数が $2$ で、$[F]$ は対応する
十五次元の混合核を動く。
\item $D$ の最小多項式の次数が $3$ で、$[F]=[gD]$、$g\in S_3$。
\end{enumerate}
これらの族は、全 Poisson スキームにおいてそれぞれ既知の成分
\[
 \operatorname{PB}(\mathbb P^2),\qquad
 L(1,1,4),\qquad L(1,1,1,3)
\]
に含まれる。従って $c=0$ に台を持つ未知の既約成分は存在しない。
\end{theorem}

\begin{proof}
$c=0$ では残る方程式は
\begin{equation}
 (x\times D)\cdot F=0.
 \label{eq:c-zero-det-v33}
\end{equation}
まず $D$ の最小多項式が三次とする。このとき各固有空間は一次元であり、
射影空間上の $x\times D=0$ の零点は有限個である。従って三つの二次式
$x\times D$ が生成するイデアルの高さは二である。
行列 $[\,x\mid D\,]$ の最大小行列式は、符号を除きこの三式である。
商体上で $x,D$ は一次独立なので、式 \eqref{eq:c-zero-det-v33} の解は
$F=fx+gD$ と一意に書ける。高さ一の素イデアルで局所化すれば、
上の最大小行列式の少なくとも一つが可逆であるため、二行を使って解くことで
$f,g$ はその局所環に属する。多項式環は正規環であり、高さ一の局所環の共通部分に
等しいので、$f,g$ は多項式である。斉次性から次数はいずれも三となる。
これは本例の Hilbert--Burch 完全性を直接示す議論である。従って全解は一意に
\[
 F=fx+gD,\qquad f,g\in S_3
\]
と書ける。Euler 商では $[F]=[gD]$ となり、核の次元は十である。
ここで $x,D$ は商体上で一次独立なので、$g\mapsto[gD]$ は単射である。

相異なる固有値を持つ $D=\operatorname{diag}(\lambda_1,\lambda_2,\lambda_3)$
の場合、$u=w-g$ と置くと、対応する水平一形式は
\[
 \alpha=(w-g)(x\times D)\cdot dx
 =xyzu\left((\lambda_3-\lambda_2)\frac{dx}{x}
 +(\lambda_1-\lambda_3)\frac{dy}{y}
 +(\lambda_2-\lambda_1)\frac{dz}{z}\right).
\]
これは $u$ に沿う留数を零とした $L(1,1,1,3)$ の一形式である。
三次の最小多項式を持つ $D$ の集合は $\mathfrak{sl}_3$ の既約な開集合であり、
その上の全解は上記の十次元ベクトル束である。
相異なる固有値の集合は稠密なので、その全体が同じ成分の閉包に含まれる。
この族の次元は $8+10=18$ である。

次に $D\ne0$ の最小多項式の次数が二の場合、Jordan 標準形により、
座標変更と非零の定数倍を除いて次の二型のみである。
\[
 D=\lambda(x\partial_x+y\partial_y-2z\partial_z)\quad(\lambda\ne0),
 \qquad D=x\partial_z.
\]
いずれも式 \eqref{eq:c-zero-det-v33} は $xF_2=yF_1$ と同値である。
従って $F_1=xh,F_2=yh$ と書け、Euler 商で一意に
\[
 F=(0,0,f),\qquad f\in S_4
\]
とできる。この核は十五次元である。
二型に対応する一形式はそれぞれ
\[
 (3\lambda zw+f)(x\,dy-y\,dx),\qquad
 (-xw+f)(x\,dy-y\,dx)
\]
である。いずれも $Q(x\,dy-y\,dx)$ という形を持ち、
因子 $x,y,Q$、重み $(1,1,4)$、留数 $(-1,1,0)$ の対数型に属する。
因子が特殊化している場合も $L(1,1,4)$ の閉包への包含が従う。

最小多項式の次数が二以下の行列の閉集合は既約で次元五である。
この集合は $I,D,D^2$ の一次従属という閉条件で定まる。
実際、階数一以下の行列 $N$ に対して
\[
 D=N-\frac{\operatorname{tr}N}{3}I
\]
とすると、$\operatorname{tr}N\ne0$ では二重固有値を持つ半単純型、
$\operatorname{tr}N=0$ では階数一の冪零型となり、上の非零二型を全て覆う。
階数一以下の行列の錐は既約で次元五であり、この写像は単射である。
単射性は、相異なる二つの階数一行列の差が非零スカラー行列にはなれないことから従う。
零行列を除いた基底上では混合核の次元は一定の十五であり、
対応する族はベクトル束、従って既約で次元 $5+15=20$ である。

最後に $D=0$ のファイバーは全 $35$ 次元の線形空間であり、
既知の $\operatorname{PB}(\mathbb P^2)$ 成分そのものである。
三つの局所閉族は $\mathcal Z$ の全ての閉点を尽くす。
十八次元族の一般点では $D$ は最小多項式三次なので、二十次元族の閉包に含まれない。
二十次元族は次元から十八次元族の閉包に含まれない。
また非零の $D$ を持つこれらの族は $D=0$ の三十五次元空間には含まれず、
逆向きの包含も次元から不可能である。
従って三つの閉包が、ちょうど $\mathcal Z_{\mathrm{red}}$ の既約成分である。

全 Poisson スキームの既約成分が $c=0$ に含まれていれば、
その既約性から上記三つの閉包のいずれかに含まれる。
いずれも既知の全スキームの既約成分に含まれるため、
成分の極大性から未知の成分は生じない。
\end{proof}

\begin{remark}[一般重みでの零三次形式の層]
同じ議論は重み $(1,1,1,m)$、$m\ge3$ にも適用できる。
この場合の混合項は $F\in S_{m+1}^{\oplus3}/xS_m$ であり、
$c=0$ の方程式は非零定数倍を除き依然として $\det(x,D,F)=0$ である。
最小多項式三次の型では $[F]=[gD]$、$g\in S_m$、
非零の二次型では $[F]=[(0,0,f)]$、$f\in S_{m+1}$ となる。
従って零三次形式の層の被約化は、ちょうど三つの既約成分を持ち、
それらの次元はそれぞれ
\[
 m^2+6m+8,\qquad 5+\binom{m+3}{2},\qquad
 8+\binom{m+2}{2}
\]
である。証明は上の次数四・三を次数 $m+1$・$m$ に置き換えるだけである。

この重みでの水平形式は、$c=0$ では
$\alpha=\tfrac m3w(x\times D)\cdot dx+(F\times x)\cdot dx$ となる。
従って後二族はそれぞれ対数型の族
$L(1,1,m+1)$ と $L(1,1,1,m)$ の閉包に含まれる。
ここでは、一般の $m$ においてこれらの対数型の閉包が全 Poisson スキームの
既約成分であることは主張しない。
定理~\ref{thm:c-zero-v33}の「$c=0$ に台を持つ未知の全スキームの成分はない」
という結論は、包含先の成分性が確立している $m=3$ に限って用いる。
\end{remark}

この定理は、$c\ne0$ の非被約三次形式や特殊 Jordan 型に関する残る分類を
自動的に解決するものではない。また $\mathcal Z_{\mathrm{red}}$ の
次元 $20,18$ の成分は、全 Poisson スキームの既約成分ではない。

\section{明示的なAS正則4-Calabi--Yau量子化}
\subsection{半古典規約と発散零代表}
通常の半古典括弧 $\{a,b\}=([a,b]/\hbar)\bmod\hbar$ を用いる。
重み付きEuler場 $E=E_x+3w\partial_w$ と標準体積形式の発散 $\delta$ に対し、
重み零多重ベクトルでは $\delta(E\wedge U)+E\wedge\delta U=6U$ である。
従ってEuler同値類の発散零代表は $\widehat P=P-E\wedge\delta P/6$ で一意に定まり、
本稿の座標では
\[
 \widehat P=\Pi+\frac16E_x\wedge D+
 \left(F-\frac{\operatorname{div}F}{6}x-\frac w2D\right)\wedge\partial_w.
\]
発散零部分はSchouten括弧で閉じるため、射影的なPoisson条件はこの代表の通常のPoisson条件と同値である。
通常の射影空間の対応は\cite[\S2.2]{Pym}、重み付き分解は\cite{TWW}に位置付けられる。
以下はこの一般の持ち上げを、明示的な反復Ore拡大で量子化する構成である。


\label{subsec:resonance-ore-cy-v32}

この23次元成分の一般点は、具体的な反復Ore拡大によって量子化できる。
構成そのものは二変数四次式が任意である場合にも成立する。

\begin{proposition}[二変数四次式のOre族]
\label{prop:binary-quartic-ore-v32}
$a,b\in\mathbb C^\times$ と $f\in\mathbb C[y,z]_4$ を取る。
$\deg x=\deg y=\deg z=1$、$\deg w=3$ とした代数
\[
 S(a,b,f)=\frac{\mathbb C\langle y,z,x,w\rangle}{
 \left(\begin{array}{c}
 yz-zy,\ xy-ayx,\ xz-bzx,\\
 wy-a^{-1}yw,\ wz-b^{-1}zw,\\
 wx-abxw+f(y,z)
 \end{array}\right)}
\]
は順序単項式 $y^iz^jx^kw^\ell$ を基底に持つ。
この代数はnoetherian整域、global dimension 4 のAS正則代数であり、
通常の非次数付きの意味で4-Calabi--Yauである。
Hilbert級数は
\[
 H_S(u)=\frac1{(1-u)^3(1-u^3)}
\]
である。特にGK次元は4である。
\end{proposition}

\begin{proof}
$R_0=\mathbb C[y,z]$、
$\sigma(y)=ay$、$\sigma(z)=bz$ とし、$A=R_0[x;\sigma]$ と置く。
$A$ の自己同型 $\tau$ と $\tau$-導分 $\partial$ を
\[
 \tau|_{R_0}=\sigma^{-1},\quad \tau(x)=abx,
 \qquad \partial|_{R_0}=0,\quad \partial(x)=-f
\]
で定める。$\tau$ は $xr=\sigma(r)x$ を保ち、
\[
 \partial(xr-\sigma(r)x)=-fr+rf=0
\]
なので $\partial$ も定義される。
従って $S=A[w;\tau,\partial]$ であり、PBW基底とHilbert級数は直ちに従う。
自己同型による反復Ore拡大はnoetherian性と整域性を保つ。

Liu--Wang--WuのOre拡大定理\cite[Theorem 3.3]{LWW} により、
$S$ はtwisted Calabi--Yau次元4である。
三変数量子多項式環 $A$ のNakayama自己同型は
\[
 \mu_A(y)=a^{-1}y,\qquad \mu_A(z)=b^{-1}z,
 \qquad \mu_A(x)=abx
\]
である。同じ定理から、$S$ のNakayama自己同型は
\[
 \mu_S|_A=\tau^{-1}\mu_A=\operatorname{id}_A,
 \qquad \mu_S(w)=u w+b_0
\]
を満たす。ただし $u\in A^\times=\mathbb C^\times$、$b_0\in A$ である。

補助的な整数次数
\[
 \deg' x=1,\qquad\deg' w=-1,\qquad\deg' y=\deg' z=0
\]
に関する $\mathbb G_m$ 作用を考える。
正の重み次数に関して $S$ は連結次数付き整域なので、単元はスカラーのみである。
従ってNakayama自己同型は内的自己同型の不定性を持たず一意であり、
この $\mathbb G_m$ 作用と可換する。
$b_0$ は補助次数 $-1$ を持つ必要があるが、$A$ の補助次数は非負である。
よって $b_0=0$ である。
$f\neq0$ のとき、最後の関係に $\mu_S$ を適用すると
\[
 (1-u)f=0
\]
を得るので $u=1$。$f=0$ のときは四変数量子多項式環であり、
交換係数の各行の積が1であることからNakayama自己同型は恒等である。
従ってすべての $f$ に対し $\mu_S=\operatorname{id}$ となる。
\cite[Lemma 1.3]{LWW}によって連結次数付きのtwisted Calabi--Yau代数はAS正則であり、
この場合はHilbert級数から有限GK次元も満たす。
\end{proof}

\begin{corollary}[共鳴成分の半古典極限]
\label{cor:resonance-quantization-v32}
$3\alpha+t=0$ とし、
\[
 \theta_i=\frac{3\alpha_i+t}{6},\qquad
 (\alpha_1,\alpha_2,\alpha_3)=(\alpha,\beta,\gamma)
\]
と置く。形式パラメータ $\hbar$ に対し
\[
 a=e^{\hbar\theta_3},\qquad b=e^{-\hbar\theta_2},\qquad
 S_\hbar=S(a,b,\hbar f)
\]
と取れば、$S_\hbar$ は共鳴標準形
$P=\Pi+f(y,z)\partial_x\wedge\partial_w$ の
unimodular liftを半古典極限に持つ、重み付き次数を保つ平坦形式量子化である。
\end{corollary}

\begin{proof}
壁の条件は $\theta_1=0$ を意味し、
\[
 \frac{\delta_2}{2}=-\theta_3,\qquad
 \frac{\delta_3}{2}=\theta_2,\qquad
 \frac{\delta_1}{2}=\theta_3-\theta_2
\]
である。従って $ab=e^{\hbar\delta_1/2}$ となる。
関係式の一次交換子は
\begin{align*}
 &\{y,z\}=0,\quad\{z,x\}=\theta_2zx,
       \quad\{x,y\}=\theta_3xy,\\
 &\{x,w\}=f-\frac{\delta_1}{2}xw,
 \quad\{y,w\}=-\frac{\delta_2}{2}yw,
 \quad\{z,w\}=-\frac{\delta_3}{2}zw.
\end{align*}
これは明示的なunimodular lift
\[
 \widehat P=\Pi+\frac16E_x\wedge D+
       \left(f\partial_x-\frac w2D\right)\wedge\partial_w
\]
と一致する。ここでは $\operatorname{div}(f\partial_x)=0$ を使った。
全次数の平坦性はOre基底による。
\end{proof}


\begin{corollary}[標準六項列とVK部分圏の実現]
\label{cor:ore-six-objects-v32}
命題~\ref{prop:binary-quartic-ore-v32} の代数 $S=S(a,b,f)$ のGorenstein parameterは六であり、
\[
 (\mathcal O_S,\mathcal O_S(1),\ldots,\mathcal O_S(5))
\]
は $D^b(\operatorname{qgr}S)$ のfull strong exceptional collectionである。
従って
\[
 B(S)=\End\bigl(\mathcal O_S\oplus\mathcal O_S(1)
              \oplus\mathcal O_S(3)\oplus\mathcal O_S(4)\bigr)
       \simeq S_{(0,1,3,4)}
\]
に対し、$\Perf(B(S))$ は $D^b(\operatorname{qgr}S)$ の許容部分圏である。
残る二つの例外対象をmutationすれば、これを一成分とする半直交分解を得る。
\end{corollary}

\begin{proof}
$S$ は連結次数付きnoetherian AS正則代数で、global dimensionは四である。
Gorenstein parameterを $\ell$ と書く。
自明加群の有限次数付き自由分解とそのAS--Gorenstein双対を比較すると、
Hilbert級数の相反関係
\[
 H_S(t^{-1})=(-1)^4t^\ell H_S(t)
\]
を得る。命題のHilbert級数では左辺は $t^6H_S(t)$ なので $\ell=6$ である。
Orlovの定理\cite[Corollary 2.7]{Orlov2009}を適用すれば、表示した六項列が
full strongであり、その自己準同型代数が次数切り出しであることが従う。
部分例外列が生成する部分圏は許容であり、strong性からその圏は $\Perf(B(S))$ に同値である。
\end{proof}

ここでの $\operatorname{qgr}S$ は有限生成次数付き右加群を有限次元加群で割った圏である。
$a=b=1,f=0$ では重み付き射影\emph{スタック}の連接層の圏を回復する。
したがって本系は、この明示族に関して、有限六項代数の形式補完に加えて
非可換重み付き射影空間の実際の標準例外列による実現を与えている。

上のAS正則性・Calabi--Yau性はこの明示Ore族に対する結論である。
一般の形式量子化が同じ性質を持つ、あるいはこの構成が全Poisson成分を量子化の同型まで
分類するという主張ではない。

\section{証明の範囲と再現}

\subsection{確定した一般現象と残る分類}

全整数 $m\ge3$ で、引戻し成分の一般点のArtin環
$K[a,b]/(a^{m+2},ab,b^2)$ と交叉環を決定した。
接線系列では、$r=m-2j>0$ の全範囲で既約成分性と
全形式局所環 $(\kappa u,\gamma^{j+1}u)$ を証明した。
一般点の長さ $j+1$ は成分に固有であり、固定した重みでも互いに異なる長さが現れる。
$m=1,2$ には二重ベクトル場の追加係数があるので、同じ成分定理の対象には含めない。
接線成分定理の $r>0$ と、量子化・中心の定理の $r\ge0$ も区別する。

全重みの非線形アフィン族について、Nakayama自己同型、通常のCY条件、
次数付き自己同型群、多項式中心、次数零の関数斜体を決定した。
多項式中心がDu Val曲面であるという結論は、その曲面の座標環を中心として含むという意味であり、
曲面自身に非可換な積を入れたという意味ではない。
非零の量子化係数 $s$ に対するこれらの代数は非PIであり、中心上有限な代数として扱わない。
同じ重みの異なる共鳴の射影圏が同値かどうかは未決である。

重み三では少なくとも十成分を認定したが、全分類と最大次元は未決である。
$c=0$ の被約部分の三成分とその大域所属、非零の接線型線形場に対する
四共鳴の形式近傍は決定した。
残る層には非被約三次形式と、その他の特異三次形式の特殊係数がある。
混合核の増大だけから新成分を数えることはしない。

三角形の三つの直線・二次曲線枝は、同じ $L(1,2,3)$ 成分に属する。
共点三直線の交叉定理は、明記した非空開集合に対するものであり、
全境界点へ同じ局所式を延長する主張ではない。
全ての次元は自己同型群とスカラー倍で商する前の値である。

\subsection{構造的証明と有限証明書}

本版の統合検証器は \nolinkurl{code/verify_v34.py} である。
全重みの共鳴に関する記号恒等式、重み三の全Schur式、
Nakayamaの規約を確認する三次syzygy、不変環の関係式・等変変換と中心式を検証する。
通常実行では入力や証明書を再生成せず、全JSON出力を保存済み証明書と比較する。
コード・入力・証明書の実行前後のhashも比較する。
新しい定理の独立監査を保存し、過去の原稿・計算入力・検証記録を保持する。

全重みの結論は有限個の検算からの外挿ではない。
混合核の二重対角表示、横断二次曲線の次数降下、余核指標、
特異点評価が、任意の $m,j$ でSchurイデアルを決める。
不変環の生成性は局所化と $x$ 飽和で、中心の生成性は全ての零重量単項式の分類で証明する。
AS正則性、CY性、例外列の存在は、本文の構造的定理と適用条件による。
補助的な有限次数の核次元検査は、これらの生成性の証明を置き換えない。

miniversal基底の同定は形式座標と量子化の選択を伴う。
固定した非零量子代数の全変形空間、全Poisson成分のAS正則量子化、
元の四対象cornerの全変形についての量子座標環の実現可能性は、本稿では決定していない。
$m>3$ の四対象cornerが無障害であることや、外側との $HH^2$ 制限が同型であることも未証明である。
重み三の標準補完に関する\cite{VKAmbient32}の主張を、
別の例外列や一般重みに拡張しない。

\begin{thebibliography}{99}
\interlinepenalty=10000
\bibitem{DoninMakarLimanov1998}J.~Donin and L.~Makar-Limanov,
\emph{Quantization of quadratic Poisson brackets on a polynomial algebra of three variables},
Journal of Pure and Applied Algebra \textbf{129} (1998), 247--261.
\url{https://arxiv.org/abs/q-alg/9504002}.

\bibitem{LiuXu1992}Z.-J.~Liu and P.~Xu,
\emph{On quadratic Poisson structures},
Letters in Mathematical Physics \textbf{26} (1992), 33--42.
\url{https://doi.org/10.1007/BF00420516}.

\bibitem{Pym}B.~Pym,
\emph{Quantum deformations of projective three-space},
Advances in Mathematics \textbf{281} (2015), 1216--1241.
\url{https://arxiv.org/abs/1403.6444}.

\bibitem{TWW}X.~Tang, X.~Wang and J.~J.~Zhang,
\emph{Twists of graded Poisson algebras and related properties},
Journal of Geometry and Physics \textbf{207} (2025), 105344.
\url{https://arxiv.org/abs/2206.05639}.

\bibitem{Molinuevo2016}A.~Molinuevo, \emph{Unfoldings and deformations of rational and logarithmic foliations},
Annales de l'Institut Fourier \textbf{66} (2016), 1583--1613.
\url{https://doi.org/10.5802/aif.3044}.

\bibitem{PereiraClosedForms}J.~V.~Pereira, \emph{Closed meromorphic 1-forms},
arXiv:2206.09745v3 (2023), \S5, Corollary~5.6.
\url{https://arxiv.org/abs/2206.09745}.

\bibitem{GAMQV2026}J.~Gargiulo Acea, A.~Molinuevo, F.~Quallbrunn and S.~Velazquez,
\emph{Stability of pullbacks of foliations on weighted projective spaces},
Journal of Pure and Applied Algebra \textbf{230} (2026), 108248.
\url{https://doi.org/10.1016/j.jpaa.2026.108248}.
\url{https://arxiv.org/abs/2212.12974}.

\bibitem{LWW}L.-Y.~Liu, S.-Q.~Wang and Q.-S.~Wu,
\emph{Twisted Calabi--Yau property of Ore extensions},
Journal of Noncommutative Geometry \textbf{8} (2014), 587--609.
\url{https://doi.org/10.4171/JNCG/165}.
\url{https://arxiv.org/abs/1205.0893}.

\bibitem{Orlov2009}D.~Orlov, \emph{Derived categories of coherent sheaves and triangulated categories of singularities},
in \emph{Algebra, Arithmetic, and Geometry: In Honor of Yu. I. Manin}, Vol.~II,
Progress in Mathematics \textbf{270}, Birkh\"auser, 2009, 503--531.
\url{https://doi.org/10.1007/978-0-8176-4747-6_16}.
\url{https://arxiv.org/abs/math/0503632}.

\bibitem{VKAmbient32}本研究、公表原稿VIII、\emph{半直交成分の変形と重み付き射影空間への補完}, v32.
\nolinkurl{vk_ambient_completions_ja_v32.pdf}.

\bibitem{Kontsevich}M.~Kontsevich,
\emph{Deformation quantization of Poisson manifolds},
Letters in Mathematical Physics \textbf{66} (2003), 157--216.
\url{https://www.ihes.fr/~maxim/TEXTS/DefQuant_final.pdf}.

\bibitem{FPSS}L.~Fu, M.~Porta, S.~Scherotzke and N.~Sibilla,
\emph{Hochschild--Kostant--Rosenberg isomorphism for derived Deligne--Mumford stacks},
arXiv:2509.00501, author version consulted 29 October 2026, Corollary~4.16.
\url{https://irma.math.unistra.fr/~lfu/articles/HKR_for_derived_DM_stacks.pdf}.

\bibitem{LiuMa2020}L.~Liu and W.~Ma,
\emph{Nakayama automorphisms of Ore extensions over polynomial algebras},
Glasgow Mathematical Journal \textbf{62} (2020), 518--530.
\url{https://doi.org/10.1017/S0017089519000259}.
\url{https://arxiv.org/abs/1808.03589}.

\bibitem{BKM2026}V.~Benedetti, C.~Kuster and A.~Muniz,
\emph{Minimal-degree foliations on cominuscule Grassmannians},
Bulletin of the Brazilian Mathematical Society, New Series \textbf{57} (2026), article~47.
\url{https://doi.org/10.1007/s00574-026-00533-3}.

\bibitem{Kaplansky1972}I.~Kaplansky, \emph{Fields and Rings}, second edition,
Chicago Lectures in Mathematics, University of Chicago Press, 1972,
Part~II, Theorem~50, p.~157.
\url{https://webhomes.maths.ed.ac.uk/~v1ranick/papers/kapfield.pdf}.

\bibitem{LecoutreSierra2019}
C.~Lecoutre and S.~J.~Sierra,
\emph{A new family of Poisson algebras and their deformations},
Nagoya Math. J. \textbf{233} (2019), 32--86,
\href{https://doi.org/10.1017/nmj.2017.29}{doi:10.1017/nmj.2017.29};
\href{https://arxiv.org/abs/1611.04306}{arXiv:1611.04306}.

\bibitem{Zhang1996}
J.~J.~Zhang,
\emph{Twisted graded algebras and equivalences of graded categories},
Proc. London Math. Soc. (3) \textbf{72} (1996), 281--311,
\href{https://doi.org/10.1112/plms/s3-72.2.281}{doi:10.1112/plms/s3-72.2.281}.

\bibitem{Maubach2000}
S.~Maubach, \emph{Triangular monomial derivations on $k[X_1,X_2,X_3,X_4]$
have kernel generated by at most four elements},
J. Pure Appl. Algebra \textbf{153} (2000), no.~2, 165--170.
\url{https://doi.org/10.1016/S0022-4049(99)00084-5}.
\bibitem{ReidDuVal}
M.~Reid, \emph{The Du Val singularities $A_n,D_n,E_6,E_7,E_8$},
author's notes, Table 1.
\url{https://mreid.warwick.ac.uk/surf/more/DuVal.pdf}.

\end{thebibliography}

\end{document}
